% =========================================================================
% Apartado metodológico --- sistema de recuperación multisalto sobre grafo de
% aserciones reificadas con memoria canónica de entidades, y comparación
% controlada con HippoRAG 2 sobre 2WikiMultihopQA.
%
% Compilar:  pdflatex metodologia_2wiki.tex  (dos pasadas)
% Requiere:  babel (spanish), amsmath, booktabs, tikz
%
% El nombre del sistema es provisional: cambiar aquí una sola vez.
% =========================================================================
\documentclass[11pt,a4paper]{article}
\usepackage[T1]{fontenc}
\usepackage[utf8]{inputenc}
\usepackage[spanish,es-noshorthands,es-tabla]{babel}
\usepackage{lmodern}
\usepackage[margin=2.4cm]{geometry}
\usepackage{amsmath,amssymb,mathtools}
\usepackage{booktabs}
\usepackage{array}
\usepackage{tikz}
\usetikzlibrary{arrows.meta,positioning,fit,backgrounds,calc,shapes.geometric}
\usepackage{caption}
\usepackage{enumitem}
\usepackage[hidelinks]{hyperref}

\newcommand{\sys}{\textsc{CatRAG-Reif}}          % nombre provisional del sistema
\newcommand{\hippo}{HippoRAG~2}
\newcommand{\op}[1]{\operatorname{#1}}
\newcommand{\R}{\mathbb{R}}
\newcommand{\norm}[1]{\lVert #1 \rVert}
\DeclareMathOperator*{\argmax}{arg\,max}

\setlength{\parskip}{2pt}

\title{Recuperación multisalto sobre un grafo de aserciones reificadas\\
con memoria canónica de entidades:\\
método, caso de uso clínico (GeSIDA) y evaluación comparada}
\author{}
\date{}

\begin{document}
\maketitle
\vspace{-3em}

% =========================================================================
\section{Introducción y estado del arte}
\label{sec:intro}

\subsection{Contexto y motivación}

La generación aumentada por recuperación (RAG) responde preguntas
apoyándose en pasajes recuperados de un corpus. En las preguntas
\emph{multisalto} ---aquellas cuya respuesta exige componer hechos de dos o
más documentos--- la recuperación por semejanza vectorial fracasa de un modo
característico: encuentra el documento nombrado en la pregunta, pero no el
documento \emph{puente} (la biografía del director cuando la pregunta habla
de una película), porque este comparte poco vocabulario y poca semántica
superficial con la consulta. El fallo resultante es sutil: las métricas de
exhaustividad parcial parecen altas, pero la \emph{cadena de evidencia}
llega incompleta al modelo lector, que entonces alucina o acierta sin
respaldo. Por ello, este trabajo adopta como métrica principal la
recuperación de la cadena completa
(\emph{full chain}, sección~\ref{sec:protocolo}), en línea con la
literatura reciente~\cite{catrag}.

\subsection{Estado del arte}

\paragraph{Recuperación densa.}
Los codificadores densos ---desde Contriever~\cite{contriever} y
GTR~\cite{gtr} hasta modelos de 7.000 millones de parámetros como
NV-Embed-v2~\cite{nvembed}--- dominan la recuperación de un salto, pero se
estancan en la composición: en 2WikiMultihopQA, el mejor codificador denso
del estudio de referencia alcanza $76{,}5$ de exhaustividad a 5
pasajes~\cite{hipporag2} y los pequeños quedan en $57$--$71$
(tabla~\ref{tab:sota}). El vector de una pregunta no puede parecerse a un
documento cuyo vínculo con ella es un hecho que está en \emph{otro}
documento.

\paragraph{RAG estructurado.}
Una segunda familia impone estructura sobre el corpus.
RAPTOR~\cite{raptor} construye un árbol jerárquico de resúmenes;
GraphRAG~\cite{graphrag} deriva comunidades de un grafo de entidades y las
resume para preguntas de síntesis global; LightRAG~\cite{lightrag} combina
un grafo plano con recuperación vectorial de doble nivel (entidades y
temas); HyperGraphRAG~\cite{hypergraphrag} generaliza la arista a
hiperarista $n$-aria, evidencia independiente de que el triplete binario es
una limitación reconocida. HopRAG~\cite{hoprag} conecta pasajes con aristas
lógicas (pseudo-consultas generadas por un modelo de lenguaje) y recorre el
grafo con un bucle recuperar--razonar--podar, a coste de múltiples llamadas
por consulta.

\paragraph{La línea HippoRAG y CatRAG.}
HippoRAG~\cite{hipporag} y HippoRAG~2~\cite{hipporag2}, inspirados en la
teoría de la indexación hipocampal de la memoria, construyen un grafo de
entidades y pasajes mediante extracción abierta y recuperan con PageRank
personalizado (PPR)~\cite{pprsurvey}: la pregunta se ancla en unos nodos
semilla y un paseo aleatorio con reinicio difunde relevancia por el grafo.
HippoRAG~2 añade la puntuación consulta--triplete, un filtro de
\emph{memoria de reconocimiento} y la fusión con recuperación densa, y
alcanza el estado del arte en 2Wiki ($90{,}4$ de exhaustividad a 5 con
NV-Embed-v2). PropRAG~\cite{proprag} sustituye los tripletes por
proposiciones y busca caminos con búsqueda en haz. CatRAG~\cite{catrag}
diagnostica en esta línea la \guillemotleft{}falacia del grafo estático\guillemotright{}
---las probabilidades de transición se fijan en indexación--- y la ataca
con tres mecanismos de consulta sobre el grafo de HippoRAG~2, que dejan
la cadena completa en 2Wiki en $67{,}6\%$ (tabla~\ref{tab:sota}). Por
ser la genealogía directa de este trabajo, la sección~\ref{sec:palancas}
describe los tres sistemas con detalle conceptual y funcional, y muestra
sobre el corpus de evaluación por qué la mejora de CatRAG es real pero
acotada.

\paragraph{Grafos como memoria de agentes.}
En paralelo, los sistemas de memoria para agentes conversacionales
---Zep/Graphiti~\cite{zep} con su grafo temporal de entidades y episodios,
MAGMA~\cite{magma} con memorias multigrafo, y la taxonomía
de~\cite{agentmemory}--- convergen en la misma pieza que aquí resulta
central: un registro canónico de entidades mantenido en el tiempo, sobre el
que se acumulan hechos con procedencia. Nuestra \emph{memoria canónica de
entidades} (sección~\ref{sec:diccionario}) es esa pieza, construida de una
vez sobre un corpus cerrado.

\subsection{De HippoRAG a CatRAG: las dos palancas del paseo con reinicio}
\label{sec:palancas}

La línea que va de HippoRAG a CatRAG es la genealogía directa de este
trabajo, y conviene describirla con más detalle del habitual en un estado
del arte, porque las aportaciones de la sección~\ref{sec:metodos} se
entienden como respuesta a lo que esa evolución deja abierto. El hilo
conductor es una distinción sencilla que los tres sistemas comparten y
que cada uno explota de forma distinta.

\paragraph{Las dos palancas.}
Un paseo aleatorio con reinicio personalizado sobre un grafo calcula el
punto fijo
\begin{equation}
\mathbf{p} \;=\; \lambda\,P^{\!\top}\mathbf{p} \;+\; (1-\lambda)\,\mathbf{v},
\qquad\text{es decir}\qquad
\mathbf{p}^{*} \;=\; (1-\lambda)\sum_{\ell\ge 0}\lambda^{\ell}\,(P^{\!\top})^{\ell}\,\mathbf{v},
\label{eq:neumann}
\end{equation}
donde $\mathbf{v}$ es el \emph{vector de reinicio} ---dónde nace la masa y
adónde vuelve en cada reinicio, con probabilidad $1-\lambda$ por paso--- y
$P$ la \emph{matriz de transición} ---adónde va la masa cuando sale de un
nodo, $P(u\to w)=w_{uw}/\sum_{w'}w_{uw'}$---. La serie de la derecha hace
visible el papel de cada objeto: $\mathbf{v}$ es la \emph{fuente} de toda
la masa, y $P$ el \emph{coeficiente} que la propaga a distancia $\ell$ con
peso $\lambda^{\ell}$. Todo mecanismo de consulta de un recuperador basado
en PPR toca una de las dos palancas ---la fuente o el coeficiente--- y la
diferencia no es cosmética: dar masa de reinicio a un nodo lo alimenta en
cada reinicio, venga de donde venga el paseo y sin quitársela a nadie;
subir el peso de una arista solo surte efecto si el paseo llega a su
origen, y lo hace a costa de las aristas hermanas, porque la fila se
renormaliza. Con esa lente, los cuatro sistemas de la línea se resumen en
la tabla~\ref{tab:palancas}.

\paragraph{HippoRAG: la pregunta entra solo por la siembra.}
HippoRAG~\cite{hipporag} construye en indexación un grafo de frases
(las entidades extraídas por extracción abierta) con tres clases de arista
y peso fijo: aristas de relación entre dos frases, con peso proporcional
al número de tripletes que las unen;\footnote{Su código inserta cada
triplete en ambos sentidos sobre un grafo no dirigido, de modo que cada
triplete aporta peso $2$ entre sus dos frases; los valores de la tabla
siguiente son los efectivos.} aristas pasaje--frase, con peso $1$; y
aristas de sinonimia entre frases, con peso igual al coseno de sus
incrustaciones cuando supera $0{,}80$. En consulta aplica reconocimiento
de entidades a la pregunta, empareja cada entidad con el nodo-frase más
afín y pone en $\mathbf{v}$ exactamente esos nodos, ponderados por su
especificidad (inversa al número de pasajes en que aparecen): la siembra
es \emph{fuerte} y es la única vía por la que la pregunta entra en el
paseo; $P$ es la del corpus. El fallo característico es de anclaje: si
la mención de la pregunta no casa con ningún nodo, no hay semilla.

\paragraph{\hippo{}: hechos, filtro y pesos fijados en indexación.}
\hippo{}~\cite{hipporag2} conserva el grafo y sus pesos y cambia la
siembra. Incrusta cada triplete $(s,r,o)$ en un almacén paralelo; en
consulta puntúa la pregunta contra todos los tripletes, toma los más
afines y un modelo de lenguaje ---la \emph{memoria de reconocimiento}---
selecciona los útiles. Las entidades de los tripletes aprobados,
\emph{sujeto y objeto}, forman la siembra de frases, y todos los pasajes
reciben además su puntuación densa:
\begin{equation}
\mathbf{v}(\varepsilon) \;=\;
\frac{1}{|T_\varepsilon|}\sum_{t\in T_\varepsilon}\frac{s_q(t)}{|P_\varepsilon|},
\qquad
\mathbf{v}(c) \;=\; 0{,}05\cdot\widehat{\cos}(q,c)\quad\text{para todo pasaje } c,
\label{eq:siembra-hippo}
\end{equation}
donde $T_\varepsilon$ son los tripletes aprobados en que participa la
frase $\varepsilon$, $s_q(t)$ la afinidad consulta--triplete y
$|P_\varepsilon|$ el número de pasajes en que aparece la frase
(corrección anti-concentrador). Si el filtro no aprueba ningún triplete,
el sistema degrada a recuperación densa pura. Dos hechos de esta
descripción, verificados sobre su código, son los que importan aquí.
Primero, \hippo{} \emph{no} aplica reconocimiento de entidades a la
pregunta: el sujeto de la pregunta entra en $\mathbf{v}$ solo si el filtro
aprueba un triplete que lo contiene, y el \emph{segundo salto} ---la
entidad puente--- entra por el mismo camino, como objeto del triplete
aprobado. Segundo, el paseo se ejecuta con los pesos guardados en el
grafo: $P$ es idéntica para todas las preguntas, y la señal de la consulta
entra únicamente por $\mathbf{v}$.

\paragraph{La falacia del grafo estático, con los números del corpus.}
CatRAG~\cite{catrag} da nombre a la consecuencia: como las transiciones
se fijaron al indexar, las aristas genéricas y los nodos concentradores
desvían la masa antes de que alcance la evidencia del segundo salto. Para
hacerlo tangible tomamos filas reales de la matriz $P$ del grafo que el
código de \hippo{} construyó sobre nuestro corpus de evaluación
(sección~\ref{sec:comparacion}; $56\,127$ nodos y $305\,826$ aristas):
\begin{center}\small\setlength{\tabcolsep}{4pt}
\begin{tabular}{llrrl}
\toprule
nodo-frase $u$ & vecino $w$ & $w_{uw}$ & $P(u\to w)$ & observación\\
\midrule
god's gift to women & michael curtiz & $2$ & $0{,}105$ & un triplete\\
god's gift to women & musicals, other countries, frank fay, \dots & $2$ & $0{,}105$ & idéntica\\
god's gift to women & pasaje \emph{God's Gift to Women} & $1$ & $0{,}053$ & \\
aldri annet enn bråk & edith carlmar & $6$ & $0{,}316$ & tres tripletes\\
edith carlmar & bedre enn sitt rykte, fjols til fjells, \dots & $2$ & $0{,}037$ & sus otras películas\\
michael curtiz & (52 vecinos, $Z_u=103$) & & $\le 0{,}039$ & su biografía a $1/103$\\
\bottomrule
\end{tabular}
\end{center}
Estas cifras las decidió el extractor en indexación ---cuántos tripletes
produjo entre cada par de frases--- y valen para cualquier pregunta. Para
\guillemotleft{}?`qué director es mayor, el de \emph{God's Gift to Women} o el de
\emph{Aldri annet enn bråk}?\guillemotright{}, la masa que sale de la primera película
va a su director con la misma probabilidad que a \guillemotleft{}musicals\guillemotright{},
mientras que de la segunda película a su directora se cruza tres veces
más fácilmente; y la biografía de Michael Curtiz, un concentrador de 52
vecinos, recibe una centésima de lo que llega a su nodo. El resultado
real de \hippo{} en esa pregunta es el previsible: las dos películas y
Edith Carlmar en los puestos 1--3, \emph{otras} películas de Carlmar en
los puestos 4--5 y Michael Curtiz ausente del top-10 (la
sección~\ref{sec:fallan-otros} la analiza junto a la salida de \sys{}).

\paragraph{CatRAG: reponderar el coeficiente y añadir anclas débiles.}
CatRAG deja intactos el grafo, la siembra de la
ecuación~\eqref{eq:siembra-hippo} y el filtro de \hippo{}, y añade tres
mecanismos de consulta. Sea $\mathcal{S}_q$ el conjunto \emph{semilla}
---las $N_{\text{sem}}=5$ frases de mayor $\mathbf{v}$--- y
$T_{\text{sem}}$ los tripletes aprobados por el filtro.
\begin{enumerate}[leftmargin=1.6em, itemsep=2pt]
\item \textbf{Anclas simbólicas débiles} (palanca $\mathbf{v}$). Recupera
  el reconocimiento de entidades sobre la pregunta de HippoRAG, pero con
  masa pequeña: cada entidad reconocida $e$ recibe
  $\mathbf{v}(e) \mathrel{+}= \epsilon/|P_e|$ con $\epsilon=0{,}2$. Cubre
  el hueco de \hippo{}: la entidad nombrada en la pregunta ---o una
  palabra de tipo, \guillemotleft{}universidad\guillemotright{}--- siembra aunque el filtro haya
  descartado su triplete.
\item \textbf{Reponderación dinámica de aristas} (palanca $P$). Para cada
  semilla $u\in\mathcal{S}_q$ se toman sus $K=15$ aristas de relación
  salientes más afines a la pregunta (coseno entre la pregunta y la
  incrustación del triplete de la arista) y un modelo de lenguaje puntúa
  cada vecino $w$ de $0$ a $10$ viendo la pregunta, $T_{\text{sem}}$, el
  triplete de enlace y un resumen $C(w)$ del vecino, generado por el
  propio modelo a partir de sus tripletes. La puntuación reescala el peso
  estático,
  \begin{equation}
  \begin{aligned}
  \hat{w}_{uw} \;&=\; \phi\bigl(\op{LLM}(q,u,w,C(w))\bigr)\cdot w_{uw},\\
  \phi:\ &0\text{--}3\mapsto 0,\quad 4\text{--}6\mapsto 0{,}2\text{--}0{,}3,\quad
  7\text{--}9\mapsto 2\text{--}3,\quad 10\mapsto 5,
  \end{aligned}
  \label{eq:catrag-aristas}
  \end{equation}
  la fila de $u$ se renormaliza y se ejecuta \emph{un solo} paseo sobre la
  matriz resultante $P_q$. La modulación es asimétrica ---solo aristas
  salientes de las semillas--- y de \emph{un salto}: no se expande a los
  vecinos de los vecinos. Sobre la fila real anterior, un $8$ para
  Michael Curtiz ($\times 2{,}5$), un $5$ para los tres intérpretes
  ($\times 0{,}25$) y $\le 3$ para el resto dejarían
  $Z_u = 5 + 1{,}5 + 1 = 7{,}5$ y $P_q(u\to\text{curtiz}) = 0{,}667$
  frente a $0{,}105$: el mismo grafo, otra matriz para esta pregunta. El
  vecino no recibe masa de reinicio; recibe más probabilidad de ser
  cruzado cuando el paseo está en la semilla, a costa de sus hermanos, y
  con él todo lo que hay detrás de esa arista.
\item \textbf{Refuerzo de pasajes con hechos verificados} (palanca $P$).
  La arista semilla$\to$pasaje que respalda un triplete aprobado se
  amplifica,
  $\hat{w}_{up} = w_{up}\,\bigl(1+\beta\,\mathbf{1}[(u,p)\text{ respaldada
  por } T_{\text{sem}}]\bigr)$ con $\beta=2{,}5$.
\end{enumerate}
El coste es de hasta cinco llamadas por consulta (una por semilla, con
sus $15$ vecinos en lote) más los resúmenes de nodo, y el propio artículo
reconoce que el sistema resulta más caro que la recuperación densa. Su
ablación en 2Wiki mide la aportación de cada mecanismo en exhaustividad a
5: $87{,}0$ el sistema completo, $86{,}1$ sin anclas, $85{,}6$ sin
reponderación y $88{,}4$ sin refuerzo de pasajes (este último, por tanto,
resta en 2Wiki); en cadena completa, $66{,}1 \to 67{,}6$ sobre \hippo{}.

\paragraph{Por qué la ganancia es acotada: alcance de un salto.}
El alcance de la reponderación es \guillemotleft{}semillas más una arista\guillemotright{}, y
\guillemotleft{}una arista desde las semillas\guillemotright{} no equivale a \guillemotleft{}el primer
salto de la pregunta\guillemotright{}: depende de quién sea semilla. Si el filtro
aprobó el triplete puente, la entidad puente ya es semilla ---como objeto
del triplete, ecuación~\ref{eq:siembra-hippo}--- y sus aristas, el
segundo salto de la pregunta, también se reponderan; el salto lo dio el
filtro y la reponderación solo lo prolonga una arista. Si el filtro
\emph{no} lo aprobó ---lo que ocurre en el ejemplo anterior, donde Curtiz
no figura en ningún triplete aprobado---, la reponderación puede
amplificar película$\to$Curtiz y el paseo llega a él, pero su fila sigue
siendo la estática ($Z_u=103$, biografía a $1/103$) y el refuerzo de
pasaje tampoco le alcanza, porque ningún triplete verificado respalda esa
arista: la cadena se queda a una arista de completarse. En términos de la
ecuación~\eqref{eq:neumann}, los tres mecanismos de CatRAG actúan sobre
el coeficiente o añaden fuentes débiles; la \emph{fuente fuerte} del
segundo salto sigue siendo el filtro de \hippo{} operando sobre tripletes
planos y entidades fragmentadas.

\begin{table}[t]
\centering\footnotesize\setlength{\tabcolsep}{3.5pt}
\caption{Por dónde entra la pregunta en el paseo con reinicio
(ecuación~\ref{eq:neumann}) en cada sistema de la línea. \guillemotleft{}Fuente\guillemotright{}
es el vector de reinicio $\mathbf{v}$; \guillemotleft{}coeficiente\guillemotright{}, la matriz de
transición $P$.}
\label{tab:palancas}
\begin{tabular}{lp{0.40\textwidth}p{0.38\textwidth}}
\toprule
sistema & fuente ($\mathbf{v}$) & coeficiente ($P$)\\
\midrule
HippoRAG & entidades de la pregunta (NER), fuerte & fija en indexación\\
\addlinespace
\hippo{} & sujeto y objeto de los tripletes aprobados por el filtro
  ($s_q/|P_\varepsilon|$) $+$ todos los pasajes ($0{,}05$) & fija en
  indexación\\
\addlinespace
CatRAG & lo de \hippo{} $+$ entidades de la pregunta como anclas débiles
  ($0{,}2/|P_e|$) & reponderada por juez LLM en las aristas salientes de
  $5$ semillas (un salto) $+$ refuerzo $\times 3{,}5$ de las aristas
  semilla$\to$pasaje verificadas\\
\addlinespace
\sys{} & entidades de los hechos aprobados a peso $1$ (compuerta dura)
  $+$ entidades enlazadas ($0{,}5$) $+$ \emph{pasaje propio} de cada una
  ($0{,}5\,\mathbf{v}$) $+$ anclas débiles ($0{,}05$) $+$ todos los
  pasajes ($0{,}05$) & modulada globalmente por afinidad a la entrada de
  \emph{todos} los nodos-hecho, $\mu_a\in[0{,}25,\,1]$, sin llamadas
  (ec.~\ref{eq:modulacion})\\
\bottomrule
\end{tabular}
\end{table}

\paragraph{Lo que la evolución deja abierto.}
Tres observaciones se siguen de este recorrido y fijan el punto de partida
de la sección~\ref{sec:metodos}. (i)~En toda la línea el segundo salto es
cuestión de \emph{fuente}: lo da el filtro al sembrar el objeto del
triplete aprobado, y ningún coeficiente lo sustituye cuando el filtro
falla. (ii)~Más allá de una arista de las semillas gobierna el grafo
estático, con sus concentradores y sus aristas genéricas; reponderar por
juez tiene un coste por arista que impide extender el alcance. (iii)~La
identidad de las entidades se resuelve en consulta con aristas de
sinonimia por umbral, de modo que cada eslabón de la cadena es una nueva
ocasión de quedar en una \guillemotleft{}isla\guillemotright{}. \sys{} responde a las tres:
resuelve la identidad en indexación (sección~\ref{sec:diccionario});
promueve el hecho a nodo, de modo que es a la vez un peaje modulable por
consulta y un punto donde converge la evidencia de todos sus
participantes (sección~\ref{sec:asercion}); modula el coeficiente
\emph{globalmente} y sin llamadas, por afinidad consulta--hecho
(sección~\ref{sec:ppr}); y hace de la siembra la palanca del segundo
salto: la entidad puente descubierta por el filtro siembra a peso
completo y su pasaje propio recibe masa directa
(sección~\ref{sec:siembra}), de modo que la cadena no depende de que el
paseo la recorra.

\subsection{Relación con el prototipo clínico GeSIDA y con CatRAG}
\label{sec:relacion}

Este trabajo no parte de cero: el sistema evaluado es la generalización de
un recuperador en producción dentro de un asistente clínico sobre las guías
GeSIDA de VIH (sección~\ref{sec:gesida}), que a su vez adoptó y extendió los
mecanismos de consulta de la línea HippoRAG~2/CatRAG. La genealogía es
explícita (sección~\ref{sec:palancas}): del lado CatRAG hereda las
anclas simbólicas débiles, la modulación de transiciones por consulta
---aquí global y sin juez por arista--- y la idea de reforzar pasajes
clave ---aquí trasladada de la arista a la siembra---; del lado
HippoRAG~2, la siembra densa de todos los pasajes, el filtro de memoria
de reconocimiento y los diales del paseo. Las aportaciones propias
están en la \emph{representación} ---el grafo de aserciones reificadas
$n$-arias en lugar de tripletes planos--- y en la \emph{resolución de
entidades} ---una memoria canónica construida en indexación en lugar de
aristas de sinonimia en consulta---, además del anclaje estructural de la
pregunta (enlazador, frontera, pasaje propio) que esa representación hace
posible. En el caso clínico, el papel de la memoria canónica lo desempeña
una terminología externa (SNOMED~CT vía UMLS); en el caso abierto de
2Wiki, donde no existe terminología, el diccionario se deriva del propio
corpus. Ambas instancias del mismo patrón se comparan en la
tabla~\ref{tab:gesida-wiki}.

\subsection{Contribuciones}

\begin{enumerate}[leftmargin=1.6em, itemsep=2pt]
\item \textbf{Representación}: un grafo de conocimiento cuya unidad es la
  \emph{aserción reificada $n$-aria} (sección~\ref{sec:asercion}) ---el
  hecho como nodo con frase autocontenida, calificadores, procedencia y un
  rol por participante--- frente al triplete-arista de la línea HippoRAG.
\item \textbf{Resolución de entidades en indexación}: una \emph{memoria
  canónica de entidades} derivada del corpus (fichas por mención,
  candidatos por dos canales vectoriales, restricciones duras y
  adjudicación por modelo de lenguaje) que elimina las aristas de sinonimia
  y sus fallos de umbral (sección~\ref{sec:diccionario}).
\item \textbf{Mecanismos de consulta estructurales} sobre esa
  representación: enlazador de menciones contra el diccionario, candidatos
  de frontera para el filtro de hechos, compuerta dura en la siembra y masa
  de reinicio directa en el \emph{pasaje propio} de cada entidad
  descubierta, más un enrutador léxico de intención con un especialista de
  descomposición (sección~\ref{sec:consulta}).
\item \textbf{Evaluación controlada}: paridad estricta con HippoRAG~2 (su
  código oficial, mismos modelos, diales pre-registrados, separación
  rigurosa de datos) y comparación pareada por pregunta con arranque
  aleatorio (\emph{bootstrap}); mejora significativa de $+24{,}8$ puntos de
  cadena completa a 5 pasajes y estado del arte publicado en 2Wiki con
  modelos pequeños (sección~\ref{sec:resultados}).
\item \textbf{Caso de uso clínico}: la instanciación del mismo patrón en un
  asistente sobre las guías GeSIDA de VIH, con anclaje terminológico
  SNOMED~CT/UMLS, capas de verificación deóntica y numérica y curación
  humana (sección~\ref{sec:gesida}).
\end{enumerate}

% =========================================================================
\section{Material y métodos}
\label{sec:metodos}

\subsection{Visión general}
\label{sec:vision}

El sistema propuesto (\sys) es un recuperador de pasajes para preguntas
multisalto basado en un grafo de conocimiento construido automáticamente a
partir del corpus. Su diseño responde a tres objetivos: (i)~recuperar la
\emph{cadena completa} de pasajes de evidencia, no solo el más afín a la
pregunta; (ii)~operar con modelos pequeños y baratos
(\texttt{gpt-4o-mini} como modelo de lenguaje y
\texttt{text-embedding-3-small}, de dimensión $1536$, como codificador
vectorial), con aproximadamente una llamada al modelo de lenguaje por
consulta; y (iii)~mantener paridad experimental estricta con el sistema de
referencia (\hippo{}), que se ejecuta con su propio código y con los mismos
modelos.

El sistema opera en dos fases (figura~\ref{fig:pipeline}). La fase de
\emph{indexación}, ejecutada una sola vez por corpus y sin acceso a las
preguntas, comprende tres pasos: extracción abierta de \emph{aserciones
reificadas} (sección~\ref{sec:openie}), construcción de una \emph{memoria
canónica de entidades} (sección~\ref{sec:diccionario}) y materialización del
\emph{grafo canónico} (sección~\ref{sec:grafo}). La fase de \emph{consulta}
encadena seis pasos: enrutado léxico de la intención, enlazado de las
menciones de la pregunta contra la memoria de entidades, cálculo de afinidad
pregunta--hecho, filtro selector de hechos mediante modelo de lenguaje,
construcción del vector de personalización (siembra) y un paseo aleatorio con
reinicio personalizado (\emph{Personalized PageRank}, PPR) cuyas transiciones
se modulan por la consulta (secciones~\ref{sec:router}--\ref{sec:ppr}).

En todo el apartado, $e(\cdot)\in\R^{1536}$ denota la incrustación vectorial
normalizada de un texto ($\norm{e(x)}_2=1$), de modo que la similitud coseno
es el producto escalar $\cos(x,y)=\langle e(x),e(y)\rangle$. El corpus es
$\mathcal{D}=\{c_1,\dots,c_N\}$ con $N=6\,119$ pasajes (título + texto).
La tabla~\ref{tab:componentes} inventaría los componentes de extremo a
extremo con sus entradas, salidas y coste; las secciones siguientes los
desarrollan en el orden en que operan.

\begin{table}[t]
\centering\footnotesize\setlength{\tabcolsep}{3.5pt}
\caption{Inventario de componentes de \sys{} de extremo a extremo. El coste
LLM se expresa en llamadas al modelo de lenguaje; ``emb.''\ son llamadas al
codificador vectorial (por lotes, coste marginal). Todas las llamadas se
guardan en caché persistente, por lo que re-ejecutar no re-factura.}
\label{tab:componentes}
\setlength{\tabcolsep}{3pt}
\begin{tabular}{lllll}
\toprule
fase & componente & entrada & salida & coste\\
\midrule
index. & extractor de aserciones (\S\ref{sec:openie}) & pasaje & $E_c$ + $A_c$ & 1 llamada/pasaje\\
index. & constructor de fichas (\S\ref{sec:diccionario}) & menciones & ficha por mención & 2 emb./mención\\
index. & generador de candidatos (\S\ref{sec:diccionario}) & fichas & pares candidatos & 0\\
index. & adjudicador de pares (\S\ref{sec:diccionario}) & pares ambiguos & veredictos & 1 llamada/8 pares\\
index. & consolidador canónico (\S\ref{sec:diccionario}) & veredictos & memoria $\mathcal{E}$ (\S\ref{sec:memoria}) & 1 emb./entrada\\
index. & materializador del grafo (\S\ref{sec:grafo}) & aserciones + dicc. & grafo $G$ & 0\\
index. & incrustadores (\S\ref{sec:dual}) & $\tau_a$, fichas, pasajes & 3 almacenes vect. & 1 emb./objeto\\
\midrule
consulta & enrutador léxico (\S\ref{sec:router}) & pregunta & tipo + estrategia & 0\\
consulta & enlazador de menciones (\S\ref{sec:linker}) & pregunta + dicc. & entidades $L(q)$ & 0--1 emb.\\
consulta & descomponedor (\S\ref{sec:afinidad}) & pregunta (compar.) & subconsultas & 1 llamada\\
consulta & afinidad consulta--hecho (\S\ref{sec:afinidad}) & vectores de consulta & $\sigma_a$ por aserción & 1 emb./consulta\\
consulta & filtro selector (\S\ref{sec:filtro}) & candidatos $\mathcal{K}(q)$ & aprobados $S(q)$ & 1 llamada\\
consulta & siembra (\S\ref{sec:siembra}) & $S(q)$, $L(q)$, denso & vector $\mathbf{v}$ & 0\\
consulta & PPR modulado (\S\ref{sec:ppr}) & $G$, $\mathbf{v}$, $\tilde\sigma$ & pasajes ordenados & 0 ($\sim$ms)\\
\bottomrule
\end{tabular}
\end{table}

\begin{figure}[t]
\centering
\begin{tikzpicture}[
  font=\footnotesize,
  caja/.style={draw, rounded corners=2pt, align=center, minimum height=8.5mm, inner sep=3pt, fill=gray!8},
  dato/.style={draw, dashed, rounded corners=2pt, align=center, inner sep=3pt},
  fl/.style={-{Stealth[length=2.2mm]}, semithick},
  node distance=3mm and 5mm]
  % ---- fase de indexacion
  \node[dato] (corpus) {Corpus\\ $6\,119$ pasajes};
  \node[caja, right=of corpus] (openie) {Extracción\\ abierta (LLM)\\ \S\ref{sec:openie}};
  \node[caja, right=of openie] (dicc) {Resolución de\\ entidades\\ \S\ref{sec:diccionario}};
  \node[dato, right=of dicc] (mem) {Memoria\\ canónica\\ \S\ref{sec:memoria}};
  \node[caja, right=of mem] (grafo) {Grafo canónico\\ $G$ \S\ref{sec:grafo}};
  \node[dato, right=of grafo] (emb) {Incrustaciones\\ (3 almacenes)};
  \draw[fl] (corpus) -- (openie);
  \draw[fl] (openie) -- (dicc);
  \draw[fl] (dicc) -- (mem);
  \draw[fl] (mem) -- (grafo);
  \draw[fl] (grafo) -- (emb);
  % ---- fase de consulta: dos ramas independientes que convergen en el filtro
  \node[dato, below=25mm of corpus] (q) {Pregunta $q$};
  \node[caja, right=of q] (router) {Enrutador\\ léxico \S\ref{sec:router}};
  \node[caja, above right=6mm and 5mm of router.east, anchor=west] (linker) {Enlazador\\ $L(q)$ \S\ref{sec:linker}};
  \node[caja, below right=6mm and 5mm of router.east, anchor=west] (afin) {Afinidad\\ $\{v_j\},\ \sigma$ \S\ref{sec:afinidad}};
  \node[caja, anchor=west] (filtro) at ($(linker.east)!0.5!(afin.east)+(5mm,0)$) {Filtro selector\\ (LLM) $S(q),\tilde\sigma$\\ \S\ref{sec:filtro}};
  \node[caja, right=of filtro] (siembra) {Siembra\\ $\mathbf{v}$ \S\ref{sec:siembra}};
  \node[caja, right=of siembra] (ppr) {PPR\\ modulado \S\ref{sec:ppr}};
  \node[dato, below=3mm of ppr] (out) {pasajes\\ ordenados};
  \draw[fl] (q) -- (router);
  \draw[fl] (router.east) -- (linker.west);
  \draw[fl] (router.east) -- (afin.west);
  \draw[fl] (linker.east) -- (filtro.west);
  \draw[fl] (afin.east) -- (filtro.west);
  \draw[fl] (filtro) -- (siembra);
  \draw[fl] (siembra) -- (ppr);
  \draw[fl] (ppr) -- (out);
  \draw[fl, gray] (linker.north) to[out=15,in=165] node[above, font=\tiny, gray, pos=0.35]{$L(q)$} (siembra.north);
  \draw[fl, gray] (afin.south) to[out=-15,in=-165] node[below, font=\tiny, gray, pos=0.35]{$\{v_j\},\sigma$} (siembra.south);
  % etiquetas de fase
  \begin{scope}[on background layer]
    \node[fit=(corpus)(emb), draw=gray!60, rounded corners=3pt, inner sep=4pt,
          label={[gray, font=\scriptsize]above left:{indexación (una vez por corpus)}}] {};
    \node[fit=(q)(linker)(afin)(ppr)(out), draw=gray!60, rounded corners=3pt, inner sep=4pt,
          label={[gray, font=\scriptsize]above left:{consulta (en línea, $\sim$1 llamada LLM)}}] {};
  \end{scope}
  % conexion entre fases
  \draw[fl, gray] (grafo.south) to[out=-90,in=90] (ppr.north);
  \draw[fl, gray] (mem.south) to[out=-90,in=90] (linker.north);
\end{tikzpicture}
\caption{Arquitectura de \sys{} en dos fases. La fase de indexación no ve
ninguna pregunta. La fase de consulta no es una cadena lineal: tras el
enrutador, la rama \emph{simbólica} (enlazador) y la rama \emph{semántica}
(afinidad) leen solo la pregunta y son independientes entre sí; convergen
en el filtro, y la siembra recibe los productos de las tres (las flechas
grises marcan que $L(q)$ y $\{v_j\},\sigma$ llegan también directamente a
la siembra). Se realiza aproximadamente una llamada al modelo de lenguaje
(el filtro; dos en comparación, donde se añade la descomposición). La
memoria canónica es un artefacto de datos (caja discontinua) que consumen
tanto el grafo como el enlazador (sección~\ref{sec:memoria}).}
\label{fig:pipeline}
\end{figure}

% =========================================================================
\subsection{La unidad semántica: aserción reificada frente a triplete}
\label{sec:asercion}

La diferencia de representación más profunda entre \sys{} y \hippo{} está en
la unidad semántica del grafo, por lo que la definimos antes de describir el
proceso constructivo.

\paragraph{Triplete (\hippo{}).}
\hippo{} extrae de cada pasaje tripletes $(s,r,o)$ ---sujeto, relación,
objeto--- y los materializa como \emph{aristas} entre nodos de frase
(entidades): el hecho \guillemotleft{}Atomised fue dirigida por Oskar Roehler\guillemotright{} es la arista
$\text{Atomised} \xrightarrow{\;directed\;} \text{Oskar Roehler}$. Esta
representación tiene tres límites estructurales. Primero, es estrictamente
binaria: un hecho con tres o más participantes (\guillemotleft{}la canción fue escrita por
$A$ y $B$ para la película $C$\guillemotright{}) debe romperse o perder participantes.
Segundo, los calificadores (fecha, lugar) no tienen sitio natural: o se
pierden o generan nodos literales. Tercero, el hecho no es un objeto
direccionable: no se le puede asociar procedencia (?`qué pasaje lo afirma?) ni
una incrustación propia sin estructuras auxiliares.\footnote{\hippo{}
mantiene de hecho un almacén paralelo de incrustaciones de tripletes para la
puntuación consulta--hecho, precisamente porque la arista no es un nodo.}

\paragraph{Aserción reificada (\sys{}).}
Una \emph{aserción} es un hecho atómico promovido a nodo de primera clase:
\begin{equation}
a = \bigl(s_a,\; r_a,\; o_a,\; Q_a,\; \tau_a\bigr),
\label{eq:asercion}
\end{equation}
donde $s_a$ y $o_a$ son las menciones de sujeto y objeto, $r_a$ una etiqueta
corta de relación (\texttt{DIRECTED}, \texttt{DATE\_OF\_BIRTH},~\dots),
$Q_a$ un diccionario de calificadores (fecha, lugar) y $\tau_a$ una
\emph{frase autocontenida} que verbaliza el hecho con nombres completos y sin
pronombres. El nodo $a$ se conecta mediante aristas tipadas con su pasaje de
origen (procedencia), con su sujeto, con su objeto y con \emph{cualquier otro
participante} mencionado en $\tau_a$ (reificación $n$-aria,
sección~\ref{sec:grafo}). La figura~\ref{fig:reificacion} contrasta ambas
representaciones sobre un hecho real del corpus.

Dos consecuencias operativas justifican el coste de la reificación:
\begin{enumerate}[leftmargin=1.5em, itemsep=1pt]
\item \textbf{Lo que se incrusta es la frase, no la etiqueta.} Con
  \emph{etiqueta} nos referimos al identificador categórico de relación
  $r_a$ (\texttt{DIRECTED}, \texttt{MOTHER\_OF}, \texttt{DATE\_OF\_BIRTH}):
  un vocabulario corto y ruidoso que el extractor asigna a cada hecho. En un
  sistema de tripletes, la representación vectorial del hecho se construye
  sobre la terna $(s,r,o)$ ---y hereda ese ruido---. Aquí, en cambio, la
  afinidad consulta--hecho (sección~\ref{sec:afinidad}) se calcula sobre
  $e(\tau_a)$, la incrustación de la frase autocontenida; la etiqueta $r_a$
  se conserva solo como metadato de auditoría y \emph{no interviene en
  ninguna puntuación}. La diferencia importa en la práctica: en el corpus
  real el extractor produce en ocasiones etiquetas invertidas ---la
  aserción real \texttt{Oskar Roehler::a6} lleva $r_a=$ \texttt{MOTHER\_OF}
  con el hijo como sujeto---, pero su frase
  $\tau_a$ (\guillemotleft{}Oskar Roehler is the son of Gisela Elsner\guillemotright{}) conserva la
  semántica correcta, que es la que ven tanto el coseno como el filtro
  selector. El sistema es así robusto a la parte más frágil de la
  extracción abierta (la normalización del vocabulario de relaciones) sin
  renunciar a la estructura.
\item \textbf{El hecho es una unión ponderable donde converge la
  evidencia.} En el grafo de búsqueda (sección~\ref{sec:grafo}) los
  participantes de un hecho apuntan \emph{hacia} su nodo-aserción y este
  desagua en su pasaje de procedencia. El nodo-hecho cumple así dos
  funciones que la arista anónima del grafo plano no puede cumplir: es un
  \emph{peaje} modulable por consulta ---entrar en el hecho es tanto más
  probable cuanto más afín es a la pregunta,
  ecuación~\ref{eq:modulacion}--- y es un \emph{punto de convergencia}: el
  pasaje que afirma un hecho aprobado recibe masa de \emph{todos} sus
  participantes a la vez, una señal de evidencia conjunta que una arista
  binaria no puede agregar. El cruce hacia el segundo participante (la
  entidad puente) no lo realiza el paseo sino la capa semántica: cuando el
  filtro aprueba un hecho, \emph{todas} sus entidades siembran
  (sección~\ref{sec:siembra}); el hecho actúa de puente en la siembra y de
  peaje en el paseo (la sección~\ref{sec:ppr} lo hace exacto).
\end{enumerate}

\begin{figure}[t]
\centering
\begin{tikzpicture}[
  font=\footnotesize,
  ent/.style={draw, ellipse, align=center, inner sep=2pt},
  aser/.style={draw, rectangle, rounded corners=1pt, align=center, inner sep=3pt, fill=gray!12},
  chk/.style={draw, rectangle, align=center, inner sep=3pt, dashed},
  fl/.style={-{Stealth[length=2mm]}, semithick},
  node distance=6mm and 8mm]
  % ------- izquierda: triplete plano
  \node[ent] (t1) {Atomised};
  \node[ent, right=14mm of t1] (t2) {Oskar\\ Roehler};
  \draw[fl] (t1) -- node[above, font=\scriptsize]{\emph{directed}} (t2);
  \node[ent, below=7mm of t1, densely dotted] (t3) {Atomised\\ (film)};
  \draw[gray, densely dotted] (t1) -- node[left, font=\scriptsize]{sinonimia} (t3);
  \node[below=2mm of t3, font=\scriptsize\itshape, align=center, xshift=8mm]
       {(a) triplete plano (\hippo{}):\\ el hecho es una arista};
  % ------- derecha: asercion reificada
  \node[aser, right=32mm of t2] (a) {aserción $a_1$\\[-1pt]
        {\scriptsize$\tau$: \guillemotleft{}Atomised was directed}\\[-2pt]
        {\scriptsize by Oskar Roehler.\guillemotright{}}};
  \node[ent, above left=5mm and 2mm of a] (e1) {can:atomised};
  \node[ent, above right=5mm and 2mm of a] (e2) {can:oskar roehler};
  \node[chk, below=6mm of a] (c1) {pasaje\\ \guillemotleft{}Atomised (film)\guillemotright{}};
  \draw[fl] (a) -- node[left, font=\scriptsize, pos=0.4]{\textsc{sujeto}} (e1);
  \draw[fl] (a) -- node[right, font=\scriptsize, pos=0.4]{\textsc{objeto}} (e2);
  \draw[fl] (a) -- node[right, font=\scriptsize]{\textsc{en\_pasaje}} (c1);
  \node[below=2mm of c1, font=\scriptsize\itshape, align=center]
       {(b) aserción reificada (\sys{}): el hecho es un nodo\\
        con frase, calificadores, procedencia y roles $n$-arios};
\end{tikzpicture}
\caption{El mismo hecho del corpus en ambas representaciones. En (a) las
variantes de nombre exigen aristas de sinonimia; en (b) las menciones ya
están resueltas al nodo canónico (\texttt{can:atomised} agrupa \guillemotleft{}Atomised\guillemotright{} y
\guillemotleft{}Atomised (film)\guillemotright{}) y el hecho es direccionable: tiene incrustación propia
$e(\tau)$, procedencia y tantos roles como participantes.}
\label{fig:reificacion}
\end{figure}

% =========================================================================
\subsection{Fase de indexación}
\label{sec:indexacion}

Las etapas de indexación forman una cadena de refinamiento con productos
explícitos: la \emph{extracción} (sección~\ref{sec:openie}) convierte cada
pasaje en menciones y aserciones; la \emph{memoria canónica}
(sección~\ref{sec:diccionario}) convierte las menciones en entidades con
identidad; el \emph{grafo} (sección~\ref{sec:grafo}) materializa
aserciones y entidades sobre los pasajes; y los \emph{incrustadores}
(sección~\ref{sec:dual}) añaden la capa vectorial alineada. Cada etapa
consume exactamente la salida de la anterior (tabla~\ref{tab:componentes})
y ninguna ve una sola pregunta.

\subsubsection{Extracción abierta de aserciones reificadas}
\label{sec:openie}

Cada pasaje $c\in\mathcal{D}$ se procesa con una única llamada al modelo de
lenguaje (temperatura $0$, salida JSON estricta) que devuelve (i)~la lista de
menciones de entidad del pasaje $E_c$ (personas, obras, lugares,
organizaciones, en la forma más completa que dé el texto) y (ii)~la lista de
aserciones $A_c$ con el esquema de la ecuación~\eqref{eq:asercion}. La
consigna dada al modelo (\emph{prompt}) exige extraer \emph{todos} los
hechos (relaciones familiares, fechas y
lugares de nacimiento y muerte, vínculos creador--obra, fechas de estreno,
nacionalidad, ocupación) y que cada $\tau_a$ sea una frase autocontenida.
Todas las respuestas se almacenan en una caché persistente indexada por
(modelo, versión de consigna, pasaje), de modo que la extracción es reanudable
y reproducible sin coste adicional.

Sobre el corpus de evaluación la extracción produce $47\,999$ aserciones
($7{,}8$ por pasaje) y $38\,934$ menciones de entidad, sin errores de
formato. Como diagnóstico del \emph{techo} de la representación se midió la
cobertura de los $508$ tripletes oro del conjunto de sondeo
(sección~\ref{sec:protocolo}) mediante solape de vocabulario ($\geq 60\%$ de
los términos): el $91{,}1\%$ de los tripletes oro queda capturado con sujeto
y objeto en una \emph{misma} aserción (criterio estricto) y el $98{,}8\%$
con ambos en el conjunto extraído (criterio laxo). La extracción no es, por
tanto, el cuello de botella del sistema.

\subsubsection{Memoria canónica de entidades}
\label{sec:diccionario}

Los grafos extraídos automáticamente sufren la \emph{fragmentación de
entidades}: la misma persona aparece como \guillemotleft{}H.\,P. Lovecraft\guillemotright{} en un pasaje y
como \guillemotleft{}Howard Phillips Lovecraft\guillemotright{} en su biografía, y ambas menciones acaban
en nodos distintos. La solución habitual ---aristas de sinonimia por umbral de
coseno en consulta, como en \hippo{}--- es frágil justo donde más se
necesita: el coseno entre esas dos formas es $0{,}781$, por debajo del
umbral estándar $0{,}80$, y el paseo aleatorio se corta en la
\guillemotleft{}isla\guillemotright{}. En el análisis de errores de la versión previa del sistema, el
$57\%$ de los fallos de cadena se atribuyó a este fenómeno.

\sys{} lo resuelve \emph{antes de construir el grafo} con una memoria
canónica de entidades (un diccionario), construida una única vez desde el
propio corpus ---nunca desde los datos de evaluación--- y sin curación
humana. Es un problema clásico de \emph{resolución de entidades}: decidir,
para cada par de menciones, si nombran el mismo objeto del mundo real, y
consolidar las que sí en una entrada con identidad. La
figura~\ref{fig:diccionario} muestra el flujo completo con las cifras reales
del corpus de evaluación; las secciones que siguen recorren cada paso con
trazas instrumentadas y terminan con la atribución de cada fusión a su
mecanismo, la contribución medida al resultado y las limitaciones
conocidas.

\paragraph{Vocabulario.}
Como el apartado usa términos de resolución de entidades que no son de uso
general, los fijamos antes. Una \emph{forma} es la cadena con que un pasaje
nombra algo (\guillemotleft{}Oskar Roehler\guillemotright{}, \guillemotleft{}Klaus Roehler\guillemotright{}). Una \emph{ficha} es la
unidad de trabajo: un par (pasaje, forma) con su contexto; la misma persona
nombrada en dos pasajes produce dos fichas. Un \emph{canal} es una de las dos
representaciones vectoriales de la ficha (nombre o contexto). Las
\emph{vecinas} de una ficha son las fichas más próximas por coseno en un
canal. Un \emph{candidato} es un par de fichas que pasa el filtro barato y
merece ser juzgado: el paso de candidatos no decide nada, decide \emph{qué
se decide}. Una \emph{banda} es el destino de un candidato: fusión
automática, adjudicación por el modelo de lenguaje o descarte. Un
\emph{grupo} es un conjunto de fichas unidas por cierre transitivo, y una
\emph{entrada canónica} es la entidad que resume ese grupo. Una
\emph{colisión} es lo contrario de una fusión: dos entradas \emph{distintas}
cuyo nombre canónico coincide (dos personas llamadas John Middleton), que
reciben identificadores diferenciados.

\begin{figure}[p]
\centering
\begin{tikzpicture}[
  font=\scriptsize,
  paso/.style={draw, rounded corners=2pt, align=center, inner sep=3pt, fill=gray!8, text width=62mm},
  dato/.style={draw, dashed, rounded corners=2pt, align=center, inner sep=3pt, text width=62mm},
  rama/.style={draw, rounded corners=2pt, align=center, inner sep=2.5pt, text width=29mm, minimum height=10mm},
  fin/.style={draw, densely dotted, rounded corners=2pt, align=center, inner sep=2.5pt, text width=29mm, minimum height=10mm, text=black!70},
  nota/.style={font=\tiny, text=black!75, align=left, text width=40mm},
  fl/.style={-{Stealth[length=2mm]}, semithick},
  et/.style={font=\tiny, fill=white, inner sep=1pt},
  node distance=4mm and 5mm]
  % ------------------------------------------------ columna central
  \node[dato] (corpus) {\textbf{Entrada} (\S\ref{sec:openie}): $6\,119$ pasajes
    $\to$ $38\,934$ menciones de entidad y $47\,999$ aserciones con frase $\tau_a$};
  \node[paso, below=of corpus] (fichas) {\textbf{Paso 1 --- una ficha por (pasaje, forma)}\\
    forma $f$; descripción $d_m$ = frases $\tau_a$ del pasaje en que participa
    ($\le 400$ car.); años de nacimiento/muerte; ?`es el título?\\
    $+\,1\,342$ menciones sintéticas de título $\;\Rightarrow\;$ \textbf{$40\,276$ fichas}};
  \node[rama, below left=4mm and 1mm of fichas.south] (nu) {canal \textbf{nombre}\\ $\nu_m = e(f)$};
  \node[rama, below right=4mm and 1mm of fichas.south] (phi) {canal \textbf{contexto}\\ $\phi_m = e(f \oplus d_m)$};
  \coordinate (mid) at ($(nu.south)!0.5!(phi.south)$);
  \node[paso, below=4mm of mid] (cand) {\textbf{Paso 2 --- candidatos}\\
    fila completa de cosenos por bloques; $15$ vecinas por canal;
    se excluye el mismo pasaje; umbrales $\nu \ge 0{,}80$ ó $\phi \ge 0{,}75$\\
    $811$ millones de pares posibles $\;\Rightarrow\;$ \textbf{$47\,063$ pares examinados}};
  \node[paso, below=of cand] (restr) {\textbf{Paso 3a --- restricciones duras}\\
    fechas de nacimiento o muerte disjuntas; ordinales nobiliarios o regnales
    distintos $\;\Rightarrow\;$ el par se descarta sin más};
  \node[rama, below=of restr] (llm) {\textbf{banda LLM}\\ $9\,772$ pares, lotes de 8, caché\\ $2\,322$ sí / $7\,450$ no};
  \node[rama, left=3mm of llm] (auto) {\textbf{fusión automática}\\ nombre base idéntico, o\\ $\nu\ge0{,}90 \wedge \phi\ge0{,}70$\\ $32\,500$ pares};
  \node[fin, right=3mm of llm] (desc) {\textbf{descarte}\\ ninguna regla; o coseno alto pero fechas u ordinal incompatibles};
  \node[paso, below=6mm of llm] (cierre) {\textbf{Paso 3c --- cierre transitivo} (unión--búsqueda)\\
    todo par aceptado une sus grupos; el mismo pasaje nunca une por coseno};
  \node[rama, left=3mm of cierre] (tit) {\textbf{Paso 3b --- título$\to$sujeto}\\ $2\,848$ uniones ($2\,431$ nuevas)};
  \node[paso, below=of cierre] (entrada) {\textbf{Paso 4 --- entrada canónica por grupo}\\
    nombre (preferencia: forma-título), alias, descripción fusionada, pasajes,
    pasaje propio $\pi(\varepsilon)$; mapa (pasaje, forma) $\to$ id; incrustación de la ficha fusionada\\
    \textbf{$26\,960$ entradas}: $21\,144$ con 1 mención, $5\,816$ con $>1$, $36$ colisiones};
  \node[dato, below=of entrada] (salida) {\textbf{Salida}: grafo canónico sin aristas de sinonimia (\S\ref{sec:grafo});
    tabla de alias del enlazador (\S\ref{sec:linker}); pasaje propio para la siembra (\S\ref{sec:siembra})};
  % ------------------------------------------------ flechas
  \draw[fl] (corpus) -- (fichas);
  \draw[fl] (fichas.south) -- (nu.north);
  \draw[fl] (fichas.south) -- (phi.north);
  \draw[fl] (nu.south) -- (cand.north);
  \draw[fl] (phi.south) -- (cand.north);
  \draw[fl] (cand) -- (restr);
  \draw[fl] (restr.south) -- (llm.north);
  \draw[fl] (restr.south) -- (auto.north);
  \draw[fl] (restr.south) -- (desc.north);
  \draw[fl] (auto.south) -- (cierre.north west);
  \draw[fl] (llm.south) -- node[et, right]{solo los \guillemotleft{}sí\guillemotright{}} (cierre.north);
  \draw[fl] (tit.east) -- (cierre.west);
  \draw[fl] (cierre) -- (entrada);
  \draw[fl] (entrada) -- (salida);
  % ------------------------------------------------ notas laterales
  \node[nota, right=3mm of fichas] {\emph{Ejemplo}: \texttt{Atomised (film)||oskar roehler}
    (desc: \guillemotleft{}Atomised was directed by Oskar Roehler. Atomised was written by
    Oskar Roehler.\guillemotright{}) y \texttt{Oskar Roehler||oskar roehler} (título; nac.\ 1959;
    desc con sus 11 hechos) son dos fichas distintas de la misma persona.};
  \node[nota, right=3mm of cand] {De $74\,580$ pares sobre umbral, el tope de 15
    deja fuera $27\,517$, todos de concentradores (\guillemotleft{}France\guillemotright{}: $4\,809$), que
    el cierre transitivo recupera. El canal contexto gasta $4{,}9$ de sus 15 huecos
    en fichas del mismo pasaje (el canal nombre, $0{,}6$).};
  \node[nota, right=3mm of desc] {\emph{Ejemplos}: dos \guillemotleft{}William Ferguson\guillemotright{}
    ($\nu = 1{,}000$; nacidos en 1809 y 1938); \guillemotleft{}Robert II\guillemotright{} y \guillemotleft{}Robert III of
    Dreux\guillemotright{} ($\nu = 0{,}873$; 1154 y 1185).};
  \node[nota, right=3mm of entrada] {\emph{Fusión}: \texttt{can:robert a. stemmle} =
    \{Robert A. Stemmle, Robert Adolf Stemmle\}, pasajes \{biografía, película\}.
    \emph{Colisión}: \texttt{can:aleksander koniecpolski} (1620--1659) y
    \texttt{can:aleksander koniecpolski\#799} (1555--1609), abuelo y nieto.};
\end{tikzpicture}
\caption{Construcción de la memoria canónica de entidades sobre el corpus de
evaluación, con las cifras reales de cada etapa. El paso 2 no decide nada:
reduce $811$ millones de pares posibles a los $47\,063$ que se juzgan. El
paso 3 combina tres fuentes de unión ---automática, adjudicada por el modelo
de lenguaje y la heurística título$\to$sujeto--- que el cierre transitivo
consolida en grupos; el paso 4 convierte cada grupo en una entrada con
identidad. Las notas laterales dan un ejemplo real de cada etapa.}
\label{fig:diccionario}
\end{figure}

\paragraph{1. Ficha por (pasaje, forma de mención).}
Para cada pasaje $p$ y cada forma $f$ de su lista de menciones (la que
devolvió el extractor de la sección~\ref{sec:openie}) se compone una ficha
con los campos de la tabla~\ref{tab:ficha}. La ficha se construye
\emph{antes} de que exista nada canónico: lo canónico es la salida del
procedimiento, la ficha es su entrada, y por eso la misma persona nombrada
en dos pasajes produce dos fichas con descripciones distintas. Además, el
título de cada pasaje se registra como mención \emph{sintética} de su
sujeto cuando el extractor no lo devolvió ya como forma: un artículo
enciclopédico nombra a su sujeto en el título (\guillemotleft{}Charlie Day\guillemotright{}) aunque su
texto use la forma legal (\guillemotleft{}Charles Peckham Day\guillemotright{}), y el resto del corpus
suele referirse a él por la forma del título. Sobre el corpus de evaluación:
$38\,934$ menciones extraídas $+\,1\,342$ títulos sintéticos $= 40\,276$
fichas. Cada ficha recibe dos incrustaciones: la del nombre solo,
$\nu_m = e(f)$ (canal \emph{nombre}), y la de la ficha completa
$t_m = f \oplus d_m$, $\phi_m = e(t_m)$ (canal \emph{contexto}).

\begin{table}[t]
\centering\footnotesize\setlength{\tabcolsep}{4pt}
\caption{Campos de una ficha, con las dos fichas reales de Oskar Roehler.
La descripción $d_m$ concatena, sin repetir, las frases $\tau_a$ del pasaje
en las que la forma es sujeto, objeto o aparece en el texto (corte a 400
caracteres).}
\label{tab:ficha}
\begin{tabular}{l>{\raggedright\arraybackslash}p{0.19\textwidth}>{\raggedright\arraybackslash}p{0.26\textwidth}>{\raggedright\arraybackslash}p{0.26\textwidth}}
\toprule
campo & qué es & \texttt{Oskar Roehler||oskar roehler} & \texttt{Atomised (film)||oskar roehler}\\
\midrule
clave & pasaje $\|$ forma plegada & (encabezado) & (encabezado)\\
forma $f$ & la mención tal cual & Oskar Roehler & Oskar Roehler\\
pasaje & de dónde procede & \emph{Oskar Roehler} & \emph{Atomised (film)}\\
es\_título & ?`coincide con el título? & sí & no\\
descripción $d_m$ & frases $\tau_a$ del pasaje & \guillemotleft{}Oskar Roehler was born on 21 January
  1959. Oskar Roehler was born in Starnberg. Oskar Roehler is a film director.
  \dots\ is the son of Klaus Roehler. \dots\ is the son of Gisela Elsner. \dots\guillemotright{}
  & \guillemotleft{}Atomised was directed by Oskar Roehler. Atomised was written by Oskar
  Roehler.\guillemotright{}\\
nac.\ / muerte & años de las aserciones \texttt{DATE\_OF\_BIRTH} y \texttt{DATE\_OF\_DEATH} & \{1959\} / \{\} & \{\} / \{\}\\
texto de $\phi_m$ & \guillemotleft{}forma --- descripción\guillemotright{} & \guillemotleft{}Oskar Roehler --- Oskar Roehler was
  born \dots\guillemotright{} & \guillemotleft{}Oskar Roehler --- Atomised was directed \dots\guillemotright{}\\
texto de $\nu_m$ & la forma sola & \guillemotleft{}Oskar Roehler\guillemotright{} & \guillemotleft{}Oskar Roehler\guillemotright{}\\
\bottomrule
\end{tabular}
\end{table}

\emph{En el ejemplo}: los pasajes \emph{Oskar Roehler} y \emph{Atomised
(film)} generan $11 + 21 = 32$ fichas. Las dos de Oskar Roehler
(tabla~\ref{tab:ficha}) muestran lo distinta que es la descripción según el
pasaje: por eso el canal nombre es el principal y el contexto solo aporta
evidencia. En \emph{Atomised (film)} no hace falta título sintético (la
forma \guillemotleft{}Atomised\guillemotright{} ya estaba entre las menciones), de modo que su
ficha-título \texttt{Atomised (film)||atomised} lleva los nueve hechos de
la película. Y hay fichas pobres: \texttt{Atomised (film)||bruno}, un
personaje, tiene por toda descripción \guillemotleft{}Atomised stars Moritz Bleibtreu as
Bruno.\guillemotright{}

\paragraph{2. Generación de candidatos por dos canales.}
Con $40\,276$ fichas hay $811$ millones de pares posibles; ni las reglas
con restricciones, ni mucho menos el modelo de lenguaje, pueden juzgarlos
todos. El paso de candidatos es el filtro que reduce ese espacio a una
lista corta. Se calcula por bloques la fila \emph{completa} de cosenos de
cada ficha contra las $40\,276$ en ambos canales, y de cada fila se toman
las $15$ vecinas más próximas por canal (unión de ambas listas: a lo sumo
$30$ parejas por ficha). Se eliminan las parejas del mismo pasaje ---en un
artículo, formas distintas denotan entidades distintas--- y las que no
alcanzan ningún umbral ($\langle\nu_m,\nu_{m'}\rangle < 0{,}80$ y
$\langle\phi_m,\phi_{m'}\rangle < 0{,}75$). Lo que queda son los
candidatos que entran en el paso~3: $47\,063$ pares, de los que $32\,500$
caen en la banda automática y $9\,772$ en la banda del modelo de lenguaje.

\emph{En el ejemplo}: las $15$ vecinas de \texttt{Atomised (film)||oskar
roehler} por el canal nombre son, en orden, \texttt{Oskar Roehler||oskar
roehler} ($\nu = 1{,}000$, $\phi = 0{,}659$: fusión automática, correcta),
\texttt{Oskar Roehler||klaus roehler} ($0{,}675$ / $0{,}563$: descartada,
correcta ---es el padre---), y después Oscar Blumenthal, Oskar Drobne,
Oscar Straus\dots\ ($\nu \approx 0{,}50$--$0{,}59$), parecidos de nombre de
pila, todas descartadas. Por el canal contexto, en cambio, sus vecinas son
\texttt{Atomised (film)||atomised} ($\phi = 0{,}730$, mismo pasaje),
Christoph Schlingensief, Bernd Eichinger, Niklaus Schilling, Germany,
Berlin\dots: los \emph{compañeros de pasaje}. Lo mismo ocurre con
\texttt{H. P. Lovecraft||howard phillips lovecraft}, cuyas vecinas de
contexto son Rhode Island ($0{,}746$), Providence ($0{,}741$) y Cthulhu
Mythos. Es una propiedad del diseño que conviene dejar dicha: como la
descripción reúne frases del pasaje, \textbf{el canal contexto mide sobre
todo \guillemotleft{}mismo pasaje\guillemotright{}, no \guillemotleft{}misma entidad\guillemotright{}}, y esas parejas se excluyen por
regla; de media, $4{,}9$ de sus $15$ huecos los ocupan fichas del mismo
pasaje (frente a $0{,}6$ en el canal nombre). Su papel efectivo es de
confirmación (los umbrales de contexto de la ecuación~\ref{eq:adjudicacion}),
no de generación: el canal nombre es el que genera.

\emph{Sobre el tope de $15$ vecinas.} Como la fila completa de cosenos ya
se calcula, el tope no ahorra coste; solo limita cuántas parejas de una
fila ya pagada se examinan. Medido: $74\,580$ pares de pasajes distintos
superan algún umbral y el tope deja fuera $27\,517$ ($23\,260$ serían
fusión automática, $3\,902$ irían al modelo de lenguaje). Todos son
concentradores ---\guillemotleft{}France\guillemotright{} $4\,809$ pares, \guillemotleft{}United States\guillemotright{} $3\,711$,
\guillemotleft{}New York City\guillemotright{} $1\,955$--- y el cierre transitivo del paso~3 los
recupera: cada mención de \guillemotleft{}France\guillemotright{} ve a otras $15$ con coseno $1{,}000$ y
la estructura unión--búsqueda encadena las $111$ en una sola entrada. El
efecto práctico del tope es, por tanto, nulo en este corpus; las entidades
que importan para el multisalto (directores, madres, películas) tienen
pocas menciones, muy por debajo de $15$. El tope es un presupuesto, no un
criterio semántico, y como diseño sería preferible pre-agrupar por nombre
plegado exacto ($40\,276$ fichas $\to$ $30\,812$ formas distintas; $4\,530$
aparecen en dos o más pasajes), aplicar los umbrales a la fila completa y
limitar por ficha únicamente la banda del modelo de lenguaje, que es la que
tiene coste.

\paragraph{3. Adjudicación en tres bandas con restricciones duras.}
Antes de decidir se aplican dos restricciones de incompatibilidad: fechas
de nacimiento o muerte disjuntas, y numerales ordinales distintos en
títulos nobiliarios o regnales (\guillemotleft{}4th Baron\guillemotright{} frente a \guillemotleft{}5th Baron\guillemotright{};
\guillemotleft{}Edward III\guillemotright{} frente a \guillemotleft{}Edward IV\guillemotright{}), que separan parientes, sucesores y
homónimos aunque el coseno sea alto. Casos reales bloqueados: \guillemotleft{}Robert II
of Dreux\guillemotright{} (n.~1154) frente a \guillemotleft{}Robert III of Dreux\guillemotright{} (n.~1185), con
$\nu = 0{,}873$ y $\phi = 0{,}822$; \guillemotleft{}Robert Petre, 7th Baron Petre\guillemotright{} (n.~1689)
frente a \guillemotleft{}Robert James Petre, 8th Baron Petre\guillemotright{} (n.~1713), con
$0{,}914$ / $0{,}876$; y dos \guillemotleft{}William Ferguson\guillemotright{} con nombre idéntico
($\nu = 1{,}000$) nacidos en 1809 y 1938. Sobre los pares compatibles:
\begin{equation}
\op{decisión}(m,m') =
\begin{cases}
\text{fusión automática} & \text{si } \op{base}(f_m)=\op{base}(f_{m'}) \\[2pt]
\text{fusión automática} & \text{si } \langle\nu_m,\nu_{m'}\rangle \geq 0{,}90
                            \;\wedge\; \langle\phi_m,\phi_{m'}\rangle \geq 0{,}70 \\[2pt]
\text{adjudica el LLM} & \text{si } \langle\nu_m,\nu_{m'}\rangle \geq 0{,}80 \\[2pt]
\text{adjudica el LLM} & \text{si } \langle\phi_m,\phi_{m'}\rangle \geq 0{,}75
                            \;\wedge\; \text{término común} \\[2pt]
\text{descartar} & \text{en otro caso,}
\end{cases}
\label{eq:adjudicacion}
\end{equation}
donde $\op{base}(\cdot)$ normaliza el nombre (minúsculas, sin paréntesis).
La banda ambigua se resuelve con el modelo de lenguaje viendo las dos fichas
y sus artículos de origen, en lotes de $8$ pares con caché persistente, con
una consigna deliberadamente estricta: parientes, sucesores, homónimos y
nuevas versiones de una obra \emph{no} son la misma entidad; una forma
corta y una forma completa del mismo nombre \emph{sí}. El juez es
conservador: de los $9\,772$ pares que vio dijo \emph{misma} en $2\,322$
($24\%$). Un rechazo real: \guillemotleft{}A Matter of Time\guillemotright{} frente a \guillemotleft{}Just a Matter of
Time\guillemotright{}, con coseno $0{,}920$ en ambos canales, son dos películas.

A las uniones automáticas y a los \guillemotleft{}sí\guillemotright{} del juez se añade una tercera fuente
de unión, la heurística \emph{título$\to$sujeto}: en cada pasaje, la
mención-título se une a la mención no-título con más términos
significativos en común (desempatando por la de descripción más larga),
siempre que sus fechas sean compatibles. Está pensada para el título
sintético: en el artículo \emph{Charlie Day} une la forma del título con
\guillemotleft{}Charles Peckham Day\guillemotright{}; en \emph{Robert North Bradbury}, con \guillemotleft{}Robert N.
Bradbury\guillemotright{}. Produce $2\,848$ uniones, de las que $2\,431$ no estaban
cubiertas por los dos canales. Todas las uniones ---automáticas, del juez y
de título--- se cierran transitivamente mediante una estructura
unión--búsqueda: si $A \sim B$ y $B \sim C$, las tres fichas forman un
grupo aunque $A$ y $C$ nunca se compararan.

\paragraph{El caso Lovecraft, resuelto paso a paso.}
El ejemplo con que abre la sección permite ver los tres pasos actuando
sobre datos reales y, de paso, corregir una intuición fácil: la isla
\emph{no} la resuelve el juez. El corpus tiene cinco fichas con
\guillemotleft{}lovecraft\guillemotright{}: tres en el artículo \emph{H. P. Lovecraft}
(\texttt{howard phillips lovecraft}, con fechas 1890--1937; \texttt{h. p.
lovecraft}, que es el título; \texttt{lovecraftian horror}) y dos en
\emph{Sonia Greene} (\texttt{h. p. lovecraft}, con descripción \guillemotleft{}Sonia Haft
Greene Lovecraft Davis was married to H. P. Lovecraft.\guillemotright{}, y la propia Sonia).
Los pares entre pasajes distintos y su destino:

\begin{center}\footnotesize
\begin{tabular}{llccl}
\toprule
ficha $A$ & ficha $B$ & $\nu$ & $\phi$ & banda\\
\midrule
\emph{H. P. Lovecraft} $\|$ howard phillips lovecraft & \emph{Sonia Greene} $\|$ h. p. lovecraft & $\mathbf{0{,}781}$ & $0{,}479$ & descartado\\
\emph{H. P. Lovecraft} $\|$ h. p. lovecraft & \emph{Sonia Greene} $\|$ h. p. lovecraft & $1{,}000$ & $0{,}625$ & \textbf{automática}\\
\emph{H. P. Lovecraft} $\|$ lovecraftian horror & \emph{Sonia Greene} $\|$ h. p. lovecraft & $0{,}725$ & $0{,}439$ & descartado\\
(los seis restantes) & & $<0{,}50$ & $<0{,}45$ & descartado\\
\bottomrule
\end{tabular}
\end{center}

Ahí está el $0{,}781$: es el coseno de \emph{nombre} entre \guillemotleft{}Howard Phillips
Lovecraft\guillemotright{} y \guillemotleft{}H. P. Lovecraft\guillemotright{}. Con $\phi = 0{,}479 < 0{,}75$ y
$\nu < 0{,}80$, la ecuación~\eqref{eq:adjudicacion} lo descarta antes de
llegar al modelo de lenguaje (comprobado en la caché: el par nunca se
consultó). La isla la cierran los otros dos mecanismos: la heurística
título$\to$sujeto une, \emph{dentro} del artículo, la ficha-título
\guillemotleft{}H. P. Lovecraft\guillemotright{} con \guillemotleft{}Howard Phillips Lovecraft\guillemotright{} (comparten el término
\guillemotleft{}lovecraft\guillemotright{}, y la ficha-título no tiene fechas que lo impidan); y la
fusión automática une la ficha-título con la mención idéntica del artículo
de Sonia Greene ($\nu = 1{,}000$). El cierre transitivo hace el resto. La
lección es doble: los dos canales vectoriales solos \emph{no} habrían
resuelto el caso emblemático, y la resolución descansa en un supuesto
enciclopédico ---el título nombra al sujeto, y los demás artículos lo
nombran por esa forma--- cuyo alcance se discute más abajo.

\paragraph{4. Entrada canónica.}
Cada grupo resultante produce una entrada
\[
\varepsilon \;=\; \bigl(\text{nombre},\; \mathcal{A}_\varepsilon,\;
\text{descripción},\; \mathcal{P}_\varepsilon,\; \pi(\varepsilon)\bigr):
\]
nombre canónico (con preferencia por la forma-título), conjunto de alias
$\mathcal{A}_\varepsilon$ (todas las formas del grupo más los títulos de
sus pasajes propios), descripción fusionada (unión sin repetición de las
descripciones, hasta 600 caracteres), pasajes donde aparece y, si existe,
su \emph{pasaje propio} $\pi(\varepsilon)$ (el artículo cuyo sujeto es
$\varepsilon$). Se guarda además el mapa (pasaje, forma) $\to$
identificador, que es lo que el grafo usa para colapsar cada mención a su
nodo, y la ficha fusionada se incrusta como representación vectorial de la
entidad. El grupo de Lovecraft produce:

\begin{quote}\small
\texttt{id}: \texttt{can:h. p. lovecraft} \quad
\texttt{nombre}: H.\,P. Lovecraft \quad
\texttt{alias}: \{H.\,P. Lovecraft; Howard Phillips Lovecraft\}\\
\texttt{descripción}: \guillemotleft{}Howard Phillips Lovecraft was born on August 20,
1890 \dots\ was an American writer of weird fiction and horror fiction
\dots\ Sonia Haft Greene Lovecraft Davis was married to H P Lovecraft.\guillemotright{}\\
\texttt{pasajes}: \{\emph{H. P. Lovecraft}; \emph{Sonia Greene}\} \quad
\texttt{pasaje propio}: \emph{H. P. Lovecraft} \quad
\texttt{menciones}: 3
\end{quote}

\guillemotleft{}Lovecraftian horror\guillemotright{} queda, correctamente, como entrada aparte. Si dos
grupos distintos reciben el mismo nombre canónico, el segundo lleva un
sufijo: es una \emph{colisión}, y hay $36$ en el diccionario. Son homónimos
que las restricciones mantuvieron separados ---\texttt{can:aleksander
koniecpolski} (1620--1659) y \texttt{can:aleksander koniecpolski\#799}
(1555--1609), abuelo y nieto; \texttt{can:john middleton} (arquitecto,
n.~1820) y \texttt{can:john middleton\#1029} (futbolista, n.~1955)--- o dos
obras del mismo título citadas en un mismo artículo (\emph{The Time, the
Place and the Girl}, 1929 y 1946), que la regla del mismo pasaje no une.
Las colisiones son, por tanto, el diccionario \emph{negándose} a fundir:
no deben confundirse con su rendimiento, que resume la
tabla~\ref{tab:dicc-cifras}.

\begin{table}[t]
\centering\footnotesize\setlength{\tabcolsep}{5pt}
\caption{El diccionario sobre el corpus de evaluación: distribución de las
entradas por número de menciones fundidas y ejemplos reales de fusión con
formas de nombre distintas.}
\label{tab:dicc-cifras}
\begin{tabular}{lr@{\quad}l>{\raggedright\arraybackslash}p{0.34\textwidth}}
\toprule
menciones por entrada & entradas & ejemplo de fusión (id canónico) & alias\\
\midrule
1 & $21\,144$ & \texttt{can:robert a. stemmle} & Robert A. Stemmle; Robert Adolf Stemmle\\
2 & $3\,604$ & \texttt{can:ernest e. bryan} & Ernest Bryan; Ernest E. Bryan\\
3 & $1\,051$ & \texttt{can:gabriel cañellas} & Gabriel Cañellas; Gabriel Cañellas Fons\\
4 & $447$ & \texttt{can:love, honor, and oh baby!} & tres grafías del título\\
5 & $213$ & \texttt{can:united states} & U.S.; United States; United States of America\\
6--9 & $323$ & \texttt{can:h. p. lovecraft} & H. P. Lovecraft; Howard Phillips Lovecraft\\
$\ge 10$ & $178$ & \texttt{can:france} & France ($111$ menciones)\\
\midrule
\multicolumn{4}{l}{total: $26\,960$ entradas; $5\,816$ con más de una mención ($19\,132$ menciones fundidas);}\\
\multicolumn{4}{l}{$3\,450$ con dos o más formas de nombre distintas ($719$ con tres o más);}\\
\multicolumn{4}{l}{$31\,839$ alias; $5\,601$ con pasaje propio; $36$ colisiones.}\\
\bottomrule
\end{tabular}
\end{table}

El caso \texttt{can:robert a. stemmle} es el puente de libro: el artículo
de la película \emph{You Can No Longer Remain Silent} nombra al director
por su nombre legal completo y su biografía lleva las iniciales en el
título; antes del diccionario eran dos nodos sin arista fiable, ahora son
un nodo con dos alias, dos pasajes y la biografía como pasaje propio. En el
conjunto de sondeo, el procedimiento resolvió $178$ de las $188$
\guillemotleft{}islas\guillemotright{} detectadas en el análisis de errores.

\paragraph{Qué mecanismo fusiona qué.}
Como las tres fuentes de unión se aplican en cascada sobre la misma
estructura unión--búsqueda, su contribución se puede atribuir
reconstruyendo el diccionario con cada fuente añadida por turno
(tabla~\ref{tab:dicc-atribucion}; la reconstrucción reproduce el
diccionario publicado al entero: $26\,961$ entradas frente a $26\,960$).
La lectura es clara: la fusión automática hace el grueso (la misma cadena
en pasajes distintos), el juez aporta poco en número ($1\,164$ entradas
menos) y la heurística título$\to$sujeto \textbf{fusiona más que el
juez}: $2\,220$ entradas finales con varias menciones existen solo gracias
a ella. Es el mecanismo que resuelve Lovecraft y, como se verá, el que
concentra los errores conocidos.

\begin{table}[t]
\centering\footnotesize\setlength{\tabcolsep}{5pt}
\caption{Atribución de las fusiones del diccionario a cada fuente de unión,
reconstruyendo el cierre transitivo de forma acumulativa.}
\label{tab:dicc-atribucion}
\begin{tabular}{lrrr}
\toprule
fuentes de unión aplicadas & entradas & con $>1$ mención & fichas fundidas\\
\midrule
solo automática (nombre idéntico; $\nu\ge0{,}90 \wedge \phi\ge0{,}70$) & $30\,546$ & $4\,536$ & $14\,266$\\
$+$ juez (banda ambigua: $2\,322$ \guillemotleft{}sí\guillemotright{} de $9\,772$) & $29\,382$ & $4\,464$ & $15\,358$\\
$+$ título$\to$sujeto ($2\,848$ uniones) & $26\,961$ & $5\,815$ & $19\,130$\\
\bottomrule
\end{tabular}
\end{table}

\paragraph{Contribución medida al resultado.}
El efecto del diccionario sobre la recuperación está medido de forma
pareada sobre el banco de pruebas, comparando la versión con diccionario
(v3) con la inmediatamente anterior (v2), que resolvía la sinonimia en
consulta con aristas por umbral de coseno $0{,}80$ ---el esquema de
\hippo{}--- y era idéntica en todo lo demás (tabla~\ref{tab:dicc-impacto}).
Es el mayor salto de la secuencia de versiones de la
sección~\ref{sec:comparacion}. La reserva honesta es que mide el
\emph{paquete} que el diccionario hace posible ---el propio diccionario, el
enlazador por alias exacto (\S\ref{sec:linker}), los candidatos de frontera
y la masa en el pasaje propio (\S\ref{sec:siembra})--- y no el diccionario
aislado; pero el enlazador y el pasaje propio \emph{consumen} el diccionario
y no existen sin él.

\begin{table}[t]
\centering\footnotesize\setlength{\tabcolsep}{5pt}
\caption{Comparación pareada sobre las $1\,000$ preguntas del banco: versión
con memoria canónica (v3) frente a la misma versión con sinonimia por
coseno en consulta (v2). Diferencias en puntos con intervalo de confianza
del $95\%$ por bootstrap pareado.}
\label{tab:dicc-impacto}
\begin{tabular}{lcccc}
\toprule
métrica & v2 (sinonimia $0{,}80$) & v3 (diccionario) & $\Delta$ [IC $95\%$] & gana / pierde\\
\midrule
$\mathrm{R@}2$ & $70{,}4$ & $73{,}0$ & $+2{,}6$ $[+1{,}2, +3{,}9]$ sig. & $115$ / $67$\\
$\mathrm{R@}5$ & $90{,}9$ & $96{,}0$ & $+5{,}2$ $[+4{,}2, +6{,}2]$ sig. & $152$ / $12$\\
$\mathrm{FC@}5$ & $78{,}2$ & $90{,}6$ & $+12{,}4$ $[+10{,}2, +14{,}7]$ sig. & $136$ / $12$\\
$\mathrm{FC@}10$ & $85{,}8$ & $93{,}4$ & $+7{,}6$ $[+5{,}7, +9{,}5]$ sig. & $83$ / $7$\\
\midrule
$\mathrm{FC@}5$ comparación-puente & $55{,}7$ & $83{,}8$ & $+28{,}1$ $[+22{,}6, +34{,}0]$ sig. & \\
$\mathrm{FC@}5$ composicional & $79{,}2$ & $91{,}0$ & $+11{,}9$ $[+8{,}5, +15{,}5]$ sig. & \\
$\mathrm{FC@}5$ comparación & $96{,}3$ & $99{,}2$ & $+2{,}9$ $[+0{,}4, +5{,}3]$ sig. & \\
$\mathrm{FC@}5$ inferencia & & & $+1{,}9$ $[-4{,}6, +8{,}3]$ n.s. & \\
\bottomrule
\end{tabular}
\end{table}

\paragraph{Limitaciones conocidas y validación pendiente.}
La instrumentación anterior saca a la luz defectos que el agregado oculta
y que conviene dejar registrados.
\begin{enumerate}[leftmargin=1.6em, itemsep=1pt]
\item \textbf{Sobreaplicación de la heurística título$\to$sujeto.} Está
  pensada para el título sintético, pero se ejecuta también cuando el
  título ya es una mención propia del pasaje ($1\,549$ de las $2\,848$
  uniones); en ese caso excluye la ficha-título del bucle y une al sujeto
  con \emph{otra} mención que comparta apellido, franquicia o lugar, y las
  fechas rara vez lo impiden porque la mención secundaria no suele
  traerlas. El resultado es una mezcla de uniones legítimas (\guillemotleft{}Ryan
  Adams\guillemotright{} $\leftrightarrow$ \guillemotleft{}David Ryan Adams\guillemotright{}, el propio Lovecraft) y de
  exactamente lo que la consigna del juez prohíbe: parientes
  (\texttt{can:oskar roehler} lleva como alias a Klaus Roehler, su padre;
  \guillemotleft{}Ermengarde\guillemotright{} $\leftrightarrow$ \guillemotleft{}Hugh of Tours\guillemotright{}), secuelas (\emph{Wrong
  Turn 2/4/5} $\leftrightarrow$ \emph{Wrong Turn}) y personas con su empresa
  (\guillemotleft{}Pietro Salini\guillemotright{} $\leftrightarrow$ \guillemotleft{}Salini Costruttori SpA\guillemotright{}). Ninguno de esos
  pares pasó por el juez. Un segundo mecanismo de sobre-fusión, detectado
  en la traza de la sección~\ref{sec:siembra}, es la banda automática por
  nombre base idéntico cuando una de las formas lleva un desambiguador
  entre paréntesis: la mención sintética de título \guillemotleft{}Christopher
  Newton (criminal)\guillemotright{} no lleva fechas, su base sin paréntesis coincide
  con \guillemotleft{}Christopher Newton\guillemotright{} y la restricción de fechas (1969 frente a
  1936) no llega a actuar; título$\to$sujeto y el cierre transitivo hacen
  el resto, y el asesino y el director teatral comparten entrada y pasaje
  propio. El arreglo es que la ficha-título herede las fechas del sujeto
  de su artículo y que un nombre base idéntico con desambiguador vaya a la
  banda del juez. El arreglo previsto es no unir automáticamente en
  el caso no sintético y enviar el par a la banda del modelo de lenguaje
  (unos $200$ lotes), lo que exige regenerar diccionario, grafo y medición.
\item \textbf{El canal contexto mide pasaje, no entidad}, y su tope de
  vecinas se consume en parejas que luego se excluyen (véase el paso~2).
  Un rediseño debería construir la descripción solo con las frases en que
  la mención es sujeto u objeto y excluir el mismo pasaje \emph{antes} de
  tomar las vecinas.
\item \textbf{Infra-fusión residual.} Las variantes que ninguna semejanza
  de superficie ni de contexto breve une ---alias translingües (\guillemotleft{}I'll Be
  Going Now\guillemotright{} / \guillemotleft{}Tolgo il disturbo\guillemotright{}, sección~\ref{sec:falla-aun}),
  transliteraciones, acrónimos--- quedan como islas ($10$ de $188$ en el
  sondeo). Son el caso que una capa de identidad \emph{externa} resolvería.
\item \textbf{Calidad de las fusiones sin medir.} Se conoce el impacto
  del diccionario (tabla~\ref{tab:dicc-impacto}) pero no su precisión ni su
  exhaustividad como resolutor de entidades. El conjunto de datos ofrece
  un oro barato para medirlas: cada pregunta trae los identificadores
  Wikidata de sus entidades (\texttt{entity\_ids}) y de los sujetos y
  objetos de sus tripletes de evidencia (\texttt{evidences\_id}). Con ellos
  se puede comprobar, sobre los entes de evidencia del banco, si menciones
  del mismo QID caen en la misma entrada y las de QID distinto en entradas
  distintas, sin etiquetado manual. Esa auditoría debe preceder a cualquier
  decisión sobre cuánta curación adicional necesita el diccionario.
\item \textbf{Alcance y generalización.} Conviene separar tres capas. El
  \emph{patrón} ---ficha con contexto, dos canales, candidatos baratos,
  restricciones duras, adjudicación por lotes con caché y cierre
  transitivo--- es general y conservador, y es el mismo esqueleto del
  enlazador terminológico del caso clínico (sección~\ref{sec:gesida-index}).
  Lo que \emph{no} viaja es lo que da la mayor parte de la potencia en
  2Wiki: el supuesto de que el título nombra al sujeto y de que cada
  entidad tiene un artículo (de ahí el pasaje propio), y unas restricciones
  duras propias de biografías enciclopédicas (fechas vitales, ordinales
  nobiliarios). En un corpus sin esa estructura ---guías clínicas,
  prensa--- solo quedan los canales vectoriales y el juez, cuyos umbrales
  dependen además del codificador y del idioma ($0{,}80$ con
  \texttt{text-embedding-3-small} en inglés no vale para otro modelo ni
  para el español) y deben recalibrarse con una muestra etiquetada por
  corpus. Cuando existe una terminología externa, el paso cambia de
  naturaleza ---de \emph{agrupar menciones entre sí} a \emph{anclar
  menciones a un concepto}--- y el diccionario derivado del corpus queda
  como respaldo para lo que la terminología no cubre: en el asistente
  GeSIDA, el $70\%$ de las entidades sin identificador SNOMED~CT.
\end{enumerate}

\subsubsection{La memoria canónica como artefacto}
\label{sec:memoria}

Los cuatro pasos anteriores producen un objeto concreto que conviene mirar
por separado del procedimiento que lo construye, porque es lo que el resto
del sistema consume: la \emph{memoria canónica} es el artefacto que queda
entre la resolución de entidades y el grafo (figura~\ref{fig:pipeline}), y
en la práctica son cuatro almacenes en \texttt{artifacts/wiki2/}. La
figura~\ref{fig:memoria} los muestra con una entrada real
---\texttt{can:robert a. stemmle}, el caso puente de la
tabla~\ref{tab:dicc-cifras}--- y con quién usa cada pieza; la
tabla~\ref{tab:memoria} resume el esquema.

\begin{figure}[t]
\centering
\begin{tikzpicture}[
  font=\scriptsize,
  alm/.style={draw, rounded corners=2pt, align=left, inner sep=3pt, fill=gray!8, text width=96mm},
  cons/.style={draw, rounded corners=2pt, align=center, inner sep=3pt, text width=36mm},
  fl/.style={-{Stealth[length=2mm]}, semithick},
  node distance=3mm and 6mm]
  \node[alm] (ent) {\textbf{Entradas canónicas} --- \texttt{entidades\_<split>.json}, $26\,960$ entradas\\[1pt]
    {\tiny\color{black!80}
     \texttt{id}: \texttt{can:robert a. stemmle} \quad \texttt{nombre}: Robert A. Stemmle \quad \texttt{menciones}: 3\\
     \texttt{alias}: \{Robert A. Stemmle; Robert Adolf Stemmle\}\\
     \texttt{descripción}: \guillemotleft{}Robert Adolf Stemmle was born on 10 June 1903 in Magdeburg
     \dots\ was a German film director \dots\ directed 46 films between 1934 and 1970.
     You Can No Longer Remain Silent was directed by Robert A Stemmle.\guillemotright{}\\
     \texttt{pasajes}: \{Robert A. Stemmle; You Can No Longer Remain Silent\} \quad
     \texttt{pasaje\_propio}: Robert A. Stemmle}};
  \node[alm, below=of ent] (mapa) {\textbf{Mapa de menciones} --- \texttt{entidades\_mapa\_<split>.json},
    $40\,276$ claves (pasaje$\,\|\,$forma) $\to$ id\\[1pt]
    {\tiny\color{black!80}
     \texttt{Robert A. Stemmle||robert a. stemmle} $\to$ \texttt{can:robert a. stemmle} (ficha-título)\\
     \texttt{Robert A. Stemmle||robert adolf stemmle} $\to$ \texttt{can:robert a. stemmle} (título$\to$sujeto)\\
     \texttt{You Can No Longer Remain Silent||robert a. stemmle}\\ \hspace*{1em}$\to$ \texttt{can:robert a. stemmle} (nombre exacto)}};
  \node[alm, below=of mapa] (vec) {\textbf{Almacén vectorial} --- \texttt{emb\_fichas\_<split>.npz}, $26\,960 \times 1536$\\[1pt]
    {\tiny\color{black!80} fila 861 $=$ $e(\cdot)$ de \guillemotleft{}Robert A. Stemmle (Robert A. Stemmle, Robert Adolf Stemmle)
     --- Robert Adolf Stemmle was born on 10 June 1903 \dots\guillemotright{}}};
  \node[alm, below=of vec] (cache) {\textbf{Caché del juez} --- \texttt{entidades\_cache.sqlite},
    $12\,536$ pares (clave$_A$\,\#\#\,clave$_B$) $\to$ 0/1; solo interviene al reconstruir};
  % ---- consumidores
  \node[cons, right=of ent] (linker) {\textbf{Enlazador} (\S\ref{sec:linker})\\ escalón 1: tabla alias plegado $\to$ id;\\ escalón 2: vectores de ficha};
  \node[cons, below=of linker] (siembra) {\textbf{Siembra} (\S\ref{sec:siembra})\\ el pasaje propio $\pi(\varepsilon)$ recibe masa de reinicio (ec.~\ref{eq:pasaje-propio})};
  \node[cons, below=of siembra] (grafo) {\textbf{Grafo canónico} (\S\ref{sec:grafo})\\ un nodo por entrada; el mapa colapsa cada mención a su nodo};
  \node[cons, below=of grafo] (anclas) {\textbf{Anclas débiles} (\S\ref{sec:siembra})\\ dos entradas por coseno con cada vector de consulta};
  \draw[fl] (ent.east) -- (linker.west);
  \draw[fl] (ent.east) -- (siembra.west);
  \draw[fl] (mapa.east) -- (grafo.west);
  \draw[fl] (vec.east) -- (linker.west);
  \draw[fl] (vec.east) -- (anclas.west);
  \begin{scope}[on background layer]
    \node[fit=(ent)(cache), draw=gray!60, dashed, rounded corners=3pt, inner sep=4pt,
          label={[gray, font=\scriptsize, anchor=south west]north west:{memoria canónica de entidades (artefacto de indexación)}}] {};
  \end{scope}
\end{tikzpicture}
\caption{Anatomía de la memoria canónica construida sobre el corpus de
evaluación, con la entrada real \texttt{can:robert a. stemmle}: el
director aparece con las iniciales en el título de su biografía y con el
nombre completo en el texto, y solo con las iniciales en el artículo de la
película. La ficha-título y la forma completa se unen dentro de la
biografía (título$\to$sujeto); la mención de la película se une por nombre
exacto. A la derecha, quién consume cada almacén.}
\label{fig:memoria}
\end{figure}

\begin{table}[t]
\centering\footnotesize\setlength{\tabcolsep}{4pt}
\caption{Esquema de la memoria canónica: qué guarda cada campo y qué
componente lo consume.}
\label{tab:memoria}
\begin{tabular}{l>{\raggedright\arraybackslash}p{0.44\textwidth}>{\raggedright\arraybackslash}p{0.30\textwidth}}
\toprule
campo / almacén & qué es & quién lo consume\\
\midrule
\texttt{id} & identificador estable \texttt{can:}$+$nombre plegado (sufijo \texttt{\#n} en colisiones) & grafo (nombre del nodo), mapa\\
\texttt{nombre} & forma canónica; preferencia por la forma-título & enlazador (tabla de alias)\\
\texttt{alias} & todas las formas del grupo más los títulos de sus pasajes propios & enlazador, escalón 1 (coincidencia exacta tras plegado)\\
\texttt{descripción} & unión sin repetición de las descripciones de las fichas (hasta 600 car.) & solo a través del vector\\
\texttt{pasajes} & pasajes donde aparece alguna mención del grupo & auditoría; no interviene en consulta\\
\texttt{pasaje\_propio} & el artículo cuyo sujeto es la entidad, si existe & siembra (ec.~\ref{eq:pasaje-propio})\\
\texttt{menciones} & número de fichas fundidas & señal para la auditoría\\
mapa & (pasaje$\,\|\,$forma plegada) $\to$ \texttt{id}, una clave por ficha & constructor del grafo\\
vectores & $e(\text{nombre (alias) --- descripción})$ por entrada & enlazador, escalón 2; anclas débiles\\
caché del juez & veredicto 0/1 por par de claves & reconstrucción reproducible\\
\bottomrule
\end{tabular}
\end{table}

\paragraph{Cómo se consume.}
La memoria interviene en tres momentos. Al \emph{construir el grafo}, el
mapa colapsa las $40\,276$ menciones en $26\,960$ nodos-entidad: por eso
el grafo canónico no necesita aristas de sinonimia. En \emph{consulta}, el
enlazador busca cada tramo de la pregunta en la tabla de alias (una
consulta exacta de tabla, sin vectores) y, si falla, contra los vectores
de ficha; y la siembra deposita masa de reinicio en el pasaje propio de
cada entidad enlazada o aprobada por el filtro. En la traza de la
sección~\ref{sec:traza}, \guillemotleft{}Atomised\guillemotright{} resuelve por alias exacto a
\texttt{can:atomised} y el pasaje propio de \texttt{can:oskar roehler}
---descubierto por el filtro--- es lo que convierte la biografía del
director en semilla.

\paragraph{Ciclo de vida actual.}
Conviene decirlo con claridad: en la versión evaluada la memoria es
\emph{enteramente automática} y se regenera de cero en cada ejecución.
Tiene tres carencias que la instrumentación de la
sección~\ref{sec:diccionario} ha hecho visibles. No guarda
\emph{procedencia}: una vez que la estructura unión--búsqueda funde dos
fichas, nadie sabe si el alias entró por nombre exacto, por el juez o por
la heurística de título ---y así es como el alias \guillemotleft{}Klaus Roehler\guillemotright{} pudo
entrar en la entrada de su hijo sin que nada lo señalara---. No tiene
\emph{estado}: no distingue una entrada automática de una auditada o
curada. Y no tiene \emph{ciclo de vida}: una corrección manual no
sobreviviría a la siguiente reconstrucción. Una consecuencia práctica: el
enlazador construye la tabla de alias por orden de carga, de modo que en
las $36$ colisiones (y en cualquier alias compartido) gana la entrada que
se leyó primero. En el asistente clínico (sección~\ref{sec:gesida-index})
los anclajes terminológicos sí tienen bandeja de revisión con estados
(pendiente, curado, rechazado, sin anclar), registro de cada decisión y
escrituras atómicas; ese es el patrón que se reutilizaría.

\paragraph{Ciclo de curación previsto.}
La figura~\ref{fig:curacion} resume el diseño. La idea central es que
nadie va a revisar $27\,000$ entradas: la memoria tiene que \emph{señalar}
las que lo merecen, un juez las re-adjudica con más contexto, un humano
ve solo lo dudoso, y cada decisión se registra de forma que la siguiente
reconstrucción la respete. Lo que se cura, y cómo:
\begin{itemize}[leftmargin=1.4em, itemsep=1pt]
\item \emph{Alias}: cada alias pasa de cadena a objeto con origen (exacto,
  automático, juez, título, manual), evidencia (cosenos, veredicto) y
  estado (automático, confirmado, dudoso, rechazado). El enlazador solo usa
  los no rechazados.
\item \emph{Descripciones}: se regeneran de forma determinista a partir de
  las aserciones del pasaje propio (primero las de rol sujeto: nacimiento,
  muerte, ocupación, obras), guardando el identificador de cada frase para
  que sean auditables; admiten sobreescritura manual, que fuerza a
  recalcular el vector.
\item \emph{Pasaje propio y colisiones}: un grupo con dos fichas-título de
  pasajes distintos es sospechoso de haber fundido dos artículos; una
  colisión se marca como ambigüedad explícita para que el enlazador
  desambigüe (por coseno de la ficha con la pregunta) en lugar de elegir
  por orden de carga.
\item \emph{Identidad externa}: un campo para el identificador Wikidata
  (QID) aquí y SNOMED~CT (SCTID) en el caso clínico, que es a la vez oro
  de auditoría y la única vía para los alias que ninguna semejanza
  superficial une.
\item \emph{Decisiones como restricciones}: una curación se guarda como
  par \emph{forzado} o \emph{prohibido} entre claves de mención
  (pasaje$\,\|\,$forma), más las sobreescrituras de alias y descripción, y
  se aplica \emph{antes} del cierre transitivo en la siguiente
  reconstrucción. Así la regeneración es determinista (pasada automática
  $+$ decisiones), \guillemotleft{}separar $A$ de $B$\guillemotright{} impide que el cierre los vuelva a
  unir por otro camino, y las decisiones sobreviven a una re-extracción
  porque la clave es estable.
\end{itemize}

\begin{figure}[t]
\centering
\begin{tikzpicture}[
  font=\scriptsize,
  hecho/.style={draw, rounded corners=2pt, align=center, inner sep=3pt, fill=gray!8, text width=34mm},
  prev/.style={draw, densely dotted, rounded corners=2pt, align=center, inner sep=3pt, text width=34mm},
  dato/.style={draw, dashed, rounded corners=2pt, align=center, inner sep=3pt, text width=34mm},
  fl/.style={-{Stealth[length=2mm]}, semithick},
  et/.style={font=\tiny, fill=white, inner sep=1pt},
  node distance=6mm and 5mm]
  \node[hecho] (mem) {\textbf{Memoria canónica $v_n$}\\ entradas, alias, descripciones, mapa, vectores\\ (hoy: pasada automática de \S\ref{sec:diccionario})};
  \node[prev, right=of mem] (sen) {\textbf{Señales}\\ oro externo (QID); unión de título no sintética; apellido igual con nombre distinto; sufijo de empresa; secuelas; alias que es título de otro pasaje; colisiones};
  \node[prev, right=of sen] (juez) {\textbf{Re-adjudicación}\\ juez LLM con la consigna estricta y las primeras frases de ambos pasajes; lotes de 8 con caché};
  \node[prev, right=of juez] (band) {\textbf{Bandeja humana} (opcional)\\ solo dudosos y conflictos con el oro; estados y registro como en el caso clínico};
  \node[prev, below=of band] (dec) {\textbf{Decisiones registradas}\\ pares forzados / prohibidos entre claves (pasaje$\,\|\,$forma); alias y descripciones sobreescritos; identidad externa};
  \node[prev, below=of juez] (regen) {\textbf{Regeneración determinista}\\ pasada automática $+$ decisiones \emph{antes} del cierre; se rehacen solo descripciones, vectores y nodos de las entradas tocadas};
  \node[dato, below=of mem] (cons) {\textbf{Consumidores}\\ grafo canónico (\S\ref{sec:grafo}); enlazador (\S\ref{sec:linker}); siembra por pasaje propio (\S\ref{sec:siembra})};
  \draw[fl] (mem) -- (sen);
  \draw[fl] (sen) -- (juez);
  \draw[fl] (juez) -- (band);
  \draw[fl] (band) -- (dec);
  \draw[fl] (dec) -- (regen);
  \draw[fl] (regen.west) -| node[et, pos=0.75, left]{$v_{n+1}$} (sen.south);
  \draw[fl] (mem) -- (cons);
\end{tikzpicture}
\caption{Ciclo de curación previsto para la memoria canónica. Las cajas
rellenas existen en la versión evaluada; las punteadas son diseño. La
memoria señala las entradas sospechosas, el juez las re-adjudica, un
revisor ve solo lo dudoso, y las decisiones se guardan como restricciones
sobre claves de mención que la siguiente reconstrucción aplica antes del
cierre transitivo, de modo que la memoria $v_{n+1}$ es reproducible y las
correcciones no se pierden.}
\label{fig:curacion}
\end{figure}

\paragraph{Pipeline de evoluciones de la memoria canónica.}
Reunimos aquí, ordenadas por valor esperado frente a esfuerzo, todas las
evoluciones que el análisis de este apartado y del anterior deja
planteadas. Cada una se validaría primero en el sondeo y se mediría una
sola vez sobre el banco, como el resto del sistema
(sección~\ref{sec:protocolo}); ninguna está implementada en la versión
evaluada.
\begin{enumerate}[leftmargin=1.8em, itemsep=1pt, label=E\arabic*.]
\item \textbf{Auditoría con oro externo.} Medir precisión y exhaustividad
  de las fusiones con los identificadores Wikidata del propio conjunto de
  datos (\texttt{entity\_ids}, \texttt{evidences\_id}): mismo QID en
  entradas distintas es infra-fusión; QID distintos en la misma entrada,
  sobre-fusión. Sin coste de API. Es la primera porque decide el alcance de
  todas las demás (limitación~4 de la sección~\ref{sec:diccionario}).
\item \textbf{Procedencia y señales en el esquema.} Registrar en la
  consolidación el origen de cada unión (la información existe en el
  proceso; hoy se pierde al fundir), convertir los alias en objetos con
  origen, evidencia y estado, y calcular las señales estructurales de la
  figura~\ref{fig:curacion}. No cambia ningún resultado; hace la memoria
  auditable.
\item \textbf{Decisiones registradas y regeneración determinista.} Pares
  forzados y prohibidos sobre claves de mención, sobreescrituras de alias y
  descripción, aplicadas antes del cierre transitivo; reconstrucción
  incremental de las entradas tocadas.
\item \textbf{Re-adjudicación de las uniones señaladas y re-medición.}
  Enviar al juez las $1\,549$ uniones de título fuera de su caso previsto
  y el resto de señales (unos $250$ lotes, céntimos), regenerar
  diccionario, grafo y vectores, y medir sobre el banco (${\approx}$1--2
  USD; cuenta como una configuración más). Es el arreglo de la
  limitación~1 hecho como curación registrada, no como otro parche
  heurístico.
\item \textbf{Candidatos por umbral, no por tope.} Pre-agrupar por nombre
  plegado exacto con tabla de dispersión ($40\,276$ fichas $\to$ $30\,812$
  formas; $4\,530$ en dos o más pasajes) manteniendo las restricciones
  duras; aplicar los umbrales a la fila completa de cosenos excluyendo
  antes el mismo pasaje; limitar por ficha solo la banda del juez, que es
  la que tiene coste.
\item \textbf{Canal contexto que mida entidad y no pasaje.} Construir la
  descripción solo con las frases en que la mención es sujeto u objeto y
  excluir las fichas del mismo pasaje antes de tomar las vecinas
  (limitación~2).
\item \textbf{Enlazador consciente de colisiones.} La memoria expone la
  ambigüedad (varias entradas por alias) y el enlazador desambigua por
  coseno de la ficha con la pregunta, o siembra ambas con peso repartido,
  en lugar de elegir por orden de carga.
\item \textbf{Capa de identidad externa.} Resolución título $\to$ QID
  contra Wikidata, cacheada, como identidad primaria de la entrada: es la
  única vía para los alias translingües y las transliteraciones
  (limitación~3), y el equivalente exacto del SCTID clínico.
\item \textbf{Bandeja de curación unificada.} Reutilizar la bandeja del
  asistente clínico (estados, registro de decisiones, escrituras atómicas)
  para ambas memorias, solo donde vaya a haber un revisor; y, en sentido
  inverso, llevar al caso clínico el diccionario derivado del corpus como
  respaldo para el $70\%$ de entidades sin SCTID (limitación~5).
\end{enumerate}

\subsubsection{Grafo canónico y grafo de búsqueda}
\label{sec:grafo}

El grafo de conocimiento es $G=(V,E)$ con tres clases de nodo,
$V = \mathcal{C} \cup \mathcal{E} \cup \mathcal{A}$: pasajes
($|\mathcal{C}|=6\,119$), entidades canónicas ($|\mathcal{E}|=26\,940$) y
aserciones ($|\mathcal{A}|=47\,999$). Cada mención de sujeto, objeto o
participante se resuelve a su entrada canónica mediante el mapa
(pasaje, forma) $\to$ id construido en la sección~\ref{sec:diccionario}, de
modo que todas las variantes de nombre colapsan en un solo nodo y el grafo
canónico \emph{no necesita aristas de sinonimia}. Las aristas tipadas y sus
recuentos reales se dan en la tabla~\ref{tab:grafo}. La reificación es
$n$-aria: además del sujeto y el objeto, todo participante del pasaje cuya
forma aparezca en la frase $\tau_a$ recibe una arista de rol; sin ello, en
hechos con tres o más participantes la entidad adicional queda inalcanzable
desde el hecho. Los literales (fechas, oficios) viajan dentro de la
aserción y no generan nodos.

\begin{table}[t]
\centering\small
\caption{Esquema y tamaño del grafo canónico sobre el corpus de evaluación
($81\,058$ nodos, $165\,240$ aristas). La columna \guillemotleft{}orientación en búsqueda\guillemotright{}
indica el sentido de la arista en el grafo dirigido $G^{\!\rightarrow}$ sobre
el que corre el PPR.}
\label{tab:grafo}
\begin{tabular}{llrl}
\toprule
arista & semántica & n.\textsuperscript{o} & orientación en búsqueda\\
\midrule
aserción $\to$ pasaje   & procedencia del hecho          & $47\,999$ & aserción $\to$ pasaje\\
aserción $\to$ entidad  & rol de sujeto                  & $47\,466$ & entidad $\to$ aserción (invertida)\\
aserción $\to$ entidad  & rol de objeto / participante   & $32\,964$ & entidad $\to$ aserción (invertida)\\
entidad $\to$ pasaje    & la entidad se menciona en      & $36\,811$ & entidad $\to$ pasaje\\
\bottomrule
\end{tabular}
\end{table}

Para la consulta se deriva un grafo dirigido de búsqueda
$G^{\!\rightarrow}$ invirtiendo las aristas de rol (entidad $\to$
aserción) y conservando el resto de orientaciones. El resultado, verificado
sobre el grafo real, es un grafo \emph{por capas} estricto ---un grafo
acíclico dirigido de profundidad dos---:
\[
\underbrace{\text{entidad}}_{\text{grado de entrada }0}
\;\longrightarrow\;
\underbrace{\text{aserción}}_{\text{grado de salida }1}
\;\longrightarrow\;
\underbrace{\text{pasaje}}_{\text{sumidero}},
\qquad\qquad
\text{entidad} \;\longrightarrow\; \text{pasaje},
\]
con $80\,430$ aristas entidad$\to$aserción, $47\,999$ aserción$\to$pasaje
(cada aserción desagua exactamente en su pasaje de procedencia) y
$36\,811$ entidad$\to$pasaje (menciones); los $6\,119$ nodos-pasaje no
tienen salida. Esta estructura tiene consecuencias exactas sobre el paseo
---incluida una forma cerrada del punto fijo--- que se desarrollan en la
sección~\ref{sec:ppr}. Todas las aristas tienen peso base $1$; el peso
efectivo de las que entran en nodos-aserción se modula por consulta
(ecuación~\ref{eq:modulacion}).

\subsubsection{Representación dual: el grafo simbólico y sus almacenes
vectoriales}
\label{sec:dual}

Conviene despejar una posible confusión: el grafo canónico \emph{no}
contiene vectores, pero el índice completo sí. La representación es dual.
El grafo aporta la parte \emph{simbólica} ---nodos con identidad,
aristas tipadas, procedencia--- y, alineados con él por identificador de
nodo, tres almacenes vectoriales aportan la parte \emph{semántica}
(tabla~\ref{tab:embeddings}): todo objeto direccionable del índice
(aserción, entidad canónica, pasaje) tiene exactamente una incrustación.
Durante la construcción del diccionario se calculan además dos
incrustaciones \emph{transitorias} por mención (los canales de nombre y de
ficha, sección~\ref{sec:diccionario}), que no forman parte del índice
final.

\begin{table}[t]
\centering\footnotesize\setlength{\tabcolsep}{3.5pt}
\caption{Almacenes vectoriales que acompañan al grafo (modelo
\texttt{text-embedding-3-small}, $d=1536$, vectores normalizados). El grafo
GraphML es puramente simbólico; la alineación es por identificador de nodo.}
\label{tab:embeddings}
\begin{tabular}{llp{0.40\textwidth}}
\toprule
objeto & texto incrustado & uso en consulta\\
\midrule
aserción $a$ ($47\,999$) & la frase $\tau_a$ & afinidad $\sigma_a$
  (ec.~\ref{eq:afinidad}); candidatos del filtro; modulación $\mu$
  (ec.~\ref{eq:modulacion})\\
entidad canónica $\varepsilon$ ($26\,960$) & ficha: nombre (alias) ---
  descripción & respaldo del enlazador ($\cos\ge 0{,}80$); anclas
  simbólicas débiles\\
pasaje $c$ ($6\,119$) & título.\ texto & presa densa de la siembra
  (fuente~b de \S\ref{sec:siembra})\\
\midrule
mención $m$ (transitoria) & nombre $f_m$ \,/\, ficha $t_m$ & solo en
  indexación: candidatos a fusión por dos canales\\
\bottomrule
\end{tabular}
\end{table}

\paragraph{Dónde y cómo se resuelve la sinonimia.}
La recuperación de entidades en consulta la hace el enlazador
(sección~\ref{sec:linker}) en dos escalones: coincidencia \emph{exacta} de
alias tras plegado ---posible precisamente porque el diccionario ya agrupó
las variantes--- y, si falla, vecindad vectorial contra las fichas. No hay
distancia de edición (SymSpell, Levenshtein) en ningún punto del sistema
evaluado: las variantes de superficie que una distancia de edición
capturaría (\guillemotleft{}H.\,P.\guillemotright{} frente a \guillemotleft{}Howard Phillips\guillemotright{}) quedan resueltas
\emph{en indexación} por la fusión canónica, y las que no lo están
(erratas del usuario, formas nunca vistas) deberían absorberlas el
escalón vectorial ---aunque, medido sobre el banco, ese escalón apenas se
activa en la versión evaluada (sección~\ref{sec:linker})---. Merece la pena distinguir los tres regímenes de sinonimia que
conviven en este trabajo: (i)~la variante \emph{no canónica} del grafo
(usada en la versión v2 y conservada como ablación) resuelve en
\emph{consulta} con aristas \textsc{sinonimo\_de} por umbral de coseno
($\ge 0{,}80$), el esquema de \hippo{}; (ii)~la variante canónica (v3)
resuelve en \emph{indexación} con el diccionario ---candidatos
vectoriales, restricciones duras y adjudicación LLM--- y elimina esas
aristas; (iii)~el prototipo clínico (sección~\ref{sec:gesida}) resuelve
contra una terminología \emph{externa} (SNOMED~CT vía UMLS) con un
enlazador en cascada que sí incluye un escalón difuso léxico (FTS5 de
trigramas con coeficiente de Dice), el papel que en otros sistemas
desempeña SymSpell.

% =========================================================================
\subsection{Fase de consulta}
\label{sec:consulta}

Las seis etapas de consulta también encadenan productos explícitos, y
conviene fijar el contrato antes de entrar en cada una: el
\emph{enrutador} produce la estrategia; el \emph{enlazador} produce las
entidades enlazadas $L(q)$ (consultando el diccionario); la
\emph{afinidad} produce la puntuación $\sigma_a$ de cada hecho
(consultando el almacén de aserciones); el \emph{filtro} produce el par
$(S(q),\,\tilde\sigma)$ ---qué hechos son útiles y con qué peso---; la
\emph{siembra} combina $S(q)$, $L(q)$ y la señal densa en el vector
$\mathbf{v}$; y el \emph{paseo} convierte $(\mathbf{v},\,\tilde\sigma)$ en
el orden final sobre los pasajes. Para que el encadenamiento se vea en
acción, seguiremos un \emph{ejemplo conductor} a lo largo de las seis
etapas: una pregunta real de comparación-puente del banco de pruebas,
ejecutada instrumentada tal como se evaluó,
\begin{quote}
\emph{\guillemotleft{}Which film has the director who is older, God's Gift to Women or
Aldri Annet Enn Bråk?\guillemotright{}}
\end{quote}
cuya evidencia oro son \emph{cuatro} pasajes: las dos películas y las
biografías de sus directores (Michael Curtiz y Edith Carlmar). Es el tipo
de pregunta más difícil del banco ---la familia en la que los sistemas
publicados no superan el $12{,}3\%$ de cadena completa
(sección~\ref{sec:fallan-otros})--- y por eso el mejor banco de pruebas
del encadenamiento entre etapas. Al final de cada etapa, una nota
\guillemotleft{}en el ejemplo\guillemotright{} registra qué produjo y qué entrega a la siguiente;
la sección~\ref{sec:traza} recorre después, con el mismo detalle, un caso
composicional.

\paragraph{Qué depende de qué.}
Las etapas no forman una cadena lineal sino un grafo de dependencias con
dos ramas independientes (figura~\ref{fig:pipeline} y
tabla~\ref{tab:dependencias}). El enrutador solo lee la pregunta y decide
la estrategia. A partir de ahí, la rama \emph{simbólica} ---el enlazador,
que consulta el diccionario--- y la rama \emph{semántica} ---los vectores
de consulta y la afinidad, que consultan el almacén de aserciones---
consumen únicamente la pregunta y son independientes entre sí: ninguna
usa la salida de la otra, y pueden ejecutarse en cualquier orden o en
paralelo. Convergen en el filtro, la primera etapa que necesita las dos:
su lista de candidatos $\mathcal{K}(q)$ se compone del Top-40 por
$\sigma$ (rama semántica) y de la frontera de $L(q)$ (rama simbólica).
La siembra es la única etapa que recibe los tres productos ---$L(q)$;
$\{v_j\}$ con $\sigma$; y $(S(q),\tilde\sigma)$--- y el paseo recibe la
siembra y $\tilde\sigma$. En la estrategia de descomposición el filtro no
corre, $\tilde\sigma=\sigma$, y la siembra recibe directamente los
productos de las dos ramas.

\begin{table}[t]
\centering\footnotesize\setlength{\tabcolsep}{3.5pt}
\caption{Dependencias entre las etapas de consulta: qué consume y qué
produce cada una, y quién consume su producto. Enlazador y afinidad son
independientes; el filtro es la primera etapa que necesita ambas.}
\label{tab:dependencias}
\begin{tabular}{lp{0.30\textwidth}p{0.17\textwidth}p{0.33\textwidth}}
\toprule
etapa & consume & produce & quién lo consume\\
\midrule
enrutador & la pregunta & estrategia & afinidad (?`descomponer?), filtro (?`corre?)\\
enlazador & la pregunta, la memoria canónica (alias, fichas) & $L(q)$ & filtro (frontera); siembra (anclas fuertes, pasaje propio)\\
afinidad & la pregunta, el almacén de aserciones & $\{v_j\}$, $\sigma$ & filtro (Top-40); siembra (compuerta, presa densa con $v_0$, anclas débiles); paseo ($\mu$, si no hay filtro)\\
filtro & $\mathcal{K}(q)=\op{Top}_{40}(\sigma)\cup\op{frontera}(L(q))$ & $S(q)$, $\tilde\sigma$ & siembra (compuerta dura); paseo ($\mu$)\\
siembra & $L(q)$; $\{v_j\}$; $(S(q),\tilde\sigma)$; pasaje propio de la memoria & $\mathbf{v}$ & paseo\\
paseo & $\mathbf{v}$, $\tilde\sigma$, el grafo $G^{\!\rightarrow}$ & orden de pasajes & ---\\
\bottomrule
\end{tabular}
\end{table}

\subsubsection{Enrutador léxico de intención}
\label{sec:router}

Las preguntas de 2WikiMultihopQA pertenecen a cuatro tipos con estructuras
de evidencia distintas (composicional, comparación, inferencia y
comparación-puente). Un clasificador de reglas léxicas de coste cero las
distingue evaluando tres reglas \emph{en cascada}, de la más específica a
la más general, con la composicional como clase por defecto:
\begin{enumerate}[leftmargin=1.6em, itemsep=1pt]
\item \textbf{Comparación-puente}: disyunción (\emph{or}/\emph{both})
  $+$ un papel creativo (director, intérprete, autor\dots) $+$ un tipo de
  obra (\emph{film}, \emph{novel}\dots). La combinación delata que hay que
  resolver el papel \emph{en cada rama} antes de comparar.
\item \textbf{Comparación}: marcas comparativas directas
  (\emph{are both}, \emph{which \dots first/older/longer},
  \emph{same nationality/country/year}) o la disyunción sola: las dos
  entidades a compilar ya están nombradas en la pregunta.
\item \textbf{Inferencia}: vocabulario de parentesco de segundo grado
  (\emph{grandfather}, \emph{in-law}, \emph{maternal}\dots), que exige
  componer dos relaciones familiares.
\item \textbf{Composicional} (por defecto): ninguna marca anterior.
\end{enumerate}
La tabla~\ref{tab:router-ejemplos} muestra una pregunta real del banco por
tipo con la marca que dispara. Sobre el conjunto de calibración
($11\,376$ preguntas; nunca el banco de pruebas) el clasificador acierta
el $96{,}7\%$. El tipo decide la estrategia: \emph{comparación} usa la
variante de descomposición (sección~\ref{sec:afinidad}), donde las
subconsultas por entidad rinden más que el filtro; los otros tres tipos
usan el filtro selector (sección~\ref{sec:filtro}). Por eso la confusión
residual entre composicional e inferencia es benigna: ambas rutan a la
misma estrategia.

\begin{table}[t]
\centering\footnotesize\setlength{\tabcolsep}{3.5pt}
\caption{Una pregunta real del banco de pruebas por tipo, con la regla
léxica que la clasifica y la estrategia resultante. Las reglas se evalúan
en cascada: en el ejemplo conductor disparan las reglas 1 y 2 a la vez y
gana la más específica.}
\label{tab:router-ejemplos}
\begin{tabular}{p{0.16\textwidth}p{0.42\textwidth}p{0.20\textwidth}l}
\toprule
tipo & pregunta real & marca que dispara & estrategia\\
\midrule
comparación-puente & \emph{Which film has the director who is older,
  God's Gift to Women or Aldri Annet Enn Bråk?} & \emph{or} $+$
  \emph{director} $+$ \emph{film} (regla 1) & filtro\\
comparación & \emph{Are Christopher Newton (Criminal) and Frances M.
  Vega of the same nationality?} & \emph{same nationality} (regla 2) &
  descomposición\\
inferencia & \emph{Who is Raghnall Mac Ruaidhrí's paternal grandfather?}
  & \emph{paternal grandfather} (regla 3) & filtro\\
composicional & \emph{Who is the mother of the director of film Atomised
  (Film)?} & ninguna (por defecto) & filtro\\
\bottomrule
\end{tabular}
\end{table}

\emph{En el ejemplo conductor}: \guillemotleft{}or\guillemotright{} $+$ \guillemotleft{}director\guillemotright{} $+$ \guillemotleft{}film\guillemotright{}
disparan la regla 1 $\Rightarrow$ comparación-puente $\Rightarrow$
estrategia de filtro selector. Lo que se entrega a las etapas siguientes
es solo esa decisión; el texto de la pregunta sigue intacto.

\paragraph{Alcance y generalización.}
Estas reglas son deliberadamente específicas del banco ---la taxonomía de
cuatro tipos es la del propio conjunto de datos, y las marcas léxicas
explotan que sus preguntas se generan por plantillas--- y no
generalizarían a otro idioma ni a preguntas de formulación libre. Tres
hechos acotan el riesgo metodológico de esta elección. Primero, el
enrutador decide poco: tres de los cuatro tipos rutan a la misma
estrategia (el filtro, que es la ruta general), de modo que un error de
clasificación solo cuesta la especialización de comparación. Segundo, su
contribución al error está medida: clasifica mal $29$ de las $1\,000$
preguntas del banco y estas fallan a la tasa base ($5$ de los $94$ fallos
de cadena de la versión final estaban mal enrutados, $\approx$ tasa base),
así que no es causa medible de fallo (sección~\ref{sec:falla-aun}).
Tercero, su ganancia también: frente a la estrategia única de filtro
aportó $+2{,}0$ / $+3{,}7$ puntos ($\mathrm{R@}5$ / $\mathrm{FC@}5$),
concentrados en comparación. Es, en suma, una optimización de coste cero,
no una pieza estructural; en un despliegue abierto se sustituye por un
enrutador con modelo de lenguaje ---exactamente lo que hace el prototipo
clínico (sección~\ref{sec:gesida-consulta}), con respaldo heurístico--- al
precio de una llamada más, sin tocar el resto de la arquitectura. La misma
consideración vale para la otra heurística léxica del sistema, los tramos
capitalizados del enlazador (sección~\ref{sec:linker}), cuyo sustituto
general (reconocimiento de entidades por modelo de lenguaje) figura en las
líneas de mejora de la sección~\ref{sec:falla-aun}.

\subsubsection{Enlazador de menciones}
\label{sec:linker}

El enlazador identifica las entidades nombradas en la pregunta y las
resuelve contra el diccionario. Primero se extraen los \emph{tramos}
candidatos: secuencias maximales de palabras capitalizadas (admitiendo
conectores en minúscula), excluido el vocabulario interrogativo; los
tramos coordinados (\guillemotleft{}$A$ and $B$\guillemotright{}) se dividen si el conjunto no es un
alias conocido. Cada tramo se resuelve después en una cascada de dos
escalones de naturaleza distinta:

\paragraph{Escalón 1: resolución simbólica (alias exacto).}
El tramo se pliega (minúsculas, sin paréntesis ni puntuación) y se busca en
una tabla forma plegada $\to$ identificador canónico construida con todos
los alias del diccionario. Es una consulta exacta de tabla ---sin vectores,
coste constante--- y es la vía \emph{principal}: resulta posible
precisamente porque la fusión canónica de la
sección~\ref{sec:diccionario} ya agrupó las variantes de nombre en
indexación, de modo que \guillemotleft{}Atomised (film)\guillemotright{} y \guillemotleft{}Atomised\guillemotright{} apuntan a la
misma entrada. Si el tramo lleva paréntesis y no casa, se reintenta sin
ellos. En este banco, donde las preguntas nombran los títulos casi
literalmente, este escalón resuelve la gran mayoría de los tramos.

\paragraph{Escalón 2: respaldo vectorial sobre las fichas canónicas.}
Solo los tramos que no casan con ningún alias pasan al segundo escalón: se
incrusta el tramo (una llamada al codificador) y se compara por coseno
contra las incrustaciones de las \emph{fichas canónicas}
(tabla~\ref{tab:embeddings}: nombre, alias y descripción de cada una de
las $26\,960$ entradas); se toma la más afín y se acepta solo si
$\cos \geq 0{,}80$. Si ningún candidato alcanza el umbral, el tramo queda
sin enlazar ---un fallo silencioso tolerable, porque el filtro (vía el
Top-40 por coseno) y la presa densa siguen operando sin él---. Obsérvese
que ambos escalones devuelven siempre una entrada \emph{canónica}: lo que
cambia entre ellos es el mecanismo de búsqueda (tabla exacta frente a
vecindad vectorial), no el espacio sobre el que se busca.

\paragraph{Medición del escalón 2 sobre el banco.}
Instrumentando el enlazador sobre las $1\,000$ preguntas, el recorte
produce $1\,534$ tramos: $1\,259$ ($82\%$) resuelven por alias exacto y
$275$ ($18\%$) llegan al escalón 2. De estos, el umbral $0{,}80$ acepta
\textbf{tres} (coseno máximo observado $0{,}850$; mediana $0{,}61$). En
la práctica, por tanto, el respaldo vectorial está \emph{inactivo} en la
versión evaluada, y las preguntas afectadas las sostienen el Top-40 del
filtro y la presa densa. La señal, sin embargo, apunta bien: la ficha más
afín pertenece a un pasaje oro en $216$ de los $275$ casos ($79\%$). Las
causas son dos y distintas. La primera es el \emph{recorte}: el artículo
inicial se pierde porque \guillemotleft{}the\guillemotright{}, \guillemotleft{}a\guillemotright{} y \guillemotleft{}an\guillemotright{} figuran en el
vocabulario interrogativo que la regla de mayúsculas excluye
(\guillemotleft{}The Goose Woman\guillemotright{} $\to$ \guillemotleft{}Goose Woman\guillemotright{}; $120$ de los $275$ tramos
casarían por alias exacto con su artículo), y los dos puntos, la \guillemotleft{}A\guillemotright{}
interior y la coordinación parten títulos (\guillemotleft{}Gaby: A True Story\guillemotright{}
$\to$ \guillemotleft{}Gaby\guillemotright{} $+$ \guillemotleft{}True Story\guillemotright{}; \guillemotleft{}Interview With A Hitman\guillemotright{}
$\to$ \guillemotleft{}Interview With\guillemotright{} $+$ \guillemotleft{}Hitman\guillemotright{}). La segunda es el
\emph{objeto de comparación}: un nombre suelto contra una ficha que lleva
$300$ caracteres de descripción rara vez alcanza $0{,}80$. El tramo
\guillemotleft{}Goose Woman\guillemotright{} contra la ficha \guillemotleft{}The Goose Woman (The Goose Woman)
--- The Goose Woman is a 1925 silent film drama directed by Clarence
Brown\dots\guillemotright{} da $0{,}657$ (rechazado, siendo la ficha correcta); el mismo
tramo contra la incrustación del \emph{nombre solo} \guillemotleft{}The Goose
Woman\guillemotright{} da $0{,}933$.

\paragraph{Qué es denso aquí y qué no.}
Conviene no confundir este respaldo vectorial con la recuperación densa
clásica. El enlazador ancla \emph{entidades}: su espacio son las $26\,960$
fichas del diccionario y su entrada es el tramo aislado. La \emph{presa
densa} de la siembra (fuente~b de la sección~\ref{sec:siembra}) es otro
mecanismo, posterior e independiente: puntúa \emph{pasajes} (los $6\,119$
vectores de título$+$texto) contra la pregunta \emph{entera}, corre
siempre y para todos los pasajes, y aporta masa con peso pequeño
($0{,}05$). En la práctica, los dos eslabones débiles medidos del
enlazador son el \emph{recorte} del tramo y el \emph{objeto} contra el
que se compara (párrafo anterior), no la densidad: en el estudio de
errores retenido, los fallos atribuibles al enlazador ($3$ de $15$)
nacieron de tramos rotos por la regla de mayúsculas (\guillemotleft{}Magician
(1958\guillemotright{}), que tampoco alcanzan el umbral contra la ficha correcta. De ahí
que las evoluciones previstas (párrafo siguiente) sean sustituir la
detección de tramos por reconocimiento de entidades y comparar de nombre
a nombre, no añadir más densidad.

\paragraph{Evoluciones del enlazador.}
Las alternativas para el escalón 2 se han medido sobre los mismos $275$
tramos con dos criterios: si el mejor candidato pertenece a un pasaje oro
(calidad del \emph{señalador}) y cuántos aciertos y errores produce un
umbral de aceptación (calidad del \emph{decisor});
la tabla~\ref{tab:enlazador-alt} las resume. Tres conclusiones. Primera,
el problema no es de tipografía sino de \emph{truncamiento}: la distancia
de edición al estilo SymSpell (distancia absoluta $\le 2$) acepta $78$
tramos y acierta $33$ ---\guillemotleft{}Gaby\guillemotright{} $\to$ \guillemotleft{}Gabby\guillemotright{}---, porque un
artículo o una palabra perdidos distan cuatro o más caracteres de la
forma correcta; es la herramienta equivocada para este fallo. Segunda,
BM25 sobre los alias es el mejor señalador ($227$ de $275$ con el oro en
primer lugar) pero un mal decisor: su puntuación no acota bien la
aceptación ($167$ correctos frente a $27$ erróneos al $0{,}8$ de su
máximo posible, porque premia coincidencias parciales como \guillemotleft{}Gaby\guillemotright{}
$\to$ \guillemotleft{}Gaby Brimer\guillemotright{}); y como sustituto del escalón 1 no aporta nada:
coincide con el alias exacto en $1\,239$ de los $1\,259$ tramos resueltos,
y las $20$ discrepancias son tramos con paréntesis roto (\guillemotleft{}Possession
(1922\guillemotright{}) en los que BM25 casa con el año. Tercera, las dos medidas de
semejanza \emph{de nombre a nombre} ---el coeficiente de Dice sobre
trigramas de caracteres, el mismo escalón difuso del enlazador clínico
(sección~\ref{sec:gesida-index}), y el coseno contra la incrustación del
nombre solo--- son las que rinden, y se complementan: la léxica captura
truncamientos y erratas sin coste; la vectorial, variantes que no
comparten superficie. Su unión (Dice $\ge 0{,}80$ o coseno de nombre
$\ge 0{,}85$) enlaza $187$ de los $275$ tramos con $179$ aciertos y $8$
errores, frente a los $3$ de la versión evaluada; los errores son
casi-homónimos (\guillemotleft{}Two Brides\guillemotright{} $\to$ \guillemotleft{}The Two Brides\guillemotright{} cuando el
oro es \emph{Two Brides and a Baby}) que la ficha con descripción puede
desempatar. De los $88$ que quedan sin enlazar, $79$ son tramos rotos de
una o dos palabras (\guillemotleft{}Hitman\guillemotright{}, \guillemotleft{}True Story\guillemotright{}, \guillemotleft{}Dr\guillemotright{}): ninguna
medida de semejanza los rescata, porque el error está antes. Con esto,
las evoluciones previstas, en orden:
\begin{enumerate}[leftmargin=1.8em, itemsep=1pt, label=L\arabic*.]
\item \textbf{Reconocimiento de entidades sobre la pregunta} en lugar del
  recorte por mayúsculas: un modelo de lenguaje (una llamada por
  consulta, como hacen HippoRAG y CatRAG; sección~\ref{sec:palancas}) o
  un reconocedor local sin coste de API. Es la única corrección general
  del recorte ---artículos, dos puntos, coordinación, paréntesis--- y
  ataca los $79$ tramos rotos. Un parche de prefijos (reintentar el alias
  anteponiendo \guillemotleft{}the\guillemotright{}) resolvería $168$ de los $275$ tramos en este
  banco, pero es específico del inglés y de los títulos enciclopédicos, y
  se descarta por no generalizable.
\item \textbf{Escalón 2 de nombre a nombre}: Dice sobre trigramas
  $\ge 0{,}80$ o coseno de nombre $\ge 0{,}85$ contra los $31\,837$ alias,
  con la ficha con descripción como desempate entre alias casi iguales
  de entradas distintas. Con ello el enlazador converge con el clínico
  (exacto $\to$ trigramas $\to$ vector $\to$ juez) y el umbral deja de
  depender de la longitud de la descripción.
\item \textbf{BM25 solo como generador de candidatos}: los $k$ mejores
  alias para el desempate o para un juez, nunca como criterio de
  aceptación.
\item \textbf{Enlazador consciente de colisiones} (E7 de la
  sección~\ref{sec:memoria}).
\end{enumerate}
Cada una se validará en el sondeo y se medirá una sola vez en el banco.

\begin{table}[t]
\centering\footnotesize\setlength{\tabcolsep}{3.5pt}
\caption{Alternativas para el escalón 2 del enlazador, medidas sobre los
$275$ tramos del banco que no casan por alias exacto. \guillemotleft{}Oro en
1.\textsuperscript{o}\guillemotright{}: el mejor candidato pertenece a un pasaje oro.
Entre paréntesis, correctos\,/\,erróneos entre los aceptados. Ninguna
alternativa está en la versión evaluada.}
\label{tab:enlazador-alt}
\scriptsize\setlength{\tabcolsep}{3pt}
\begin{tabular}{llccrl}
\toprule
método & objeto comparado & oro en 1.\textsuperscript{o} & aceptación & acepta & coste\\
\midrule
coseno contra ficha (evaluado) & nombre (alias) --- descripción & $216$ & $\ge 0{,}80$ & $3$ ($3/0$) & 1 incr.\\
ídem, umbral rebajado & nombre (alias) --- descripción & $216$ & $\ge 0{,}70$ & $43$ ($42/1$) & 1 incr.\\
coseno contra nombre solo & alias & $222$ & $\ge 0{,}85$ & $152$ ($146/6$) & 1 incr.\\
Dice sobre trigramas & alias plegado & $222$ & $\ge 0{,}80$ & $180$ ($174/6$) & 0\\
BM25 & alias tokenizado & $227$ & $\ge 0{,}8$ del máx. & $194$ ($167/27$) & 0\\
edición absoluta (SymSpell) & alias plegado & $167$ & $\le 2$ & $78$ ($33/45$) & 0\\
edición normalizada & alias plegado & $167$ & $\ge 0{,}80$ & $105$ ($90/15$) & 0\\
\midrule
unión: Dice $\ge 0{,}80$ $\vee$ cos.\ nombre $\ge 0{,}85$ & alias & --- & --- & $187$ ($179/8$) & 1 incr.\\
\bottomrule
\end{tabular}
\end{table}

El resultado de la etapa es el conjunto de entidades enlazadas
$L(q) \subseteq \mathcal{E}$, que alimenta tres mecanismos: los candidatos
de frontera del filtro (sección~\ref{sec:filtro}) y las anclas fuertes y la
masa del pasaje propio en la siembra (sección~\ref{sec:siembra}).

\emph{En el ejemplo conductor}: la regla de capitalización detecta dos
tramos ---\guillemotleft{}God'S Gift To Women\guillemotright{} y \guillemotleft{}Aldri Annet Enn Bråk\guillemotright{}; el
\guillemotleft{}or\guillemotright{} minúsculo no es conector y los separa--- y ambos resuelven por
alias exacto tras el plegado, sin necesidad del respaldo vectorial:
$L(q) = \{\texttt{can:god's gift to women},\;
\texttt{can:aldri annet enn br\aa k}\}$. La etapa entrega dos anclas
exactas en el grafo; nótese que los \emph{directores} aún no aparecen por
ningún sitio: descubrirlos es trabajo de las dos etapas siguientes.

\subsubsection{Vectores de consulta y afinidad pregunta--hecho}
\label{sec:afinidad}

Sea $\{v_0,\dots,v_J\}$ el conjunto de vectores de consulta. Por defecto
$J=0$ y $v_0 = e(q)$. En las preguntas de comparación, una llamada al modelo
de lenguaje descompone la pregunta en $2$--$4$ subconsultas de un salto
(\guillemotleft{}fecha de nacimiento del director de $A$\guillemotright{} $\to$ \guillemotleft{}director de $A$\guillemotright{};
la entidad puente desconocida se denota con el comodín $[X]$), y cada
subconsulta aporta un vector adicional. La afinidad de cada aserción con la
consulta es el mejor coseno sobre los vectores:
\begin{equation}
\sigma_a \;=\; \max_{0\le j\le J}\;\bigl\langle v_j,\, e(\tau_a)\bigr\rangle .
\label{eq:afinidad}
\end{equation}

\paragraph{Qué es $J$ y por qué el máximo.}
$J$ es el número de vectores de consulta \emph{adicionales}: $v_0$ existe
siempre y, en la estrategia de descomposición, cada subconsulta aporta
uno más. Una pregunta real de comparación del banco lo muestra: para
\guillemotleft{}\emph{Are Christopher Newton (Criminal) and Frances M. Vega of the
same nationality?}\guillemotright{} (oro: las dos biografías) el modelo devolvió tres
subconsultas ---\guillemotleft{}What is the nationality of Christopher Newton?\guillemotright{},
\guillemotleft{}What is the nationality of Frances M. Vega?\guillemotright{} y \guillemotleft{}Are the
nationalities of Christopher Newton and Frances M. Vega the same?\guillemotright{}---,
luego $J=3$ y cuatro vectores. La pregunta entera, $v_0$, mezcla las dos
ramas y puntúa flojo: su mejor hecho, \guillemotleft{}Frances Marie Vega was of
Puerto Rican descent\guillemotright{}, queda a $0{,}525$ y el mejor de Newton a
$0{,}513$. Cada subconsulta puntúa \emph{su} rama con nitidez ($v_1$: los
hechos de Newton a $0{,}69$; $v_2$: los de Vega a $0{,}73$). Con el
máximo de la ecuación~\eqref{eq:afinidad}, los ocho hechos de mayor
$\sigma$ salen alternando ramas ---cinco de Vega por $v_2$, tres de
Newton por $v_1$--- en lugar de que la rama léxicamente dominante deje
sin candidatos a la otra. La tercera subconsulta, redundante con $v_0$,
no gana el máximo en ningún hecho relevante: una subconsulta inútil no
cuesta nada. En esta estrategia no corre el filtro
(sección~\ref{sec:router}), de modo que $\tilde\sigma=\sigma$ y la
siembra usa la compuerta blanda; la sección~\ref{sec:siembra} desglosa
el vector de siembra que resulta de esta misma pregunta.

\emph{En el ejemplo conductor} (estrategia de filtro: $J=0$, un solo
vector): el máximo de $\sigma$ lo alcanza el hecho puente de la primera
película, \guillemotleft{}God's Gift to Women was directed by Michael Curtiz\guillemotright{}
($\sigma=0{,}554$), pero los siete puestos siguientes son hechos de esa
misma película (estreno, reparto, obra de base). El hecho puente de la
\emph{otra} rama entra más abajo y verbalizado sin la palabra
\guillemotleft{}directed\guillemotright{}: \guillemotleft{}Edith Carlmar is known for the film `Aldri annet enn
bråk'\guillemotright{} ($\sigma=0{,}480$). Es el sesgo característico del coseno:
favorece la rama léxicamente dominante de la pregunta y puede matar de
hambre a la otra. La etapa entrega el vector completo de afinidades
$\sigma$; corregir ese sesgo es exactamente el trabajo de la siguiente.

\subsubsection{Filtro selector de hechos (memoria de reconocimiento)}
\label{sec:filtro}

El coseno de la ecuación~\eqref{eq:afinidad} es una señal de \emph{semejanza},
no de \emph{utilidad}: en una pregunta composicional, el hecho puente
(\guillemotleft{}$X$ dirigió la película $P$\guillemotright{}) compite con hechos superficialmente
parecidos de otras películas. Siguiendo la idea de memoria de
reconocimiento de \hippo{}, un modelo de lenguaje actúa de juez selectivo
sobre una lista corta de candidatos. La lista combina dos fuentes:
\begin{equation}
\mathcal{K}(q) \;=\;
\underbrace{\op{Top}_{40}\{\sigma_a\}}_{\text{candidatos por coseno}}
\;\cup\;
\underbrace{\bigcup_{\varepsilon\in L(q)}
\op{Top}_{15}\bigl\{\sigma_a : a \ni \varepsilon\bigr\}}_{\text{frontera de las entidades enlazadas}},
\label{eq:candidatos}
\end{equation}
donde la \emph{frontera} de una entidad enlazada $\varepsilon$ es el
conjunto de hechos que inciden en ella en el grafo canónico ---las
aserciones de las que $\varepsilon$ es sujeto, objeto o participante,
es decir, sus vecinos a una arista en $G^{\!\rightarrow}$---, de los que
se toman los $15$ de mayor $\sigma$. La lista es una \emph{unión}, no
una intersección: al Top-40 se le añaden los hechos de frontera que aún
no estaban en él, sin duplicar y con un tope de $20$ añadidos en total.
Las dos fuentes son de naturaleza distinta, y por eso se suman: el
Top-40 selecciona por \emph{semejanza con la pregunta} entre los
$47\,999$ hechos, sin mirar el grafo; la frontera selecciona por
\emph{estructura} ---qué afirma el corpus sobre las entidades que la
pregunta nombra---, y el coseno solo ordena dentro de cada entidad. La
frontera garantiza que el hecho puente de una entidad enlazada llegue al
filtro aunque esté verbalizado sin el vocabulario de la pregunta.

\emph{Medido sobre el banco} (tomando como hecho oro toda aserción de un
pasaje oro que solapa al menos el $60\%$ de los términos de un triplete
de evidencia; $966$ preguntas tienen alguno): $|\mathcal{K}(q)|$ vale
$41{,}4$ de media, es decir, la frontera añade $1{,}4$ hechos nuevos por
pregunta; el Top-40 ya contiene algún hecho oro en $944$ preguntas; la
frontera aporta un hecho oro de primer salto que el Top-40 no tenía en
$51$ ($5\%$: $22$ composicionales, $15$ de comparación-puente, $10$ de
inferencia, $4$ de comparación), y solo en $19$ preguntas $\mathcal{K}(q)$
no contiene ningún hecho oro. De los $3\,371$ hechos oro que inciden en
una entidad enlazada, el $88\%$ está en el Top-40, el $2{,}5\%$ entra por
la frontera y el $9\%$ queda fuera de $\mathcal{K}(q)$ (entidades con más
de $15$ hechos, o sin enlazar: $74$ preguntas no enlazan ninguna
entidad). La frontera es, por tanto, una red de seguridad estructural
de coste cero, útil en una minoría medible, no el mecanismo principal.
En el ejemplo conductor, la frontera de \texttt{can:aldri annet enn
bråk} son sus $7$ hechos incidentes, cinco de ellos fuera del Top-40
(rangos $114$ a $1\,416$ por coseno: \guillemotleft{}was edited by Edith
Carlmar\guillemotright{}, \guillemotleft{}A 1955 Danish remake\dots\guillemotright{}); el hecho puente de esa
rama ya estaba en el Top-40 (rango $5$), de modo que allí la frontera
amplía el contexto del juez pero no decide. En la traza de la
sección~\ref{sec:traza} ocurre lo mismo: $8$ hechos nuevos de
\texttt{can:atomised}, con el puente en rango $1$.

El filtro recibe la pregunta y las frases $\tau_a$ numeradas y devuelve una
selección $S(q) \subseteq \mathcal{K}(q)$ con $|S(q)| \le 4$. La consigna,
con tres demostraciones fijas, codifica dos políticas aprendidas del
análisis de fallos: \emph{conservar los saltos intermedios} (el hecho \guillemotleft{}la
película fue dirigida por $X$\guillemotright{} es útil aunque la pregunta pida la familia de
$X$) y, cuando un papel de la pregunta (intérprete, director, fundador)
tiene varios candidatos en los hechos, \emph{conservarlos todos}, con
prohibición explícita de resolver la ambigüedad con conocimiento paramétrico
del modelo. Esta última política corrige un modo de fallo observado: un juez
sin ella descartaba el hecho puente verdadero y \guillemotleft{}respondía de memoria\guillemotright{}. La
selección define la afinidad corregida
\begin{equation}
\tilde{\sigma}_a \;=\;
\begin{cases}
1 & \text{si } a \in S(q)\\
\sigma_a & \text{en otro caso.}
\end{cases}
\label{eq:sigma-tilde}
\end{equation}

\emph{En el ejemplo conductor}: la lista de candidatos $\mathcal{K}(q)$
reúne el Top-40 por coseno y los $16$ hechos de frontera de las dos
películas enlazadas ($9$ de una, $7$ de la otra), de los que $7$ no
estaban en el Top-40: $47$ candidatos numerados. El filtro selecciona
tres hechos:
\begin{center}\small
\begin{tabular}{lll}
\toprule
hecho aprobado & $\sigma$ & entidades que aporta\\
\midrule
\guillemotleft{}God's Gift to Women was directed by Michael Curtiz\guillemotright{} & $0{,}554$ &
  la película 1 y \textbf{Curtiz}\\
\guillemotleft{}Edith Carlmar is known for the film `Aldri annet enn bråk'\guillemotright{} &
  $0{,}480$ & la película 2 y \textbf{Carlmar}\\
\guillemotleft{}Edith Carlmar was Norway's first female film director\guillemotright{} &
  $0{,}400$ & Carlmar y Norway\\
\bottomrule
\end{tabular}
\end{center}
Aquí se ve la política de \guillemotleft{}conservar todos los candidatos de un
papel\guillemotright{} haciendo su trabajo: pese a que el coseno favorecía una sola rama,
el filtro retiene el hecho de dirección de \emph{ambas} películas ---el
segundo, además, verbalizado sin la palabra \guillemotleft{}directed\guillemotright{}, que un patrón
léxico no habría reconocido---. También se ve su imperfección tolerable:
el tercer hecho es superfluo y arrastrará una entidad ruidosa
(\texttt{can:norway}); las etapas siguientes lo absorben.

\paragraph{Qué entrega esta etapa.}
La salida del filtro es el par $(S(q),\,\tilde\sigma)$, y cada mitad
alimenta una etapa distinta: $S(q)$ decide \emph{quién siembra} ---en la
sección~\ref{sec:siembra}, las entidades de los hechos aprobados reciben
masa de reinicio, con lo que los \emph{directores}, que no estaban en la
pregunta ni en $L(q)$, entran por primera vez en escena--- y
$\tilde\sigma$ decide \emph{cuánto pesa cada hecho como peaje del paseo}
---en la sección~\ref{sec:ppr}, vía la modulación $\mu$---. Este doble
destino es el eslabón que une las tres últimas etapas.

\subsubsection{Siembra: el vector de personalización}
\label{sec:siembra}

El vector de personalización $\mathbf{v}\in\R^{|V|}$ concentra la masa de
reinicio del paseo. Recibe los productos de las tres etapas anteriores
---la selección $S(q)$ del filtro, las entidades enlazadas $L(q)$ del
enlazador y los vectores de consulta--- y los combina en cuatro fuentes
(figura~\ref{fig:traza}):

\paragraph{(a) Compuerta dura sobre los hechos aprobados.}
Solo siembran las entidades que participan en hechos seleccionados por el
filtro:
\begin{equation}
\mathbf{v}(\varepsilon) \;=\;
\frac{1}{|S_\varepsilon|}\sum_{a\in S_\varepsilon}
\op{clip}_{[0,1]}(\tilde{\sigma}_a),
\qquad
S_\varepsilon = \{a \in S(q) : \varepsilon \text{ participa en } a\},
\label{eq:compuerta}
\end{equation}
que con la ecuación~\eqref{eq:sigma-tilde} vale $1$ para toda entidad de un
hecho aprobado. Si el filtro no aprueba nada, no hay semillas de hechos y la
ordenación degrada con suavidad al recuperador denso (fuente~b). En la
variante de descomposición (sin filtro) la compuerta es blanda: siembran las
entidades de los $15$ hechos de mayor $\sigma$, con peso igual a la media de
$\sigma$ normalizado dividida por la frecuencia documental de la entidad
(corrección anti-concentrador, que evita que entidades muy conectadas
acaparen la masa), reteniendo las $5$ mejores.

\paragraph{(b) Presa densa sobre todos los pasajes.}
Todo nodo-pasaje recibe masa proporcional a su similitud densa con la
pregunta, normalizada al rango $[0,1]$ y escalada por
$w_{\text{pas}} = 0{,}05$:
$\mathbf{v}(c) \mathrel{+}= w_{\text{pas}}\cdot
\widehat{\langle v_0, e(c)\rangle}$. Esta fusión con el recuperador denso
sigue el diseño de \hippo{} (todos los pasajes como semillas) y garantiza
que el sistema nunca rinda por debajo del recuperador denso puro.

\paragraph{(c) Anclas simbólicas débiles.}
Las dos entidades más afines a cada vector de consulta (por coseno con la
ficha) reciben una masa pequeña $\epsilon = 0{,}05$, que preserva
alcanzabilidad cuando ni el filtro ni el enlazador aciertan.

\paragraph{(d) Anclas del enlazador y pasaje propio.}
Toda entidad enlazada desde la pregunta recibe al menos
$w_{\text{enl}} = 0{,}5$:
$\mathbf{v}(\varepsilon) \leftarrow
\max\{\mathbf{v}(\varepsilon),\, w_{\text{enl}}\}$ para
$\varepsilon \in L(q)$. Además, para \emph{toda} entidad sembrada (enlazada
o aprobada por el filtro) con pasaje propio $\pi(\varepsilon)$:
\begin{equation}
\mathbf{v}\bigl(\pi(\varepsilon)\bigr) \;\mathrel{+}=\;
w_{\text{enl}}\cdot \mathbf{v}(\varepsilon).
\label{eq:pasaje-propio}
\end{equation}
Este término codifica el sesgo enciclopédico \guillemotleft{}el artículo de una entidad
\emph{es} la entidad\guillemotright{}: si el hecho aprobado introduce al director $X$, la
biografía de $X$ ---que suele contener la respuesta del segundo salto--- recibe
masa de reinicio directa, sin depender de que el paseo la alcance. En el
análisis de errores previo, la ausencia de este término explicaba el
$57\%$ de los fallos de cadena: las entidades concentradoras repartían su
masa entre todos los pasajes que las mencionan y el pasaje propio quedaba
en rangos 5--10.

Finalmente $\mathbf{v} \leftarrow \mathbf{v} / \norm{\mathbf{v}}_1$.

\emph{En el ejemplo conductor}: las cinco entidades de los hechos
aprobados quedan sembradas a peso $\approx 1$ (las dos películas, los dos
directores y la ruidosa \texttt{can:norway}; las anclas débiles añaden
$0{,}05$ a dos de ellas), y el término de pasaje propio deposita masa
directa en los \emph{cuatro} pasajes oro: $0{,}575$ y $0{,}545$ en las dos
películas, $0{,}555$ y $0{,}546$ en las dos biografías. Obsérvese el
relevo entre etapas: los directores entraron por el filtro, y este término
convierte ese descubrimiento en masa sobre sus artículos ---los pasajes
del segundo salto--- sin esperar al paseo. \texttt{can:norway} no tiene
pasaje propio en el corpus, así que su siembra solo podrá diluirse. La
etapa entrega el vector $\mathbf{v}$ normalizado.

\paragraph{De los tres productos al vector, paso a paso.}
La siembra no ve la pregunta: ve tres productos de las etapas anteriores.
Conviene fijar qué es cada uno y cómo se produjo en la pregunta de
comparación de la sección~\ref{sec:afinidad}, antes de recorrer el
algoritmo sobre ella.
\begin{description}[leftmargin=1.6em, itemsep=1pt, style=unboxed]
\item[P1, del enlazador (\S\ref{sec:linker}): $L(q)$.] La regla de
  mayúsculas produce un único tramo, \guillemotleft{}Christopher Newton (Criminal)
  and Frances M Vega\guillemotright{}; como no es alias, se parte por la coordinación
  en dos; cada mitad se pliega (minúsculas, sin paréntesis ni
  puntuación: \guillemotleft{}christopher newton\guillemotright{}, \guillemotleft{}frances m vega\guillemotright{}) y casa en
  la tabla de alias, sin respaldo vectorial:
  $L(q)=\{\texttt{can:christopher newton},\ \texttt{can:frances m.
  vega}\}$. Obsérvese que el plegado descarta el desambiguador
  \guillemotleft{}(Criminal)\guillemotright{}: aunque el diccionario tuviera dos entradas, el
  enlazador no podría distinguirlas (evolución L4).
\item[P2, de la afinidad (\S\ref{sec:afinidad}): $\{v_j\}$ y $\sigma$.]
  Cuatro vectores ($v_0$ más las tres subconsultas) y, para cada uno de
  los $47\,999$ hechos, $\sigma_a=\max_j\langle v_j,e(\tau_a)\rangle$,
  con mínimo $-0{,}074$ y máximo $0{,}730$.
\item[P3, del filtro (\S\ref{sec:filtro}): $S(q)$ y $\tilde\sigma$.] En
  la estrategia de descomposición el filtro no corre: $S(q)$ está vacío,
  $\tilde\sigma=\sigma$ y la fuente~(a) usa la compuerta blanda. (En la
  estrategia de filtro, $S(q)$ son los hechos aprobados y
  $\tilde\sigma_a=1$ para ellos; el ejemplo conductor recorre ese caso.)
\end{description}
Con esos tres productos, $\mathbf{v}$ se construye en cuatro pasos y una
normalización:
\begin{enumerate}[leftmargin=1.6em, itemsep=2pt]
\item \textbf{Hechos $\to$ entidades (fuente~a, compuerta blanda).} Se
  normaliza $\sigma$ al rango $[0,1]$,
  $n_a=(\sigma_a-\min\sigma)/(\max\sigma-\min\sigma)$, y se toman los
  $15$ hechos de mayor $\sigma$ (tabla~\ref{tab:siembra-pasos}). Cada
  entidad participante acumula $n_a/\op{frec}(\varepsilon)$ por hecho,
  donde $\op{frec}$ es el número de pasajes que la mencionan, y su peso
  es la media sobre sus hechos. Así, \texttt{can:frances m. vega} suma
  $3{,}817$ en $4$ hechos $\to 0{,}954$; \texttt{can:david newton} suma
  $4{,}403$ en $5$ $\to 0{,}881$; \texttt{can:christopher newton} suma
  $2{,}706$ en $6$ hechos, pero cada aportación va dividida por su
  frecuencia $2$ $\to 0{,}451$; los lugares entran por un solo hecho
  (\texttt{oakland} $0{,}294$, \texttt{north carolina} $0{,}127$), y
  \texttt{california} ($\op{frec}=62$) y \texttt{new york city}
  ($\op{frec}=96$) quedan aplastados por la corrección
  anti-concentrador ($0{,}014$ y $0{,}009$). Se retienen las $5$
  mejores.
\item \textbf{Presa densa (fuente~b).} $\langle v_0,e(c)\rangle$ para los
  $6\,119$ pasajes (mínimo $-0{,}118$; máximo $0{,}445$, en
  \emph{Christopher Newton}), normalizado al rango $[0,1]$ y escalado
  por $0{,}05$: \emph{Frances M. Vega} $0{,}443\to0{,}996\to0{,}050$;
  \emph{Christopher Newton (criminal)} $0{,}428\to0{,}970\to0{,}049$;
  \emph{David Newton (artist)} $0{,}372\to0{,}871\to0{,}044$. Como la
  media normalizada es $0{,}45$, la fuente suma $138{,}1$ unidades
  repartidas por todo el corpus.
\item \textbf{Anclas débiles (fuente~c).} Para cada $v_j$, las dos fichas
  canónicas más afines reciben $+0{,}05$: $v_0\to$
  \texttt{christopher newton} ($0{,}492$), \texttt{frances m. vega}
  ($0{,}490$); $v_1\to$ \texttt{christopher newton} ($0{,}678$),
  \texttt{david newton} ($0{,}580$); $v_2\to$ \texttt{frances m. vega}
  ($0{,}729$), \texttt{iraq war} ($0{,}544$); $v_3\to$ \texttt{frances
  m. vega} ($0{,}525$), \texttt{christopher newton} ($0{,}410$). Newton
  y Vega acumulan $+0{,}150$ cada uno.
\item \textbf{Enlazador y pasaje propio (fuente~d).} Las dos entidades de
  $L(q)$ ya superan $0{,}5$, así que el máximo no actúa. Después, toda
  entidad sembrada con pasaje propio deposita
  $0{,}5\cdot\mathbf{v}(\varepsilon)$ en él: \emph{Frances M. Vega}
  recibe $0{,}5\times1{,}104=0{,}552$; \emph{David Newton (artist)},
  $0{,}5\times0{,}931=0{,}465$; \emph{Christopher Newton} ---el pasaje
  propio de la entrada fundida, es decir, la biografía del director
  teatral---, $0{,}5\times0{,}601=0{,}301$. Los lugares no tienen pasaje
  propio, y \emph{Christopher Newton (criminal)} no es pasaje propio de
  ninguna entrada: se queda con su presa densa.
\item \textbf{Normalización.}
  $\norm{\mathbf{v}}_1=2{,}71+138{,}06+0{,}40+1{,}32=142{,}49$ y
  $\mathbf{v}\leftarrow\mathbf{v}/142{,}49$. El resultado es la
  tabla~\ref{tab:siembra-ejemplo}.
\end{enumerate}

\begin{table}[t]
\centering\scriptsize\setlength{\tabcolsep}{3pt}
\caption{Paso 1 de la siembra en la pregunta de comparación: los $15$
hechos de mayor $\sigma$, con su afinidad normalizada $n_a$, el vector
de consulta que dio el máximo y sus entidades canónicas (entre
paréntesis, la frecuencia documental que divide su aportación). Los seis
hechos de \texttt{christopher newton} mezclan al actor canadiense y al
asesino: la entrada está fundida.}
\label{tab:siembra-pasos}
\begin{tabular}{rcccp{0.40\textwidth}p{0.24\textwidth}}
\toprule
& $\sigma_a$ & $n_a$ & $v_j$ & hecho $\tau_a$ & entidades (frec.)\\
\midrule
1 & $0{,}730$ & $1{,}000$ & $v_2$ & Frances Marie Vega was born on September 2, 1983 & frances m. vega (1)\\
2 & $0{,}723$ & $0{,}991$ & $v_2$ & Frances Marie Vega was of Puerto Rican descent & frances m. vega (1)\\
3 & $0{,}692$ & $0{,}953$ & $v_1$ & Christopher J. Newton was born on November 13, 1969 & christopher newton (2)\\
4 & $0{,}667$ & $0{,}922$ & $v_1$ & Christopher J. Newton was an American murderer & christopher newton (2)\\
5 & $0{,}664$ & $0{,}919$ & $v_2$ & Frances Marie Vega was a United States Army soldier & frances m. vega (1); united states army (9)\\
6 & $0{,}656$ & $0{,}908$ & $v_1$ & Christopher Newton is Canadian & christopher newton (2)\\
7 & $0{,}655$ & $0{,}907$ & $v_2$ & Frances Marie Vega died on November 2, 2003 & frances m. vega (1)\\
8 & $0{,}649$ & $0{,}899$ & $v_1$ & David Christopher Newton was an American sculptor, painter, \dots & david newton (1)\\
9 & $0{,}645$ & $0{,}895$ & $v_1$ & Christopher Newton is an actor & christopher newton (2)\\
10 & $0{,}641$ & $0{,}890$ & $v_1$ & David Christopher Newton lived and worked in North Carolina & david newton (1); north carolina (7)\\
11 & $0{,}634$ & $0{,}881$ & $v_1$ & David Christopher Newton was born in December 1953 in Oakland, California & david newton (1); california (62); oakland (3)\\
12 & $0{,}625$ & $0{,}870$ & $v_1$ & Christopher J. Newton died on May 24, 2007 & christopher newton (2)\\
13 & $0{,}623$ & $0{,}868$ & $v_1$ & David Christopher Newton lived and worked in Massachusetts & david newton (1); massachusetts (8)\\
14 & $0{,}622$ & $0{,}865$ & $v_1$ & David Christopher Newton lived and worked in New York & david newton (1); new york city (96)\\
15 & $0{,}621$ & $0{,}865$ & $v_1$ & Christopher J. Newton was a murderer & christopher newton (2)\\
\bottomrule
\end{tabular}
\end{table}

\paragraph{Un vector de siembra, desglosado.}
La tabla~\ref{tab:siembra-ejemplo} muestra $\mathbf{v}$ nodo a nodo y
fuente a fuente para la pregunta de comparación de la
sección~\ref{sec:afinidad}, instrumentada exactamente como se evaluó (el
desglose reproduce la salida del sistema con diferencia $<10^{-15}$), y
debajo el desglose agregado del ejemplo conductor. Cuatro lecturas.
(i)~\emph{La masa total la lleva la presa densa}: de las $142{,}5$
unidades sin normalizar, $138{,}1$ ($97\%$) son la fuente~(b),
repartidas entre $6\,118$ pasajes ($0{,}023$ de media, es decir
$1{,}6\cdot10^{-4}$ tras normalizar); en el ejemplo conductor, el $94\%$.
Las fuentes estructurales (a), (c) y (d) suman el $3$--$6\%$ restante,
pero concentrado en una decena de nodos: cada entidad sembrada por hechos
recibe entre $25$ y $50$ veces la masa de un pasaje medio, y cada pasaje
propio, unas $25$--$30$. La presa densa es, en suma, un \emph{suelo}
uniforme que garantiza que ningún pasaje tenga masa cero y que el sistema
no rinda por debajo del recuperador denso, no una señal que compita con
las semillas.
(ii)~\emph{La compuerta blanda no discrimina}: sin filtro, el hecho
\guillemotleft{}David Christopher Newton was an American sculptor\guillemotright{}
($\sigma=0{,}649$ con $v_1$) siembra a un tercer homónimo,
\texttt{can:david newton}, con $0{,}881$ ---más que al propio Christopher
Newton, $0{,}451$, penalizado por su frecuencia documental $2$--- y su
biografía acaba segunda en el paseo. Es el precio de la variante de
descomposición y un argumento para llevar también a ella el filtro
(sección~\ref{sec:falla-aun}).
(iii)~\emph{El pasaje propio hereda los errores del diccionario}: la
entrada \texttt{can:christopher newton} funde al director teatral
canadiense (n.~1936, artículo \emph{Christopher Newton}) con el asesino
Christopher J. Newton (n.~1969, artículo \emph{Christopher Newton
(criminal)}), y su pasaje propio es la biografía equivocada; por eso el
pasaje oro del criminal recibe únicamente su presa densa ($0{,}049$)
(sección~\ref{sec:diccionario}, limitación~1).
(iv)~\emph{Qué decide el orden}: la cadena completa se cumple
($\mathrm{FC@}5=1$: \emph{Frances M. Vega} primera, \emph{Christopher
Newton (criminal)} cuarta), pero el pasaje del criminal llega por el
canal de menciones de la entidad fundida y por su presa densa, no por la
siembra estructural; en el ejemplo conductor, en cambio, los cuatro
pasajes oro reciben masa directa de la fuente~(d) ($0{,}50$--$0{,}53$
cada uno; $0{,}0043$--$0{,}0045$ normalizados) y salen en los puestos
1--4.

\begin{table}[t]
\centering\footnotesize\setlength{\tabcolsep}{4pt}
\caption{El vector de siembra $\mathbf{v}$ de la pregunta \guillemotleft{}Are
Christopher Newton (Criminal) and Frances M. Vega of the same
nationality?\guillemotright{} (estrategia de descomposición, sin filtro), nodo a nodo
y fuente a fuente: (a)~hechos (compuerta blanda), (b)~presa densa,
(c)~anclas débiles, (d)~ancla del enlazador y pasaje propio. Debajo, el
desglose agregado del ejemplo conductor (estrategia de filtro). Valores
sin normalizar salvo la última columna.}
\label{tab:siembra-ejemplo}
\begin{tabular}{lrrrrrr}
\toprule
nodo & (a) & (b) & (c) & (d) & total & $\mathbf{v}$ norm.\\
\midrule
\texttt{can:frances m. vega} & $0{,}954$ & & $0{,}150$ & & $1{,}104$ & $0{,}0077$\\
\texttt{can:david newton} (homónimo) & $0{,}881$ & & $0{,}050$ & & $0{,}931$ & $0{,}0065$\\
pasaje \emph{Frances M. Vega} (oro) & & $0{,}050$ & & $0{,}552$ & $0{,}602$ & $0{,}0042$\\
\texttt{can:christopher newton} (fundida) & $0{,}451$ & & $0{,}150$ & & $0{,}601$ & $0{,}0042$\\
pasaje \emph{David Newton (artist)} & & $0{,}044$ & & $0{,}465$ & $0{,}509$ & $0{,}0036$\\
pasaje \emph{Christopher Newton} (director teatral) & & $0{,}050$ & & $0{,}301$ & $0{,}351$ & $0{,}0025$\\
\texttt{can:oakland} & $0{,}294$ & & & & $0{,}294$ & $0{,}0021$\\
\texttt{can:north carolina} & $0{,}127$ & & & & $0{,}127$ & $0{,}0009$\\
\texttt{can:iraq war} & & & $0{,}050$ & & $0{,}050$ & $0{,}0004$\\
pasaje \emph{Christopher Newton (criminal)} (oro) & & $0{,}049$ & & & $0{,}049$ & $0{,}0003$\\
otros $6\,114$ pasajes (solo presa densa) & & $137{,}9$ & & & $137{,}9$ & $0{,}968$\\
\midrule
suma por fuente & $2{,}71$ & $138{,}06$ & $0{,}40$ & $1{,}32$ & $142{,}49$ & $1$\\
\midrule
\multicolumn{7}{l}{\emph{ejemplo conductor} (filtro): $5$ entidades a peso $1$ (dos con ancla débil, $1{,}05$); $4$ pasajes propios a $0{,}50$--$0{,}53$}\\
suma por fuente & $5{,}00$ & $119{,}82$ & $0{,}10$ & $2{,}05$ & $126{,}97$ & $1$\\
\bottomrule
\end{tabular}
\end{table}

\subsubsection{Paseo aleatorio con reinicio y transiciones moduladas}
\label{sec:ppr}

\paragraph{De la adyacencia a la matriz de transición.}
Sea $A\in\R^{|V|\times|V|}$ la matriz de adyacencia de
$G^{\!\rightarrow}$ ($A_{uw}=1$ si existe la arista $u\to w$). Antes del
paseo, las \emph{columnas} correspondientes a nodos-aserción se reescalan
por la afinidad del hecho con la consulta:
\begin{equation}
\tilde{A} \;=\; A\,\op{diag}(\mu),
\qquad
\mu_u \;=\;
\begin{cases}
\dfrac14 + \dfrac34\,\op{clip}_{[0,1]}(\tilde{\sigma}_u)
   & \text{si } u \in \mathcal{A}\\[6pt]
1  & \text{en otro caso,}
\end{cases}
\label{eq:modulacion}
\end{equation}
y sobre la matriz modulada se normaliza por filas:
\begin{equation}
P \;=\; D^{-1}\tilde{A},
\qquad
D \;=\; \op{diag}\bigl(\tilde{A}\,\mathbf{1}\bigr),
\label{eq:transicion}
\end{equation}
con la fila de todo nodo sin salida (los pasajes) idénticamente nula. El
efecto conjunto de ambas operaciones se ve mejor escribiendo las
transiciones de un nodo-entidad $\varepsilon$ con hechos incidentes
$F(\varepsilon)$ y pasajes mencionantes $M(\varepsilon)$:
\begin{equation}
P(\varepsilon\to a) \;=\; \frac{\mu_a}{Z_\varepsilon}
\;\;(a\in F(\varepsilon)),
\qquad
P(\varepsilon\to c) \;=\; \frac{1}{Z_\varepsilon}
\;\;(c\in M(\varepsilon)),
\qquad
Z_\varepsilon \;=\; \sum_{a'\in F(\varepsilon)}\!\mu_{a'} \;+\;
|M(\varepsilon)|,
\label{eq:transiciones}
\end{equation}
mientras que cada aserción transita con probabilidad $1$ a su único pasaje
de procedencia, $P(a\to\pi(a))=1$. Dos observaciones. Primera: la
modulación es una \emph{competición local} ---reparte la masa de salida de
cada entidad entre \emph{sus} hechos y menciones, de modo que lo que
importa son los cocientes entre los $\mu$ de hechos hermanos, no su escala
global (los hechos aprobados, con $\mu=1$, ganan a sus hermanos de
$\mu<1$; ninguno baja de $1/4$, lo que preserva la exploración)---.
Segunda: la modulación actúa a la \emph{entrada} del hecho, no a su
salida: decide cuánto cuesta cruzar el peaje, no adónde va lo que lo
cruza. Esta modulación por consulta del propio grafo ---heredada del
recuperador en producción del que deriva \sys{}--- es el tercer punto de
inyección de la señal de la pregunta, junto a la siembra y el filtro.

\paragraph{Punto fijo, sumideros y coste.}
Con $P$ se itera el sistema de PageRank personalizado con redistribución
de la masa de los nodos sin salida $\mathcal{S}$ (\emph{dangling}; aquí,
todos los pasajes):
\begin{equation}
\mathbf{p}^{(t+1)} \;=\;
\lambda\Bigl(P^{\!\top}\mathbf{p}^{(t)}
 \;+\; \bigl(\textstyle\sum_{u\in\mathcal{S}} \mathbf{p}^{(t)}_u\bigr)\,\mathbf{v}\Bigr)
\;+\; (1-\lambda)\,\mathbf{v},
\qquad \mathbf{p}^{(0)}=\mathbf{v},
\label{eq:ppr}
\end{equation}
con factor de amortiguación $\lambda = 0{,}5$ (el valor publicado por
\hippo{}, adoptado por paridad; tabla~\ref{tab:hiper}), hasta convergencia
$\norm{\mathbf{p}^{(t+1)}-\mathbf{p}^{(t)}}_1 < 10^{-9}$ o un máximo de
$80$ iteraciones. La iteración conserva la masa
($\norm{\mathbf{p}^{(t)}}_1 = 1$: el término de redistribución devuelve al
vector de reinicio exactamente la masa que los sumideros absorberían) y es
una contracción de factor $\lambda$, de modo que el error decae como
$\lambda^{t}$ y con $\lambda=0{,}5$ bastan $\sim\!30$ iteraciones. Cada
iteración es un producto matriz dispersa--vector de coste $O(|E|)$
($165\,240$ aristas): el paseo completo cuesta milisegundos en CPU, y todo
el coste real por consulta está en las llamadas de la
tabla~\ref{tab:componentes}, no aquí.

\paragraph{Forma cerrada en el grafo por capas.}
La estructura por capas de $G^{\!\rightarrow}$ (sección~\ref{sec:grafo})
hace el punto fijo especialmente transparente. Como los nodos-entidad no
tienen aristas de entrada, su masa estacionaria es proporcional a su
siembra, $\mathbf{p}^{*}(\varepsilon) = \beta\,\mathbf{v}(\varepsilon)$
con $\beta = \lambda m^{*} + (1-\lambda)$, donde $m^{*}$ es la masa
estacionaria total de los sumideros; y como toda trayectoria muere en un
pasaje a lo sumo dos aristas después de nacer en una entidad, la
puntuación de un pasaje admite una descomposición \emph{exacta} en tres
canales:
\begin{equation}
\mathbf{p}^{*}(c) \;=\; \beta\Biggl[\;
\underbrace{\mathbf{v}(c)\vphantom{\sum_{\varepsilon}}}_{\text{siembra directa}}
\;+\;
\lambda \underbrace{\sum_{\varepsilon:\; c\,\in M(\varepsilon)}
\frac{\mathbf{v}(\varepsilon)}{Z_\varepsilon}}_{\text{menciones (1 arista)}}
\;+\;
\lambda^{2} \underbrace{\sum_{a:\;\pi(a)=c}\;
\sum_{\varepsilon\,\in\, a}
\frac{\mathbf{v}(\varepsilon)\,\mu_a}{Z_\varepsilon}}_{\text{hechos (2 aristas, modulado)}}
\;\Biggr].
\label{eq:tres-canales}
\end{equation}
La constante global $\beta$ no altera el orden. Tres lecturas exactas se
siguen de aquí. (i)~Las contribuciones decaen geométricamente con la
longitud del camino, $\lambda^{\ell}$, y en el grafo reificado un
\guillemotleft{}salto\guillemotright{} de la pregunta cuesta dos aristas: la evidencia mediada por
hechos entra con $\lambda^{2}=0{,}25$. (ii)~El canal de hechos de un pasaje
suma sobre \emph{todos} los participantes de cada hecho que lo afirma: un
hecho aprobado ($\mu_a=1$) canaliza hacia su pasaje la masa de todas sus
entidades a la vez (evidencia convergente). (iii)~\emph{El paseo no
encadena saltos en esta representación}: al no tener entradas los
nodos-entidad, ninguna trayectoria pasa de un participante a otro; el
encadenamiento del segundo salto lo realiza la siembra ---la compuerta
dura siembra a la entidad puente y la ecuación~\ref{eq:pasaje-propio} a su
pasaje propio--- y el paseo agrega los tres canales y desempata. La
implementación conserva el solucionador iterativo general
(ecuación~\ref{eq:ppr}) porque las variantes con ciclos no admiten esta
forma cerrada: el grafo no canónico de la ablación tiene aristas de
sinonimia bidireccionales entre entidades, y el grafo clínico de la
sección~\ref{sec:gesida} añade \textsc{anclada\_a} e \textsc{is\_a}
bidireccionales; en ellos $\lambda$ sí gobierna la profundidad efectiva
del paseo (longitud media de trayectoria $1/(1-\lambda)$ aristas entre
reinicios). La salida del sistema es en todos los casos la restricción de
$\mathbf{p}^{*}$ a los nodos-pasaje, ordenada de forma decreciente; se
devuelven los $k$ primeros títulos.

\emph{En el ejemplo conductor}: el paseo devuelve los cuatro pasajes oro
en los puestos 1--4 ---Edith Carlmar ($0{,}0070$), God's Gift to Women
($0{,}0070$), Aldri annet enn bråk ($0{,}0057$), Michael Curtiz
($0{,}0054$)--- con el primer distractor (otra película de Carlmar,
\emph{Altid ballade}) cinco veces por debajo ($0{,}0013$); la entidad
ruidosa \texttt{can:norway} resulta inocua, porque su masa se reparte
entre las muchísimas salidas de un nodo concentrador
(ec.~\ref{eq:transiciones}) y no tiene pasaje propio que sembrar. Cadena
completa de cuatro en el top-4 ($\mathrm{FC@}5=1$) en una pregunta en la
que los sistemas publicados no pasan del $12{,}3\%$ de media. Vista
entera, la cadena de relevos es nítida: el enrutador eligió la estrategia,
el enlazador ancló las dos películas, la afinidad puntuó los hechos, el
filtro descubrió a los dos directores (corrigiendo el sesgo del coseno
hacia una rama), la siembra convirtió ese descubrimiento en masa sobre las
dos biografías, y el paseo agregó los canales y ordenó.

\begin{table}[t]
\centering\small
\caption{Hiperparámetros de \sys{} y procedencia de cada valor. Ningún valor
se ajustó sobre el banco de pruebas: los diales del PPR se adoptaron de los
valores publicados de \hippo{} (decisión pre-registrada) y el resto se fijó
sobre los conjuntos de sondeo y calibración.}
\label{tab:hiper}
\footnotesize\setlength{\tabcolsep}{3.5pt}
\begin{tabular}{llll}
\toprule
componente & parámetro & valor & procedencia\\
\midrule
diccionario & umbral candidato (ficha) & $0{,}75$ & sondeo\\
            & umbral variante de nombre & $0{,}80$ & sondeo\\
            & fusión automática (nombre / ficha) & $0{,}90$ / $0{,}70$ & sondeo\\
            & fusión automática sin LLM (ficha) & $0{,}92$ & sondeo\\
            & vecinos por canal / lote LLM & $15$ / $8$ & sondeo\\
enlazador   & umbral de aceptación por ficha & $0{,}80$ & sondeo\\
filtro      & candidatos por coseno / por frontera & $40$ / $15\cdot|L(q)|$ ($\le 20$) & sondeo\\
            & selección máxima & $4$ hechos & sondeo\\
siembra     & peso denso de pasaje $w_{\text{pas}}$ & $0{,}05$ & \hippo{} (pre-registrado)\\
            & pasajes sembrados & todos & \hippo{} (pre-registrado)\\
            & ancla débil $\epsilon$ / ancla del enlazador $w_{\text{enl}}$ & $0{,}05$ / $0{,}5$ & sondeo\\
PPR         & amortiguación $\lambda$ & $0{,}5$ & \hippo{} (pre-registrado)\\
            & rango de modulación $\mu$ & $[0{,}25,\,1]$ & producción (heredado)\\
            & iteraciones / tolerancia & $\le 80$ / $10^{-9}$ & fijo\\
modelos     & lenguaje / incrustación & \texttt{gpt-4o-mini} / \texttt{te3-small} & paridad con \hippo{}\\
\bottomrule
\end{tabular}
\end{table}

\subsubsection{Ejemplo numérico: de la matriz de transición al orden final}
\label{sec:ppr-ejemplo}

Para hacer tangible el cálculo completo, reducimos la traza real de la
sección~\ref{sec:traza} a un grafo de nueve nodos que conserva su
estructura exacta (figura~\ref{fig:ppr-toy}): tres entidades
---$E_1=$ \texttt{can:atomised}, $E_2=$ \texttt{can:oskar roehler} y el
distractor $E_3=$ \texttt{can:allison anders}---, tres aserciones
---$a_1=$ \guillemotleft{}Atomised was directed by Oskar Roehler\guillemotright{} (aprobada por el
filtro: $\tilde\sigma=1$, $\mu_{a_1}=1$), $a_2=$ \guillemotleft{}Allison Anders is an
independent film director\guillemotright{} ($\sigma=0{,}503$, luego
$\mu_{a_2}=0{,}25+0{,}75\cdot 0{,}503=0{,}627$) y $a_3=$ \guillemotleft{}Atomised stars
Franka Potente\guillemotright{} ($\sigma=0{,}463$, $\mu_{a_3}=0{,}597$)--- y tres pasajes
$C_1$--$C_3$ (los artículos de la película y de los dos directores). Las
aristas de búsqueda son: roles invertidos $E_1\to a_1$, $E_1\to a_3$,
$E_2\to a_1$, $E_3\to a_2$; procedencias $a_1\to C_1$, $a_3\to C_1$,
$a_2\to C_3$; y menciones $E_1\to C_1$, $E_2\to C_1$, $E_2\to C_2$,
$E_3\to C_3$.

\paragraph{Construcción de $P$, fila a fila.}
Para $E_1$, la fila de $\tilde A$ (ec.~\ref{eq:modulacion}) vale $1$ en
$a_1$ (hecho aprobado), $0{,}597$ en $a_3$ y $1$ en $C_1$; su grado de
salida modulado es $Z_{E_1}=2{,}597$ y la fila de $P$
(ec.~\ref{eq:transiciones}) queda $P(E_1\!\to\!a_1)=0{,}385$,
$P(E_1\!\to\!a_3)=0{,}230$, $P(E_1\!\to\!C_1)=0{,}385$: el hecho aprobado
gana la competición local a su hermano no aprobado. Análogamente
$P(E_2\!\to\!a_1)=P(E_2\!\to\!C_1)=P(E_2\!\to\!C_2)=1/3$
(sus tres salidas tienen peso $1$) y, para el distractor,
$P(E_3\!\to\!a_2)=0{,}627/1{,}627=0{,}385$,
$P(E_3\!\to\!C_3)=0{,}615$. Las filas de las aserciones son
deterministas ($a_1\to C_1$, $a_3\to C_1$, $a_2\to C_3$, con probabilidad
$1$) y las de los pasajes, nulas (sumideros).

\paragraph{Siembra e iteración.}
El vector de personalización toma los valores reales de la traza,
normalizados: $\mathbf{v} = (E_1\!:0{,}328;\; E_2\!:0{,}313;\;
C_1\!:0{,}180;\; C_2\!:0{,}166;\; C_3\!:0{,}013)$; ni $E_3$ ni las
aserciones siembran. Iterando la ecuación~\eqref{eq:ppr} con
$\lambda=0{,}5$, la masa de $C_1$ evoluciona
$0{,}180 \to 0{,}238 \to 0{,}282 \to 0{,}263 \to \dots \to 0{,}266$
(convergencia a $10^{-9}$ en unas 20 iteraciones), con masa estacionaria
de sumideros $m^{*}=0{,}432$ y constante $\beta = 0{,}716$. La
tabla~\ref{tab:ppr-canales} muestra la descomposición exacta de la
ecuación~\eqref{eq:tres-canales}, que reproduce la iteración al cuarto
decimal.

\begin{table}[t]
\centering\footnotesize
\caption{Descomposición en tres canales (ec.~\ref{eq:tres-canales}) del
ejemplo numérico, verificada contra la iteración de potencias
(ec.~\ref{eq:ppr}). $\lambda=0{,}5$, $\beta=0{,}716$.}
\label{tab:ppr-canales}
\begin{tabular}{lccccc}
\toprule
pasaje & directo $\mathbf{v}(c)$ & $\lambda\cdot$menciones &
$\lambda^{2}\cdot$hechos & $\mathbf{p}^{*}(c)=\beta\cdot\text{suma}$ & rango\\
\midrule
$C_1$ \emph{Atomised (film)} & $0{,}1798$ & $0{,}1153$ & $0{,}0765$ & $0{,}2660$ & 1\\
$C_2$ \emph{Oskar Roehler}   & $0{,}1657$ & $0{,}0521$ & $0$        & $0{,}1559$ & 2\\
$C_3$ \emph{Allison Anders}  & $0{,}0134$ & $0$        & $0$        & $0{,}0096$ & 3\\
\bottomrule
\end{tabular}
\end{table}

\begin{figure}[t]
\centering
\begin{tikzpicture}[
  font=\scriptsize,
  ent/.style={draw, ellipse, align=center, inner sep=1.5pt},
  aser/.style={draw, rectangle, rounded corners=1pt, align=center, inner sep=2.5pt, fill=gray!12},
  chk/.style={draw, rectangle, align=center, inner sep=2.5pt},
  fl/.style={-{Stealth[length=1.8mm]}, semithick},
  et/.style={font=\tiny, fill=white, inner sep=0.5pt},
  peso/.style={font=\tiny, text=black!70}]
  % entidades
  \node[ent, label={[peso]left:{$\mathbf{v}=0{,}328$}}] (E1) at (0,0) {$E_1$};
  \node[ent, label={[peso]left:{$\mathbf{v}=0{,}313$}}] (E2) at (0,-1.5) {$E_2$};
  \node[ent, label={[peso]left:{$\mathbf{v}=0$}}] (E3) at (0,-3.0) {$E_3$};
  % aserciones
  \node[aser, label={[peso]above:{$\mu=0{,}597$}}] (A3) at (3.6,0.8) {$a_3$};
  \node[aser, label={[peso]below:{$\mu=1$ (aprobado)}}] (A1) at (3.6,-0.8) {$a_1$};
  \node[aser, label={[peso]below:{$\mu=0{,}627$}}] (A2) at (3.6,-3.0) {$a_2$};
  % pasajes
  \node[chk, label={[peso]right:{$\mathbf{p}^{*}=0{,}266$ (1.\textsuperscript{o})}}] (C1) at (7.4,0) {$C_1$};
  \node[chk, label={[peso]right:{$\mathbf{p}^{*}=0{,}156$ (2.\textsuperscript{o})}}] (C2) at (7.4,-1.5) {$C_2$};
  \node[chk, label={[peso]right:{$\mathbf{p}^{*}=0{,}010$ (3.\textsuperscript{o})}}] (C3) at (7.4,-3.0) {$C_3$};
  % aristas E->A
  \draw[fl] (E1) -- node[et, pos=0.45]{$0{,}230$} (A3);
  \draw[fl] (E1) -- node[et, pos=0.45]{$0{,}385$} (A1);
  \draw[fl] (E2) -- node[et, pos=0.4]{$0{,}333$} (A1);
  \draw[fl] (E3) -- node[et, pos=0.45]{$0{,}385$} (A2);
  % aristas A->C
  \draw[fl] (A3) -- node[et, pos=0.5]{$1$} (C1);
  \draw[fl] (A1) -- node[et, pos=0.5]{$1$} (C1);
  \draw[fl] (A2) -- node[et, pos=0.5]{$1$} (C3);
  % menciones E->C (curvadas por encima/debajo)
  \draw[fl, gray] (E1) to[out=25,in=155] node[et, pos=0.25]{$0{,}385$} (C1);
  \draw[fl, gray] (E2) to[out=-15,in=-160] node[et, pos=0.3]{$0{,}333$} (C1);
  \draw[fl, gray] (E2) to[out=-35,in=180] node[et, pos=0.35]{$0{,}333$} (C2);
  \draw[fl, gray] (E3) to[out=-25,in=-155] node[et, pos=0.3]{$0{,}615$} (C3);
\end{tikzpicture}
\caption{Grafo de juguete del ejemplo numérico (nueve nodos, estructura
por capas real): entidades a la izquierda con su siembra $\mathbf{v}$,
aserciones en el centro con su modulación $\mu$, pasajes a la derecha con
su masa estacionaria. Las etiquetas de las aristas son las probabilidades
de transición de $P$ (ec.~\ref{eq:transiciones}); las grises son
menciones directas entidad$\to$pasaje.}
\label{fig:ppr-toy}
\end{figure}

\paragraph{Lectura del ejemplo.}
Tres lecciones compactas. Primera: $C_1$ domina porque es el único pasaje
que suma los tres canales, y su canal de hechos recibe a la vez de $E_1$ y
de $E_2$ a través del hecho aprobado (la evidencia convergente del punto
(ii) anterior). Segunda: $C_2$ ---el pasaje del segundo salto--- vive de
la siembra: su canal directo procede del término de pasaje propio
(ec.~\ref{eq:pasaje-propio}) y su canal de menciones, de que $E_2$ fue
sembrada por la compuerta dura; si el filtro no hubiera aprobado el hecho,
$\mathbf{v}(E_2)=0$ y $C_2$ quedaría reducido a su (mediocre) masa densa,
por debajo de los distractores ---la ilustración mínima de que el
encadenamiento pertenece a la siembra---. Tercera: $E_3$, no sembrada,
tiene $\mathbf{p}^{*}=0$ exactamente (las entidades carecen de entradas),
igual que su hecho $a_2$: los distractores no reciben masa por el mero
hecho de existir en el grafo, y $C_3$ solo conserva su pequeña presa
densa.

% =========================================================================
\subsection{Ejemplo trazado de extremo a extremo}
\label{sec:traza}

Ilustramos el flujo con una pregunta composicional real del banco de pruebas,
instrumentando la ejecución del sistema exactamente como se evaluó (los
valores numéricos son los observados):

\begin{quote}
\emph{\guillemotleft{}Who is the mother of the director of film Atomised (Film)?\guillemotright{}}\\
Evidencia oro: pasajes \emph{Atomised (film)} y \emph{Oskar Roehler};
respuesta: Gisela Elsner.
\end{quote}

\begin{enumerate}[leftmargin=1.6em, itemsep=2pt]
\item \textbf{Enrutador} (\S\ref{sec:router}): sin marcas de comparación ni
  parentesco de segundo grado $\Rightarrow$ tipo \emph{composicional}
  $\Rightarrow$ estrategia con filtro selector.
\item \textbf{Enlazador} (\S\ref{sec:linker}): tramo capitalizado
  \guillemotleft{}Atomised\guillemotright{} $\to$ plegado $\to$ coincidencia exacta con el alias
  \guillemotleft{}Atomised (film)\guillemotright{} de la entrada \texttt{can:atomised}. Por tanto
  $L(q) = \{\texttt{can:atomised}\}$.
\item \textbf{Afinidad} (ec.~\ref{eq:afinidad}): el hecho puente
  \emph{\guillemotleft{}Atomised was directed by Oskar Roehler\guillemotright{}} encabeza el coseno con
  $\sigma = 0{,}503$, pero \emph{empatado} con un distractor
  (\emph{\guillemotleft{}Allison Anders is an independent film director\guillemotright{}}, $\sigma=0{,}503$)
  y seguido de hechos de otras películas: la semejanza superficial no
  distingue utilidad.
\item \textbf{Candidatos y filtro} (ec.~\ref{eq:candidatos},
  \S\ref{sec:filtro}): $\mathcal{K}(q)$ = Top-40 por coseno $\cup$ los $15$ hechos de mayor
  $\sigma$ de los $16$ incidentes en \texttt{can:atomised} (su frontera),
  de los que $8$ no estaban en el Top-40 (rangos por coseno $65$ a
  $12\,045$): $48$ candidatos. El hecho puente ocupa el rango $1$ por
  coseno, así que aquí la frontera no decide. El filtro selecciona
  exactamente un hecho, $S(q) = \{a_1\}$ con
  $\tau_{a_1} = $ \guillemotleft{}Atomised was directed by Oskar Roehler\guillemotright{}:
  conserva el salto intermedio y descarta los directores de otras películas.
\item \textbf{Siembra} (\S\ref{sec:siembra}): la compuerta dura asigna
  $\mathbf{v}=1$ a \texttt{can:atomised} y \texttt{can:oskar roehler}
  (entidades del hecho aprobado); el ancla del enlazador y el ancla débil
  elevan \texttt{can:atomised} a $1{,}05$ (el máximo con $0{,}5$ no
  actúa; el ancla débil aporta $+0{,}05$, y otro $+0{,}05$ va a un
  distractor, \texttt{can:annabelle}); la presa densa reparte $112{,}6$
  unidades entre los $6\,119$ pasajes ($0{,}050$ a \emph{Atomised (film)},
  $0{,}030$ a \emph{Oskar Roehler}). El término de pasaje propio
  (ec.~\ref{eq:pasaje-propio}) inyecta masa directa en
  \emph{Atomised (film)} ($0{,}5\times1{,}05=0{,}525$; $0{,}575$ sumada
  la presa densa) y en \emph{Oskar Roehler} ($0{,}500$; $0{,}530$).
  Normalizado por $\norm{\mathbf{v}}_1=115{,}8$: $0{,}0091$, $0{,}0086$,
  $0{,}0050$ y $0{,}0046$ respectivamente. La figura~\ref{fig:traza}
  recoge la trazabilidad completa, etapa a etapa, y la
  figura~\ref{fig:traza-subgrafo} el subgrafo efectivo. Obsérvese el punto crítico: la
  biografía del director tiene similitud densa mediocre con la pregunta
  (la pregunta habla de una madre y de una película; la palabra \guillemotleft{}Roehler\guillemotright{} no
  aparece) --- es la cadena
  hecho aprobado $\to$ entidad $\to$ pasaje propio la que la convierte en
  semilla.
\item \textbf{PPR modulado} (ec.~\ref{eq:modulacion}--\ref{eq:ppr}): el hecho
  aprobado transita con $\mu=1$ y el paseo concentra la masa estacionaria en
  los dos pasajes oro, con un margen de un orden de magnitud sobre el primer
  distractor: $\mathbf{p} = 0{,}0084$ (\emph{Atomised (film)}),
  $0{,}0063$ (\emph{Oskar Roehler}), $0{,}0004$ (\emph{Allison Anders},
  primer no-oro). Cadena completa en los dos primeros puestos
  ($\mathrm{FC@}2 = 1$).
\end{enumerate}

\begin{figure}[p]
\centering
\begin{tikzpicture}[
  font=\scriptsize,
  etapa/.style={draw, rounded corners=2pt, fill=gray!8, align=center, inner sep=3pt, text width=22mm, minimum height=9mm},
  prod/.style={draw, dashed, rounded corners=2pt, align=left, inner sep=3pt},
  fl/.style={-{Stealth[length=2mm]}, semithick},
  gr/.style={-{Stealth[length=2mm]}, semithick, gray}]
  % ---- pregunta
  \node[prod, text width=146mm, align=center] (q) at (0,0)
    {\textbf{Pregunta} \guillemotleft{}Who is the mother of the director of film Atomised (Film)?\guillemotright{} \quad oro: \emph{Atomised (film)}, \emph{Oskar Roehler}};
  % ---- enrutador
  \node[etapa] (router) at (-6.2,-1.5) {Enrutador léxico\\ \S\ref{sec:router}};
  \node[prod, text width=108mm, anchor=west] (p0) at (-4.5,-1.5)
    {sin marcas de comparación ni de parentesco de 2.\textsuperscript{o} grado $\Rightarrow$ \textbf{composicional} $\Rightarrow$ estrategia de filtro ($J=0$: un solo vector de consulta). El texto de la pregunta sigue intacto.};
  % ---- dos ramas independientes
  \node[etapa] (linker) at (-6.2,-4.3) {Enlazador\\ (rama simbólica)\\ \S\ref{sec:linker}};
  \node[prod, text width=39mm, anchor=west] (p1) at (-4.5,-4.3)
    {tramo por mayúsculas: \guillemotleft{}Atomised\guillemotright{} (\guillemotleft{}(Film)\guillemotright{} cae con el vocabulario excluido)\\
     plegado \guillemotleft{}atomised\guillemotright{} $\to$ alias exacto (escalón 1; sin respaldo vectorial)\\
     $L(q)=\{\texttt{can:atomised}\}$};
  \node[etapa] (afin) at (1.45,-4.3) {Afinidad\\ (rama semántica)\\ \S\ref{sec:afinidad}};
  \node[prod, text width=42mm, anchor=west] (p2) at (2.95,-4.3)
    {$v_0=e(q)$; $\sigma_a=\langle v_0,e(\tau_a)\rangle$ en $47\,999$ hechos (mín $-0{,}100$, máx $0{,}503$)\\
     1.\textsuperscript{o} \guillemotleft{}Atomised was directed by Oskar Roehler\guillemotright{} $0{,}503$\\
     2.\textsuperscript{o} \guillemotleft{}Allison Anders is an independent film director\guillemotright{} $0{,}503$\\
     3.\textsuperscript{o} \guillemotleft{}Allison Anders is an American independent film director\guillemotright{} $0{,}486$};
  % ---- filtro
  \node[etapa] (filtro) at (-6.2,-7.6) {Filtro selector\\ (LLM) \S\ref{sec:filtro}};
  \node[prod, text width=108mm, anchor=west] (p3) at (-4.5,-7.6)
    {$\mathcal{K}(q)$ = Top-40 por $\sigma$ $\cup$ frontera de \texttt{can:atomised}: $16$ hechos incidentes $\to$ $15$ de mayor $\sigma$, de los que $8$ no estaban en el Top-40 (rangos $65$--$12\,045$) $\Rightarrow$ $48$ candidatos numerados\\
     juez (\texttt{gpt-4o-mini}, una llamada, $\le 4$ hechos): $S(q)=\{a_1\}$ = \guillemotleft{}Atomised was directed by Oskar Roehler\guillemotright{}; descarta a Allison Anders y a los demás directores\\
     $\tilde\sigma_{a_1}=1$; para el resto, $\tilde\sigma=\sigma$};
  % ---- siembra
  \node[etapa] (siembra) at (-6.2,-11.0) {Siembra $\mathbf{v}$\\ \S\ref{sec:siembra}};
  \node[prod, text width=108mm, anchor=west] (p4) at (-4.5,-11.0)
    {(a) compuerta dura sobre $S(q)$: \texttt{can:atomised} $1$, \texttt{can:oskar roehler} $1$ (el director entra aquí por primera vez)\\
     (b) presa densa: $6\,119$ pasajes $\times\,0{,}05\times$ coseno normalizado con $v_0$ (\emph{Atomised (film)} $0{,}050$; \emph{Oskar Roehler} $0{,}030$); suma $112{,}6$\\
     (c) anclas débiles: $+0{,}05$ a \texttt{can:atomised} y a \texttt{can:annabelle} (distractor)\\
     (d) $L(q)$: $\max(\cdot,0{,}5)$ no actúa; pasaje propio: \emph{Atomised (film)} $\mathrel{+}=0{,}5\times1{,}05$ $\to 0{,}575$; \emph{Oskar Roehler} $\mathrel{+}=0{,}5\times1{,}00$ $\to 0{,}530$\\
     $\norm{\mathbf{v}}_1=115{,}8$ $\Rightarrow$ \texttt{can:atomised} $0{,}0091$, \texttt{can:oskar roehler} $0{,}0086$, \emph{Atomised (film)} $0{,}0050$, \emph{Oskar Roehler} $0{,}0046$};
  % ---- paseo
  \node[etapa] (ppr) at (-6.2,-14.2) {PPR modulado\\ \S\ref{sec:ppr}};
  \node[prod, text width=108mm, anchor=west] (p5) at (-4.5,-14.2)
    {$\mu_{a_1}=1$ (peaje abierto); los hechos no aprobados quedan en $\mu=0{,}25+0{,}75\,\sigma$; $\lambda=0{,}5$, convergencia en $\sim 30$ iteraciones\\
     $\mathbf{p}^{*}$: \emph{Atomised (film)} $0{,}0084$ (1.\textsuperscript{o}, oro); \emph{Oskar Roehler} $0{,}0063$ (2.\textsuperscript{o}, oro); \emph{Allison Anders} $0{,}0004$ (3.\textsuperscript{o}) $\Rightarrow$ $\mathrm{FC@}2=1$};
  % ---- flechas de dependencia (columna de etapas)
  \draw[fl] (q.south -| router.north) -- (router.north);
  \draw[fl] (router.south) -- (linker.north);
  \draw[fl] (router.south) to[out=-90,in=90] (afin.north);
  \draw[fl] (linker.south) -- (filtro.north);
  \draw[fl] (afin.south) to[out=-90,in=90] (filtro.north);
  \draw[fl] (filtro.south) -- (siembra.north);
  \draw[fl] (siembra.south) -- (ppr.north);
  % ---- productos que llegan directamente a la siembra (por los margenes)
  \draw[gr] (linker.west) -- ++(-0.4,0) |- (siembra.west);
  \node[font=\tiny, gray, anchor=west] at (-7.75,-9.4) {$L(q)$};
  \draw[gr] (p2.east) -- ++(0.3,0) |- (p4.east);
  \node[font=\tiny, gray, anchor=east] at (7.85,-9.4) {$v_0$, $\sigma$};
\end{tikzpicture}
\caption{Trazabilidad completa de la pregunta de la sección~\ref{sec:traza},
etapa a etapa y con los valores observados. La estructura de dependencias
es la de la tabla~\ref{tab:dependencias}: enlazador y afinidad leen solo la
pregunta y son independientes; el filtro es la primera etapa que necesita
las dos (Top-40 por $\sigma$ y frontera de $L(q)$); la siembra recibe los
tres productos (las flechas grises por los márgenes llevan $L(q)$ y
$v_0,\sigma$ directamente a ella), y el paseo, la siembra y $\tilde\sigma$.}
\label{fig:traza}
\end{figure}

\begin{figure}[t]
\centering
\begin{tikzpicture}[
  font=\footnotesize,
  ent/.style={draw, ellipse, align=center, inner sep=2pt},
  aser/.style={draw, rectangle, rounded corners=1pt, align=center, inner sep=3pt, fill=gray!12},
  chk/.style={draw, rectangle, align=center, inner sep=3pt},
  fl/.style={-{Stealth[length=2mm]}, semithick},
  peso/.style={font=\scriptsize, text=black!70},
  node distance=5mm and 10mm]
  \node[ent, label={[peso]above:{$\mathbf{v}=1{,}05$ \ (filtro + enlace + ancla)}}] (e1) {can:atomised};
  \node[aser, right=of e1, label={[peso]above:{$\mu=1$ (hecho aprobado)}}]
       (a) {$a_1$: \guillemotleft{}Atomised was directed\\ by Oskar Roehler.\guillemotright{}};
  \node[ent, right=of a, label={[peso]above:{$\mathbf{v}=1{,}00$ \ (filtro)}}] (e2) {can:oskar roehler};
  \node[chk, below=8mm of e1, label={[peso]below:{$\mathbf{v}=0{,}575$ \ (pasaje propio + densa)}}]
       (c1) {\emph{Atomised (film)}\\ oro 1 --- rango 1};
  \node[chk, below=8mm of e2, label={[peso]below:{$\mathbf{v}=0{,}530$ \ (pasaje propio + densa)}}]
       (c2) {\emph{Oskar Roehler}\\ oro 2 --- rango 2};
  \draw[fl] (e1) -- (a);
  \draw[fl] (e2) -- (a);
  \draw[fl] (a) to[out=-150,in=30] (c1);
  \draw[fl] (e1) -- (c1);
  \draw[fl] (e2) -- (c2);
\end{tikzpicture}
\caption{Subgrafo efectivo de la traza de la sección~\ref{sec:traza} con las
masas de siembra reales (sin normalizar), en la orientación real del grafo
de búsqueda: los dos participantes apuntan al hecho aprobado y este desagua
en su pasaje de procedencia. El término de pasaje propio
(ec.~\ref{eq:pasaje-propio}) siembra directamente las dos biografías, y el
paseo agrega los canales y confirma el orden. El segundo salto de la
pregunta (la madre) lo responde el lector sobre el pasaje
\emph{Oskar Roehler}, que contiene
\guillemotleft{}Oskar Roehler is the son of Gisela Elsner\guillemotright{}.}
\label{fig:traza-subgrafo}
\end{figure}

% =========================================================================
\section{Caso de uso clínico: el asistente sobre las guías GeSIDA de VIH}
\label{sec:gesida}

El sistema evaluado en este trabajo no nació como artefacto de laboratorio:
es la generalización del recuperador de un asistente clínico en desarrollo
sobre las guías GeSIDA de VIH. Describimos aquí ese prototipo con el mismo
nivel de detalle que el método general, por tres motivos: es la
instanciación del patrón en un dominio con terminología externa
(SNOMED~CT), añade las capas de seguridad que un dominio clínico exige, y
la comparación entre ambas instancias (tabla~\ref{tab:gesida-wiki}) aclara
qué partes del método son generales y cuáles dependen del dominio.

\subsection{Contexto y requisitos}

El asistente trabaja sobre ocho guías GeSIDA en PDF ---tratamiento
antirretroviral en adultos (2022), adherencia (2020), profilaxis
postexposición, tuberculosis en VIH, VIH y embarazo, alteraciones
neurocognitivas, riesgo cardiovascular y metabólico, y medicina
preventiva/vacunas--- y se valida contra un banco propio de $505$ preguntas
clínicas en cuatro niveles de complejidad ($177/159/96/73$), de las cuales
$176$ exigen conectar dos o más guías y $73$ requieren una pausa de
supervisión médica. El principio rector de diseño es \emph{razonar la
estrategia, no el contenido}: el sistema decide qué recuperar, cómo
combinar evidencia de varias guías y cuándo detenerse, pero toda afirmación
clínica queda anclada a una cita verificable (guía, sección, página y grado
de evidencia) y la aritmética de umbrales la ejecutan mecanismos
deterministas. Tres garantías se aplican en cada consulta: trazabilidad de
citas, verificación de fidelidad de la respuesta y abstención honesta
cuando la guía no cubre la pregunta.

\begin{figure}[t]
\centering
\begin{tikzpicture}[
  font=\scriptsize,
  caja/.style={draw, rounded corners=2pt, align=center, minimum height=7mm, inner sep=2.5pt, fill=gray!8},
  dato/.style={draw, dashed, rounded corners=2pt, align=center, inner sep=2.5pt},
  fl/.style={-{Stealth[length=2mm]}, semithick},
  node distance=2.5mm and 3.5mm]
  % capa de conocimiento
  \node[dato] (pdf) {8 guías\\ GeSIDA (PDF)};
  \node[caja, right=of pdf] (chk) {troceado\\ estructural\\ 2\,694 fragm.};
  \node[caja, right=of chk] (ner) {NER + reificación\\ clínica (LLM)\\ 8\,444 aserciones};
  \node[caja, right=of ner] (rec) {recomendaciones\\ deóntica + grado\\ 1\,332};
  \node[caja, right=of rec] (snomed) {anclaje SNOMED\\ (UMLS local + juez\\ + curación)};
  \node[caja, right=of snomed] (kg) {grafo clínico\\ 22\,807 nodos\\ + incrustaciones};
  \draw[fl] (pdf) -- (chk);
  \draw[fl] (chk) -- (ner);
  \draw[fl] (ner) -- (rec);
  \draw[fl] (rec) -- (snomed);
  \draw[fl] (snomed) -- (kg);
  % capa agentica
  \node[dato, below=15mm of pdf] (q) {pregunta\\ clínica};
  \node[caja, right=of q] (intake) {acogida:\\ seudonim.\\ + reescritura};
  \node[caja, right=of intake] (router) {enrutador\\ (nivel, guías,\\ estrategia)};
  \node[caja, right=of router] (recu) {recuperación\\ grafo / densa\\ (+ ECL clase)};
  \node[caja, right=of recu] (fus) {fusión RRF +\\ reordenador CE\\ diverso};
  \node[caja, right=of fus] (sint) {síntesis\\ citada};
  \node[caja, right=of sint] (verif) {verificación:\\ juez + lentes\\ deónt./num.};
  \node[dato, below=2.5mm of verif] (resp) {respuesta\\ o abstención\\ (o pausa HITL)};
  \draw[fl] (q) -- (intake);
  \draw[fl] (intake) -- (router);
  \draw[fl] (router) -- (recu);
  \draw[fl] (recu) -- (fus);
  \draw[fl] (fus) -- (sint);
  \draw[fl] (sint) -- (verif);
  \draw[fl] (verif) -- (resp);
  \draw[fl, gray] (verif.north) to[out=160,in=20]
    node[pos=0.5, fill=white, inner sep=1pt, font=\tiny]{recuperación dirigida} (recu.north);
  \draw[fl, gray] (kg.south) to[out=-90,in=90] (recu.north east);
  \begin{scope}[on background layer]
    \node[fit=(pdf)(kg), draw=gray!60, rounded corners=3pt, inner sep=3pt,
          label={[gray, font=\tiny]above left:{capa de conocimiento (indexación)}}] {};
    \node[fit=(q)(verif)(resp), draw=gray!60, rounded corners=3pt, inner sep=3pt,
          label={[gray, font=\tiny]above left:{capa agéntica (consulta)}}] {};
  \end{scope}
\end{tikzpicture}
\caption{Arquitectura del asistente clínico GeSIDA. La capa de conocimiento
instancia el patrón de la figura~\ref{fig:pipeline} con terminología externa
(SNOMED~CT/UMLS) y una clase de nodo adicional (la recomendación normativa);
la capa agéntica añade fusión con recuperación densa, reordenación,
verificación determinista y supervisión humana.}
\label{fig:gesida}
\end{figure}

\subsection{Capa de conocimiento (indexación clínica)}
\label{sec:gesida-index}

\paragraph{Troceado estructural.}
Las guías se fragmentan respetando su estructura, no por un número ciego de
caracteres: los encabezados se detectan por numeración de sección
(\guillemotleft{}6.2.1\guillemotright{}) o por tamaño tipográfico, y las tablas ---de dosis y de
interacciones--- se detectan heurísticamente y se emiten como fragmento
propio \emph{sin partir}, porque una pauta posológica troceada es
inservible y peligrosa. Resultan $2\,694$ fragmentos (objetivo
$1\,100$ caracteres, solape $150$) con identificador canónico
\guillemotleft{}guía::página::índice\guillemotright{} (p.\,ej.\
\texttt{TAR\_ADULTOS\_2022::p6::10}), que es el anclaje de toda cita. Cada
fragmento lleva metadatos extraídos por léxicos y reglas: fármacos
antirretrovirales mencionados, condiciones, \emph{umbrales numéricos}
(filtrado glomerular, CD4, carga viral, HLA-B*5701) y grado de evidencia
GeSIDA presente en el texto.

\paragraph{Extracción clínica en dos pasos.}
Un modelo de lenguaje (DeepSeek-V3, temperatura $0$, con el fragmento
anterior como contexto) realiza primero reconocimiento de entidades con
seis clases clínicas (fármaco, enfermedad, condición clínica, población,
intervención, resultado clínico) y después la reificación: cada hecho se
promueve a \emph{aserción clínica}, la especialización de la
ecuación~\eqref{eq:asercion} para este dominio. Sobre la aserción genérica,
la clínica fija la semántica de los roles ---intervención (sujeto),
resultado (objeto) y \emph{población diana} como tercer rol de primera
clase--- y añade campos propios: tipo de aserción (recomendación
terapéutica, evidencia clínica, advertencia, diagnóstico), nivel de
evidencia, marca de negación y una ontología semicerrada de 13 relaciones
base (\texttt{TREATS}, \texttt{CONTRAINDICATES}, \texttt{PREVENTS},
\texttt{INCREASES\_RISK},~\dots). El corpus produce $8\,444$ aserciones y
$7\,737$ entidades. Como en el caso general, la frase verbalizada es lo que
se incrusta; la etiqueta de relación es metadato.

\paragraph{Recomendaciones normativas.}
Un extractor \emph{por reglas} (sin modelo de lenguaje) identifica las
frases normativas: las que llevan grado GeSIDA (\guillemotleft{}A-I\guillemotright{}\dots\guillemotleft{}C-III\guillemotright{})
o un verbo deóntico. Cada una se normaliza a una acción deóntica
---recomendado, considerar, no recomendado, generalmente no recomendado,
sin determinar--- y se materializa como nodo \emph{Recomendación} con
guía, sección, página, grado y fecha de vigencia: $1\,332$ nodos. Esta
clase de nodo no existe en el caso enciclopédico; codifica la unidad de
respuesta del dominio (la polaridad de una recomendación jamás puede
invertirse, y su vigencia es un guardarraíl de seguridad).

\paragraph{Anclaje terminológico SNOMED~CT/UMLS.}
El papel que en 2Wiki desempeña el diccionario derivado del corpus
(sección~\ref{sec:diccionario}) lo desempeña aquí una terminología externa:
una instalación local de UMLS~\cite{umls} (versión 2026AA; fuentes SCTSPA,
SNOMEDCT\_US, MSHSPA y RXNORM; 2,3 millones de términos en SQLite con
índice de trigramas), con el identificador SNOMED~CT (SCTID) como identidad
primaria y la jerarquía IS-A construida desde MRHIER ($386\,000$ pares). El
enlazador resuelve cada entidad extraída en cascada estricta:
\begin{enumerate}[leftmargin=1.6em, itemsep=1pt]
\item \textbf{Exacto normalizado} (plegado Unicode, sin diacríticos):
  puntuación $1{,}0$.
\item \textbf{Difuso léxico}: búsqueda FTS5 de trigramas (frase entera;
  si no, unión de palabras largas) puntuada con el coeficiente de Dice
  sobre trigramas, umbral $0{,}55$. Este es el escalón que en otros
  sistemas ocuparía SymSpell; aquí vive en la base terminológica, no en el
  código.
\item \textbf{Filtro de grupo semántico por rol}: un fármaco solo puede
  anclar a conceptos químicos; una condición, a trastornos u organismos
  (los patógenos que el NER etiqueta como enfermedad).
\item \textbf{Umbral de anclaje} $0{,}85$ y regla de ambigüedad: dos
  candidatos exactos de grupos distintos sin rol que desambigüe
  $\Rightarrow$ ninguno.
\item \textbf{Juez LLM} para lo no resuelto con candidatos: ve la frase
  real de la guía y los candidatos con su grupo y puntuación, bajo la regla
  \guillemotleft{}ante la duda, ninguno: un anclaje erróneo es peor que ninguno\guillemotright{}; solo
  se aceptan veredictos de confianza alta o media.
\item \textbf{Curación humana}: bandeja de revisión en la interfaz con
  prioridad juez $>$ manual $>$ difuso $>$ exacto (el pilotaje mostró que
  los anclajes del juez son los que más revisión necesitan); toda decisión
  queda auditada.
\end{enumerate}
Cobertura actual: $2\,326$ de $7\,737$ entidades con SCTID ($30\%$;
$1\,488$ exactas, $1\,060$ por juez, $249$ difusas). Obsérvese la simetría
con la sección~\ref{sec:diccionario}: candidatos baratos $\to$
restricciones duras $\to$ adjudicación LLM $\to$ revisión; cambia la fuente
de verdad (terminología curada frente a corpus) y, con ella, la necesidad
de curación humana.

\paragraph{Grafo clínico.}
El grafo resultante tiene $22\,807$ nodos de cinco clases ---aserción
($8\,444$), entidad ($7\,737$), concepto SNOMED ($2\,875$, prefijo
\texttt{sct:}), fragmento ($2\,419$) y recomendación ($1\,332$)--- y
$51\,071$ aristas tipadas: además de las clases del caso general
(procedencia, roles, menciones), aparecen \textsc{anclada\_a} (entidad
$\to$ concepto, $2\,326$), \textsc{is\_a} entre conceptos ($2\,649$, con
poda de concentradores de más de $300$ descendientes),
\textsc{contiene\_recomendacion} y \textsc{mencionada\_en\_rec}
($8\,253$). La sinonimia se resuelve aquí en régimen mixto
(sección~\ref{sec:dual}): aristas vectoriales a umbral $0{,}80$ más la
sinonimia \emph{verificada} ---entidades ancladas al mismo SCTID forman
estrella---, porque la fusión canónica completa exigiría la cobertura
terminológica que aún no se alcanza. Tres almacenes vectoriales acompañan
al grafo, como en el caso general: aserciones y entidades con
\texttt{text-embedding-3-small}, fragmentos con
\texttt{text-embedding-3-large}.

\subsection{Capa agéntica (consulta clínica)}
\label{sec:gesida-consulta}

La consulta la orquesta un grafo de estados (LangGraph) cuyos nodos
principales son:

\paragraph{Acogida y enrutado.}
La acogida seudonimiza datos personales y reescribe la pregunta a forma
autocontenida usando el contexto del hilo. Un enrutador LLM (con respaldo
heurístico por reglas) produce un plan de recuperación: tipo de consulta
(general o sobre un paciente), nivel estimado, guías candidatas, si es
multi-aspecto, si contiene lógica numérica, si es una decisión seria (que
exigirá pausa de supervisión) y la estrategia ---vectorial, de grafo o
paralela---. Es el análogo con modelo de lenguaje del enrutador léxico de
la sección~\ref{sec:router}: la misma decisión, con presupuesto distinto.

\paragraph{Recuperación de grafo.}
El recuperador de grafo es el mismo núcleo descrito en la
sección~\ref{sec:consulta} ---de hecho, es el original del que aquél se
portó---: siembra consulta-a-hecho con corrección anti-concentrador
(las 15 aserciones más afines aportan sus entidades, peso dividido por
frecuencia documental, 5 mejores), siembra densa de pasajes, anclas
simbólicas débiles (aquí extraídas por un NER-LLM de la consulta con dos
sinónimos por concepto, $\epsilon=0{,}05$) y PPR con la misma modulación
por consulta de la ecuación~\eqref{eq:modulacion}. Se añaden dos mecanismos
propios del dominio: \emph{semillas conceptuales} ---cada entidad de la
consulta que el enlazador SNOMED resuelve aporta su nodo concepto con peso
$1{,}0$ y sus ancestros IS-A con decaimiento $0{,}5^{d}$ (hasta 2 saltos),
todo escalado por $0{,}1$--- y \emph{sesgos por intención}: en preguntas
sobre un paciente, las aristas de población diana se refuerzan
($\times 1{,}5$); en preguntas numéricas, las transiciones hacia fragmentos
con umbrales ($\times 1{,}3$). El enrutador elige entre dos perfiles del
paseo: \emph{balanceado} (siembra densa activa, amortiguación $0{,}7$)
para precisión temprana, y \emph{grafo} (sin siembra densa, amortiguación
$0{,}85$, anclas fuertes) para cobertura profunda multi-guía. Cada
resultado lleva su traza de razonamiento legible
(\guillemotleft{}semilla [dolutegravir] $\to$ aserción [\textsc{treats}] $\to$
fragmento\guillemotright{}), que la interfaz muestra junto a la cita.

\paragraph{Fusión, reordenación y suficiencia.}
En estrategia paralela, las listas del grafo y del recuperador denso se
fusionan por rango recíproco (RRF). Un reordenador neuronal
(\emph{cross-encoder} multilingüe mmarco-MiniLM, local y sin coste de API)
reordena el conjunto con una pasada de \emph{diversidad}: primero el mejor
fragmento de cada guía implicada, después por puntuación pura ---medido:
el reordenador puro concentraba fragmentos casi idénticos de una guía y
expulsaba la cobertura multi-guía---. Un juicio de suficiencia decide
continuar, ampliar, cambiar de estrategia o abstenerse antes de generar.

\paragraph{Consultas de clase (subsunción).}
Las preguntas sobre \emph{clases} de fármacos (\guillemotleft{}?`qué inhibidores de la
integrasa se recomiendan en el embarazo?\guillemotright{}) no se responden bien por
semejanza: exigen enumerar los miembros de la clase y comprobar la
recomendación de \emph{cada uno}. Un agente de subsunción las detecta
(criba por reglas + traducción LLM a clase, población y polaridad) y
resuelve la pertenencia contra la jerarquía SNOMED por dos vías
---descendientes del concepto de clase y ancestros de cada entidad anclada
($\le 3$ saltos)---, cruza los miembros con las recomendaciones normativas
y devuelve también la \emph{exhaustividad negativa}: los miembros de la
clase \emph{sin} recomendación positiva, que en clínica es parte de la
respuesta.

\paragraph{Síntesis, verificación y supervisión.}
La síntesis genera con citas obligatorias (guía, sección, página, grado)
bajo reglas innegociables de vigencia (prioriza la guía más reciente y
señala obsolescencia) y de polaridad deóntica. La verificación combina un
juez LLM de fidelidad con dos \emph{lentes deterministas}: la deóntica
comprueba que ningún fármaco presentado como recomendado aparezca en una
recomendación citada con polaridad negativa (la inversión de polaridad es
el fallo clínicamente letal), y la numérica exige que todo umbral con
unidades de la respuesta aparezca en los fragmentos citados. Los hallazgos
de las lentes degradan el veredicto y se convierten en consultas de
\emph{recuperación dirigida} que realimentan el bucle. Las decisiones
serias pausan el flujo para supervisión médica (con reanudación que
conserva el estado), y la abstención es una salida de primera clase.

\subsection{Dos trazas clínicas}
\label{sec:gesida-trazas}

\paragraph{Indexación, paso a paso (bictegravir).}
(1)~El PDF de la guía de tratamiento antirretroviral produce, entre otros,
el fragmento \texttt{TAR\_ADULTOS\_2022::p6::10} (sección de pautas
preferentes), con metadatos de fármacos y grado.
(2)~El NER clínico extrae \guillemotleft{}Bictegravir\guillemotright{} (fármaco) y la reificación
produce una aserción clínica de tipo recomendación terapéutica con su
frase autocontenida, población diana y nivel de evidencia.
(3)~El extractor normativo detecta la frase con grado y crea el nodo
Recomendación (acción \emph{recomendado}, vigencia 2022).
(4)~El enlazador terminológico resuelve \guillemotleft{}Bictegravir\guillemotright{} por coincidencia
exacta en RXNORM: SCTID \texttt{772193003}, puntuación $1{,}0$; la prueba
de control de jerarquía comprueba que bictegravir subsume en $\le 2$ saltos
IS-A en un concepto de \guillemotleft{}inhibidor de la integrasa\guillemotright{}. El nodo
\texttt{sct:772193003} y sus ancestros entran al grafo. Un contraste con el
escalón difuso: \guillemotleft{}Tenofovir Disoproxil Fumarato\guillemotright{} ancla por trigramas con
Dice $0{,}881$ y queda marcado \emph{fuzzy} en la bandeja de curación.
(5)~Se incrustan las aserciones nuevas y se reconstruye el grafo.

\paragraph{Consulta, paso a paso (clase + deóntica).}
Ante \guillemotleft{}?`qué inhibidores de la integrasa se recomiendan en el
embarazo?\guillemotright{}: el enrutador la marca multi-guía; la criba de clase la deriva
al agente de subsunción, que traduce (clase: inhibidores de la integrasa;
población: embarazo; polaridad: recomendado), enumera los miembros por la
jerarquía SNOMED y cruza con las recomendaciones de la guía de embarazo.
El resultado ---raltegravir con recomendación positiva (grado C-III) y el
resto de miembros de la clase sin recomendación positiva en el corpus
(guía de 2018)--- entra como contexto tanto en la síntesis como en el
verificador (sin lo segundo, el verificador marcaba la tabla de clase como
sin respaldo y abstenía). La respuesta final cita cada afirmación y
explicita la exhaustividad negativa. En una pregunta numérica
(\guillemotleft{}?`a partir de qué filtrado glomerular puede usarse tenofovir
alafenamida?\guillemotright{}), el camino difiere: el enrutador marca lógica numérica, el
paseo refuerza los fragmentos con umbrales ($\times 1{,}3$), y la lente
numérica comprueba a posteriori que la cifra con unidades de la respuesta
figura en los fragmentos citados; si no, la afirmación pasa a recuperación
dirigida.

\subsection{Correspondencia entre las dos instancias del patrón}

La tabla~\ref{tab:gesida-wiki} alinea, dimensión a dimensión, el asistente
clínico y el sistema evaluado en 2Wiki. La lectura conjunta deja claro qué
es \emph{patrón} y qué es \emph{dominio}: comparten la unidad semántica
reificada, la resolución de entidades previa al grafo (con fuentes de
verdad distintas), el esquema de siembra y el paseo modulado; difieren en
las capas que cada dominio exige (deóntica, vigencia, verificación y
supervisión en clínica; nada de ello en el banco enciclopédico) y en la
dirección de la transferencia pendiente: el filtro selector con compuerta
dura y la fusión canónica, validados en 2Wiki, aún no se han retro-portado
a producción.

\begin{table}[t]
\centering\footnotesize\setlength{\tabcolsep}{3.5pt}
\caption{El mismo patrón en dos dominios: el asistente clínico GeSIDA y el
sistema evaluado en 2Wiki. Las filas superiores son el patrón compartido;
las inferiores, lo que cada dominio añade o deja fuera.}
\label{tab:gesida-wiki}
\begin{tabular}{p{0.20\textwidth}p{0.37\textwidth}p{0.35\textwidth}}
\toprule
dimensión & asistente GeSIDA & sistema evaluado (2Wiki)\\
\midrule
corpus & 8 guías clínicas, $2\,694$ fragmentos estructurales & $6\,119$
  artículos enciclopédicos\\
unidad semántica & aserción clínica reificada: roles fijos
  (intervención, resultado, población diana) + deóntica, grado, vigencia,
  negación & aserción $n$-aria genérica: roles abiertos + calificadores\\
resolución de entidades & terminología externa (SNOMED~CT/UMLS):
  exacto $\to$ difuso por trigramas (Dice $\ge 0{,}55$) $\to$ grupo
  semántico $\to$ juez LLM $\to$ curación humana; $30\%$ ancladas &
  diccionario derivado del corpus: candidatos vectoriales $\to$
  restricciones duras $\to$ adjudicación LLM; sin curación\\
sinonimia residual & aristas vectoriales ($\ge 0{,}80$) + estrella por
  mismo SCTID & ninguna: la fusión canónica la elimina\\
capa conceptual & jerarquía IS-A SNOMED (peso $0{,}5$; semillas con
  decaimiento $0{,}5^{d}$); agente de subsunción para clases & no aplica\\
anclaje de la consulta & NER-LLM + enlazador terminológico & enlazador de
  alias + respaldo vectorial (sin LLM)\\
filtro de hechos & no adoptado aún (siembra blanda anti-concentrador) &
  selector con demostraciones + compuerta dura (\S\ref{sec:filtro})\\
paseo (PPR) & mismo núcleo y modulación; sesgos por intención clínica;
  perfiles balanceado/grafo & mismo núcleo; diales de paridad con
  \hippo{}\\
posproceso & fusión RRF + reordenador neuronal con diversidad por guía;
  verificación (juez + lentes) + supervisión humana & ninguno: se evalúa
  recuperación pura\\
métricas & cobertura de guías y secciones por nivel clínico & $\mathrm{R@}k$
  y $\mathrm{FC@}k$ sobre pasajes oro\\
\bottomrule
\end{tabular}
\end{table}

\subsection{Métricas propias y estado del prototipo}

Sobre el banco de $505$ preguntas, la métrica principal es la cobertura de
guías (\guillemotleft{}todas las guías esperadas en el top-$k$\guillemotright{}, el análogo clínico de
la cadena completa) con desglose por nivel. En muestra estratificada de
100: el recuperador de grafo domina la cobertura profunda
(todas@20 $96\%$ frente a $90\%$ de BM25 y $85\%$ del denso) y, en las
preguntas de nivel 4 ---las multi-guía difíciles, el análogo de las
cadenas de 4 pasajes de 2Wiki---, el perfil grafo alcanza $84{,}9\%$@20
donde el denso queda en $24{,}7\%$. La configuración de producción (fusión
+ reordenador diverso) rinde todas@5 $92\%$ y alguna@1 $95\%$ (MRR de guía
$0{,}975$). Una ablación descartó el recorrido iterativo tipo
HopRAG~\cite{hoprag} como recuperador ($53{,}4\%$@20 sin LLM; $26\%$ con
presupuesto de 20 llamadas, frente a $84{,}9\%$ del paseo global); sus
$10\,480$ aristas-pregunta se conservan como artefacto de trazabilidad.
Limitaciones actuales del prototipo: la cobertura terminológica ($30\%$ de
entidades con SCTID) acota la capa conceptual; el reordenador neuronal no
está incluido en la imagen de despliegue (degradación silenciosa al orden
de fusión); la validación clínica con patrón oro está pendiente; y las
mejoras validadas en 2Wiki ---el filtro selector con compuerta dura y la
fusión canónica de entidades--- están pendientes de retro-portar a
producción, donde el análisis de este trabajo predice que atacarían
precisamente el residuo duro medido (precisión temprana de sección).

% =========================================================================
\section{Protocolo de evaluación}
\label{sec:protocolo}

\paragraph{Conjunto de datos y tarea.}
Evaluamos sobre 2WikiMultihopQA~\cite{wiki2} en el subconjunto de
reproducción publicado con \hippo{}: $1\,000$ preguntas ($413$ composicionales, $244$ de
comparación, $235$ de comparación-puente y $108$ de inferencia) y un corpus
de $6\,119$ pasajes. El protocolo es de \emph{recuperación abierta}: se
recupera contra el corpus completo; el contexto de 10 párrafos que acompaña
a cada pregunta no se utiliza. El oro de cada pregunta son los títulos de
sus pasajes de evidencia (\texttt{supporting\_facts}): $2$ pasajes por
pregunta, salvo en comparación-puente, que exige $4$ (dos películas y dos
directores).

\paragraph{Métricas.}
Para una pregunta $q$ con oro $G_q$ y lista devuelta $T_k(q)$ (los $k$
primeros títulos):
\begin{equation}
\mathrm{R@}k \;=\; \frac{|G_q \cap T_k(q)|}{|G_q|},
\qquad\qquad
\mathrm{FC@}k \;=\; \mathbf{1}\bigl[\,G_q \subseteq T_k(q)\,\bigr].
\label{eq:metricas}
\end{equation}
$\mathrm{R@}k$ es la exhaustividad estándar; $\mathrm{FC@}k$
(\emph{full chain}, cadena completa) exige \emph{todos} los pasajes oro en el
top-$k$ y es la métrica alineada con la tarea multisalto: un lector no puede
componer la respuesta si le falta un eslabón. Ambas se agregan globalmente y
por tipo, con intervalos de confianza del $95\%$ por bootstrap
($B=2\,000$ remuestreos de preguntas, semilla fija).

\paragraph{Separación de datos.}
Las $1\,000$ preguntas del banco de pruebas son intocables: ninguna decisión
de diseño, dial ni consigna se ajustó con ellas. Del conjunto de desarrollo
etiquetado oficial ($12\,576$ preguntas, disjunto del banco) se derivan:
(i)~un \emph{sondeo} estratificado de $200$ preguntas ($50$ por tipo) con su
propio mini-corpus de $1\,558$ pasajes, usado para toda la iteración de
diseño y la validación previa a cada medición; y (ii)~un fondo de
\emph{calibración} ($11\,376$ preguntas) usado para el enrutador léxico.
Cada configuración se midió sobre el banco de pruebas \emph{una sola vez}.

\paragraph{Comparación pareada.}
Además de los agregados, todas las comparaciones entre sistemas son pareadas
por pregunta: para cada métrica $m$ se calcula la diferencia
$\delta_i = m_A(q_i) - m_B(q_i)$, su media, su intervalo de confianza del
$95\%$ por bootstrap pareado ($B=2\,000$) y el recuento de preguntas que
gana o pierde cada sistema. Una diferencia se declara significativa si el
intervalo no cruza el cero.

\paragraph{Coste.}
Todo el estudio (extracción, diccionario, grafos, todas las mediciones de
ambos sistemas) se ejecutó con \texttt{gpt-4o-mini} y
\texttt{text-embedding-3-small} por un coste total aproximado de
$12$~USD. En consulta, \sys{} realiza $\sim\!1$ llamada al modelo de
lenguaje (el filtro; $2$ en comparación por la descomposición); la
indexación requiere $1$ llamada por pasaje (extracción) más la adjudicación
del diccionario (lotes de $8$ pares). Todas las llamadas se cachean
(SQLite/npz), de modo que cualquier re-medición es reproducible a coste
cero.

% =========================================================================
\section{Resultados y discusión}
\label{sec:resultados}

\subsection{Comparación controlada con \hippo{}}
\label{sec:comparacion}

\subsubsection{Paridad experimental}

\hippo{} se ejecutó como línea de base fuerte \emph{con su propio código}
(repositorio oficial vendorizado), indexando el mismo corpus y respondiendo
las mismas preguntas, y con los mismos modelos que \sys{}
(\texttt{gpt-4o-mini} y \texttt{text-embedding-3-small}). Adicionalmente,
para aislar la contribución de la representación y no de los diales, \sys{}
\emph{adoptó los valores publicados de \hippo{}} para los tres parámetros
compartidos del paseo (amortiguación $\lambda=0{,}5$, peso denso de pasaje
$0{,}05$, siembra de todos los pasajes), como decisión pre-registrada antes
de la medición final. La comparación es, por tanto, entre representaciones y
mecanismos, no entre presupuestos de cómputo ni ajustes de hiperparámetros.

\subsubsection{Correspondencia y diferencias mecanismo a mecanismo}

Ambos sistemas comparten el esqueleto ---grafo derivado por extracción
abierta con modelo de lenguaje, puntuación consulta--hecho, filtro LLM de
hechos, fusión con recuperación densa y PPR personalizado--- lo que hace la
comparación especialmente informativa: las diferencias están localizadas en
cuatro decisiones de diseño (tabla~\ref{tab:comparacion}).

\begin{table}[t]
\centering\small
\caption{Correspondencia mecanismo a mecanismo entre \hippo{} y \sys{}.}
\label{tab:comparacion}
\begin{tabular}{>{\raggedright\arraybackslash}p{0.19\textwidth}>{\raggedright\arraybackslash}p{0.36\textwidth}>{\raggedright\arraybackslash}p{0.36\textwidth}}
\toprule
dimensión & \hippo{} & \sys{}\\
\midrule
unidad semántica &
triplete $(s,r,o)$ como arista entre nodos de frase; incrustación del
triplete en almacén paralelo &
aserción reificada como nodo (ec.~\ref{eq:asercion}): frase autocontenida
$\tau_a$, calificadores, procedencia y roles $n$-arios
(\S\ref{sec:asercion})\\
\addlinespace
resolución de entidades &
en consulta: $\sim\!172\,000$ aristas de sinonimia por umbral de coseno
($\ge 0{,}80$, hasta $2\,047$ vecinos) &
en indexación: fusión canónica por diccionario con adjudicación LLM
(\S\ref{sec:diccionario}); cero aristas de sinonimia\\
\addlinespace
anclaje de la consulta &
sin reconocimiento de entidades en consulta: siembran el sujeto y el
objeto de los tripletes aprobados, peso $s_q/|P_\varepsilon|$
(ec.~\ref{eq:siembra-hippo}); respaldo denso si el filtro no aprueba nada &
enlazador contra el diccionario (alias exacto + respaldo vectorial),
anclas fuertes y masa directa en el pasaje propio
(ec.~\ref{eq:pasaje-propio})\\
\addlinespace
filtro de hechos &
memoria de reconocimiento sobre los tripletes de mayor coseno (consigna
optimizada con DSPy) &
selector con demostraciones sobre Top-40 por coseno $\cup$ frontera de las
entidades enlazadas (ec.~\ref{eq:candidatos}); compuerta dura en la siembra\\
\addlinespace
paseo (PPR) &
$\lambda=0{,}5$, peso denso $0{,}05$, todos los pasajes sembrados; matriz
de transición fija en indexación &
idéntico (adoptado por paridad) $+$ modulación por consulta de las
transiciones por hecho (ec.~\ref{eq:modulacion})\\
\addlinespace
enrutado por intención &
no tiene &
enrutador léxico de coste cero; descomposición para comparaciones
(\S\ref{sec:router})\\
\addlinespace
coste por consulta &
$\sim\!1$ llamada LLM (filtro) &
$\sim\!1$ llamada LLM (filtro; $2$ en comparación)\\
coste de indexación &
$2$ llamadas LLM por pasaje (NER + tripletes) &
$1$ llamada por pasaje (extracción) $+$ adjudicación del diccionario\\
\bottomrule
\end{tabular}
\end{table}

Conceptualmente, las cuatro diferencias son dos apuestas simétricas.
\hippo{} mantiene un grafo \emph{barato y ruidoso} (tripletes planos,
menciones sin resolver) y compensa el ruido en el momento de la consulta con
redundancia masiva (sinonimia de $172\,000$ aristas). \sys{} invierte el
orden: paga una vez, en indexación, por un grafo \emph{limpio y expresivo}
(aserciones $n$-arias verbalizadas sobre entidades canónicas), y eso permite
que los mecanismos de consulta sean estructurales en lugar de estadísticos:
el enlazador ancla la pregunta a nodos exactos, la frontera lleva al filtro
los hechos de esos nodos aunque el coseno no los vea, y el pasaje propio
convierte cada entidad puente descubierta en un pasaje sembrado.

En términos de las dos palancas de la sección~\ref{sec:palancas}, la
diferencia es igualmente nítida. \hippo{} inyecta la pregunta solo en la
fuente $\mathbf{v}$ ---con la entidad puente a peso atenuado,
$s_q/|P_\varepsilon|$--- sobre un coeficiente $P$ fijo; CatRAG añade un
coeficiente reponderado a una arista de cinco semillas. \sys{} actúa
sobre ambas palancas, pero reparte los papeles al revés que CatRAG: el
\emph{segundo salto} es asunto de la fuente ---la compuerta dura siembra a
la entidad puente a peso $1$ y la ecuación~\ref{eq:pasaje-propio} a su
biografía--- y el coeficiente se modula globalmente, sin juez por arista y
sin límite de alcance. La forma cerrada de la
ecuación~\eqref{eq:tres-canales} muestra por qué ese reparto es el
adecuado en el grafo reificado: $\mathbf{v}$ multiplica y
$\mu_a/Z_\varepsilon$ solo pondera, de modo que ningún coeficiente rescata
a un pasaje cuyo puente no fue sembrado.

\subsubsection{Resultados de la comparación}

La tabla~\ref{tab:resultados} resume la medición única sobre el banco de
pruebas. \sys{} supera a \hippo{} en todas las métricas globales
---$\mathrm{R@}5$ $96{,}0$ frente a $85{,}9$; $\mathrm{FC@}5$ $90{,}6$ frente
a $65{,}8$--- y en los cuatro tipos de pregunta. El análisis pareado
(sección~\ref{sec:protocolo}) da $\Delta\mathrm{FC@}5 = +24{,}8$ puntos con
IC del $95\%$ $[+21{,}8,\,+27{,}8]$ (gana en $279$ preguntas, pierde en
$31$), $\Delta\mathrm{R@}5 = +10{,}2$ (sig.) y
$\Delta\mathrm{R@}2 = +2{,}3$ (sig.). Por tipos, la ventaja es
significativa en comparación-puente ($+71{,}5$), composicional ($+13{,}8$)
y comparación ($+5{,}7$), y no significativa en inferencia
($+8{,}3$, $n=108$). El contraste más informativo es
comparación-puente ($\mathrm{FC@}5$ $83{,}8$ frente a $12{,}3$): recuperar
las cadenas de $4$ pasajes exige exactamente los mecanismos en los que
difieren los sistemas (entidades puente resueltas, hechos de frontera y
masa directa en los pasajes propios de los dos directores). Como referencia
externa, el artículo de \hippo{} reporta $\mathrm{R@}5=90{,}4$ en este
conjunto usando Llama-3.3-70B y el codificador NV-Embed-v2 (7B parámetros);
\sys{} alcanza $96{,}0$ con modelos de una fracción de ese coste.

\begin{table}[t]
\centering\small
\caption{Resultados sobre el banco de pruebas ($1\,000$ preguntas, corpus de
$6\,119$ pasajes; los tres sistemas con \texttt{gpt-4o-mini} +
\texttt{text-embedding-3-small}; \hippo{} con su código oficial).
$\mathrm{FC@}5$ por tipo: comparación, composicional, inferencia y
comparación-puente. En negrita, el mejor valor de cada columna.}
\label{tab:resultados}
\footnotesize\setlength{\tabcolsep}{3.5pt}
\begin{tabular}{lcccccccc}
\toprule
 & $\mathrm{R@}2$ & $\mathrm{R@}5$ & $\mathrm{FC@}5$ & $\mathrm{FC@}10$
 & comp. & compos. & infer. & p.-comp.\\
\midrule
BM25 & --- & $65{,}8$ & $32{,}8$ & $40{,}1$ & $86{,}5$ & $19{,}4$ & $32{,}4$ & $0{,}9$\\
\hippo{} & $70{,}7$ & $85{,}9$ & $65{,}8$ & $71{,}7$ & $93{,}4$ & $77{,}2$ & $75{,}9$ & $12{,}3$\\
\sys{} & $\mathbf{73{,}0}$ & $\mathbf{96{,}0}$ & $\mathbf{90{,}6}$ & $\mathbf{93{,}4}$
       & $\mathbf{99{,}2}$ & $\mathbf{91{,}0}$ & $\mathbf{84{,}3}$ & $\mathbf{83{,}8}$\\
\midrule
\multicolumn{9}{l}{\scriptsize pareado \sys{} $-$ \hippo{} (IC $95\%$):
$\mathrm{FC@}5$ $+24{,}8$ $[+21{,}8, +27{,}8]$;
$\mathrm{R@}5$ $+10{,}2$; $\mathrm{R@}2$ $+2{,}3$ (ambas sig.)}\\
\multicolumn{9}{l}{\scriptsize por tipo ($\mathrm{FC@}5$):
p.-comp.\ $+71{,}5$ (sig.); compos.\ $+13{,}8$ (sig.);
comp.\ $+5{,}7$ (sig.); infer.\ $+8{,}3$ (n.s., $n=108$)}\\
\bottomrule
\end{tabular}
\end{table}

\paragraph{Atribución incremental.}
Aunque las ablaciones formales componente a componente quedan como trabajo
en curso, la secuencia de versiones intermedias ---cada una validada en
sondeo y medida una única vez en el banco--- acota la contribución de cada
grupo de mecanismos ($\mathrm{R@}5$ / $\mathrm{FC@}5$): grafo reificado con
afinidad por coseno, sin filtro ($70{,}2$ / $41{,}5$); $+$ filtro selector
y enrutador ($73{,}3$ / $48{,}6$); $+$ diales de paridad y sinonimia
vectorial a $0{,}80$ ($75{,}6$ / $51{,}2$); $+$ selector con
demostraciones, compuerta dura y reificación $n$-aria completa
($90{,}9$ / $78{,}2$); $+$ diccionario canónico, enlazador, frontera y
pasaje propio ($96{,}0$ / $90{,}6$). Los dos saltos grandes corresponden a
las dos ideas centrales del sistema: el filtro con compuerta dura sobre el
grafo reificado y la memoria canónica de entidades.

% -------------------------------------------------------------------------

\subsubsection{Análisis por estructura de salto: acumulación frente a colapso}
\label{sec:saltos}

Los cuatro tipos de pregunta codifican tres \emph{estructuras} de salto (en
2WikiMultihopQA no existe el mono-salto): (A)~\emph{sin puente} ---comparación,
$n{=}244$: las dos entidades vienen nombradas y la evidencia son dos búsquedas
paralelas---; (B)~\emph{puente simple} ---composicional$+$inferencia,
$n{=}521$: una entidad puente ausente de la pregunta encadena dos saltos---; y
(C)~\emph{doble puente} ---comparación-puente, $n{=}235$: dos cadenas con
puente, cuatro pasajes oro en cinco posiciones. Reagrupando la
tabla~\ref{tab:resultados} por estructura, \hippo{} se degrada con la
profundidad ($\mathrm{FC@}5$: $93{,}4 \to 77{,}0 \to 12{,}3$) mientras \sys{}
se mantiene ($99{,}2 \to 89{,}6 \to 83{,}8$).

Un modelo nulo separa dos formas de degradar. Si cada pasaje oro entrara en el
top-5 de manera \emph{independiente} con probabilidad $p$, entonces
$\mathrm{FC@}5 \approx p^{\,|G_q|}$. Ajustando $p$ en el grupo B
($p_B=\sqrt{\mathrm{FC}_B}$) y extrapolando al grupo C ($p_B^{\,4}$):

\begin{center}\small
\begin{tabular}{lcccc}
\toprule
 & $p_B$ & C predicho $p_B^{\,4}$ & C observado & observado/predicho\\
\midrule
\hippo{} & $0{,}877$ & $59{,}2$ & $12{,}3$ & $0{,}21$\\
\sys{}   & $0{,}947$ & $80{,}3$ & $83{,}8$ & $1{,}04$\\
\bottomrule
\end{tabular}
\end{center}

\sys{} queda incluso ligeramente \emph{sobre} su extrapolación: su degradación
con la profundidad es la mera \emph{acumulación} de un error residual por
eslabón ($\approx 5\%$: juicio del filtro, huecos de extracción). \hippo{}
cae $4{,}8$ veces \emph{por debajo} de la suya: su fallo en C no es
acumulativo sino \emph{estructural} ---sus recursos por consulta son fijos y
compartidos entre las dos cadenas (el filtro selecciona a lo sumo cuatro
tripletes y en la práctica se concentra en una rama; la masa del paseo se
reparte entre las dos películas nombradas; y $k=5$ deja una única posición de
holgura para cuatro oros), de modo que cuando la selección cubre una sola
cadena el puente de la otra no se siembra nunca y $\mathrm{FC}=0$ queda
garantizado por diseño. El estudio de caso de la
sección~\ref{sec:fallan-otros} muestra este mecanismo sobre una pregunta
concreta. Los mecanismos de \sys{} eliminan justamente las dependencias de la
profundidad: identidad resuelta en indexación (con sinonimia por umbral, cada
eslabón adicional es una nueva oportunidad de \guillemotleft{}isla\guillemotright{}; con fusión canónica
el riesgo por eslabón es aproximadamente constante), hechos aprobados que
transitan a peso completo, y una contabilidad de la cadena que escala con
ella (la compuerta siembra las entidades de \emph{todos} los hechos
aprobados y el pasaje propio convierte cada puente descubierto en pasaje
sembrado). Que A no alcance el $100\%$ se debe a casi-duplicados del ranking
(secuelas, homónimos); que C quede bajo B incluso en \sys{} refleja la
ausencia de holgura con cuatro oros en cinco posiciones. El conteo explícito
de saltos de MuSiQue (2/3/4 en el identificador de cada pregunta) permitirá
contrastar directamente este modelo multiplicativo en la extensión
planificada.

\subsection{Posición frente al estado del arte publicado}
\label{sec:sota}

La tabla~\ref{tab:sota} sitúa el sistema frente a los resultados publicados
en 2Wiki. El bloque superior es directamente comparable: son los números
que reporta el artículo de CatRAG~\cite{catrag} sobre el \emph{mismo}
subconjunto de $1\,000$ preguntas y con los \emph{mismos} modelos que
usamos aquí (\texttt{gpt-4o-mini} + \texttt{text-embedding-3-small}), y su
métrica FCR (\emph{Full Chain Retrieval} a 5 pasajes) es exactamente
nuestra $\mathrm{FC@}5$. Dos comprobaciones de coherencia sostienen la
comparación: nuestra reproducción independiente de \hippo{} coincide con la
suya casi al decimal ($85{,}9$ / $65{,}8$ frente a $85{,}9$ / $66{,}1$), y
los números de Contriever y GTR coinciden entre ambos artículos de
referencia. La única divergencia es BM25 ($65{,}8$ nuestro frente a
$52{,}1$ suyo), atribuible a la implementación (la nuestra indexa el título
junto al texto); por prudencia la tabla lista el valor publicado.

El resultado central: con modelos pequeños, el sistema supera en $9{,}0$
puntos de exhaustividad y en $23{,}0$ puntos de cadena completa al mejor
sistema publicado en esta configuración (CatRAG: $87{,}0$ / $67{,}6$), y
supera también el $90{,}4$ de exhaustividad que \hippo{} reporta con un
codificador de $7.000$ millones de parámetros y Llama-3.3-70B como
extractor (bloque inferior). La evaluación extremo a extremo de respuesta
(F1 con lector compartido) queda pendiente en nuestro caso; las columnas de
F1 se listan como contexto.

\begin{table}[t]
\centering\footnotesize\setlength{\tabcolsep}{4pt}
\caption{2WikiMultihopQA, subconjunto de $1\,000$ preguntas. Bloque
superior: números publicados por~\cite{catrag} con \texttt{gpt-4o-mini} +
\texttt{text-embedding-3-small} (los mismos modelos de este trabajo); FCR
$=$ cadena completa a 5 pasajes $= \mathrm{FC@}5$; F1 con lector
Llama-3.3-70B. Bloque inferior: números publicados por~\cite{hipporag2} con
NV-Embed-v2 (7B) como recuperador. Las filas \guillemotleft{}(nuestra medición)\guillemotright{} son
de este trabajo; el resto se cita de los artículos y no se re-ejecutó aquí.}
\label{tab:sota}
\begin{tabular}{lccc}
\toprule
sistema & $\mathrm{R@}5$ & $\mathrm{FC@}5$ (FCR) & F1\\
\midrule
\multicolumn{4}{l}{\emph{mismos modelos que este trabajo
(\texttt{gpt-4o-mini} + \texttt{te3-small})~\cite{catrag}}}\\
BM25~\cite{catrag} & $52{,}1$ & $20{,}5$ & $39{,}9$\\
Contriever~\cite{contriever} & $57{,}5$ & $27{,}2$ & $41{,}9$\\
GTR~\cite{gtr} & $67{,}9$ & $35{,}8$ & $52{,}8$\\
denso \texttt{te3-small} & $70{,}8$ & $41{,}6$ & $56{,}9$\\
RAPTOR~\cite{raptor} & $69{,}8$ & $40{,}1$ & $56{,}7$\\
LightRAG~\cite{lightrag} & --- & --- & $49{,}7$\\
HyperGraphRAG~\cite{hypergraphrag} & --- & --- & $61{,}7$\\
PropRAG~\cite{proprag} & $83{,}5$ & $62{,}8$ & $63{,}9$\\
\hippo{}~\cite{hipporag2} & $85{,}9$ & $66{,}1$ & $68{,}1$\\
CatRAG~\cite{catrag} & $87{,}0$ & $67{,}6$ & $69{,}7$\\
\hippo{} (nuestra reproducción) & $85{,}9$ & $65{,}8$ & ---\\
\textbf{\sys{} (nuestra medición)} & $\mathbf{96{,}0}$ & $\mathbf{90{,}6}$ & pendiente\\
\midrule
\multicolumn{4}{l}{\emph{con NV-Embed-v2 (7B) + Llama-3.3-70B~\cite{hipporag2}}}\\
NV-Embed-v2~\cite{nvembed} & $76{,}5$ & --- & $61{,}5$\\
GraphRAG~\cite{graphrag} & --- & --- & $58{,}6$\\
HippoRAG~\cite{hipporag} & $90{,}4$ & --- & $71{,}8$\\
\hippo{}~\cite{hipporag2} & $90{,}4$ & --- & $71{,}0$\\
\bottomrule
\end{tabular}
\end{table}

% -------------------------------------------------------------------------
\subsection{Dónde fallan los demás y por qué este sistema prospera}
\label{sec:fallan-otros}

Los números agregados de la tabla~\ref{tab:sota} esconden tres modos de
fallo distintos, uno por familia.

\paragraph{Recuperación densa y léxica: no puede componer.}
El vector de la pregunta no puede parecerse al documento puente. El caso
extremo es comparación-puente: BM25 recupera la mitad de la cadena
($\mathrm{R@}5$ $50{,}9$ en ese segmento de nuestra medición) pero
prácticamente nunca la completa ($\mathrm{FC@}5$ $0{,}9$): encuentra las
dos películas nombradas y ningún director. Es el patrón
\guillemotleft{}exhaustividad parcial alta, cadena rota\guillemotright{} que motiva la métrica FCR.

\paragraph{La línea HippoRAG: la deriva a una sola rama.}
\hippo{} resuelve el patrón anterior en cadenas de 2 (composicional
$77{,}2$ de $\mathrm{FC@}5$) pero se hunde en las de 4 ($12{,}3$). Un caso
real de nuestra medición pareada lo ilustra ---es el mismo ejemplo
conductor que la sección~\ref{sec:consulta} sigue etapa a etapa---.
Pregunta:
\guillemotleft{}\emph{Which film has the director who is older, God's Gift to Women or
Aldri Annet Enn Bråk?}\guillemotright{} (oro: las dos películas y sus directores, Michael
Curtiz y Edith Carlmar). El top-5 de \hippo{} contiene las dos películas y
a Edith Carlmar, y sus puestos 4--10 son \emph{otras películas de Edith
Carlmar} (\emph{Altid ballade}, \emph{Bedre enn sitt rykte}) y directores
ajenos: el paseo entró en la rama de Carlmar, cuyo nodo-frase concentra
aristas, y la masa nunca alcanzó a Michael Curtiz (ausente del top-10).
La causa es cuantificable en su matriz de transición, fijada en
indexación (sección~\ref{sec:palancas}): de \emph{Aldri annet enn bråk} a
Edith Carlmar el paseo cruza con probabilidad $0{,}316$ (tres tripletes
entre ambas), de \emph{God's Gift to Women} a Michael Curtiz con
$0{,}105$ (un triplete, la misma probabilidad que hacia
\guillemotleft{}musicals\guillemotright{}), y desde Carlmar cada una de sus otras películas
recibe $0{,}037$, más de lo que la biografía de Curtiz recibe de su propio
nodo ($1/103$). Es la deriva hacia concentradores que el propio artículo
de CatRAG diagnostica, agravada porque las dos ramas compiten por la
misma masa.
\sys{} devuelve los cuatro pasajes oro en los puestos 1--4: el filtro
aprueba \emph{los dos} hechos de dirección (la consigna obliga a conservar
todos los candidatos de un papel), la compuerta dura siembra a los dos
directores con peso $1$, y el término de pasaje propio
(ec.~\ref{eq:pasaje-propio}) deposita masa de reinicio directamente en las
dos biografías, sin depender de que el paseo las alcance. En el agregado,
este mecanismo explica el grueso de la ventaja pareada: $+71{,}5$ puntos
de $\mathrm{FC@}5$ en comparación-puente (gana $107$ preguntas, pierde
$5$).

\paragraph{CatRAG: el recorrido correcto sobre el grafo equivocado.}
CatRAG comparte nuestro diagnóstico (grafo estático, concentradores,
cadena rota) y ataca el \emph{recorrido}: anclas, aristas reponderadas por
consulta, refuerzo de pasajes. Su ganancia es real pero acotada
($\mathrm{FC@}5$ $66{,}1 \to 67{,}6$; su ablación atribuye $0{,}9$ y
$1{,}4$ puntos de exhaustividad a las anclas y a la reponderación, y el
refuerzo de pasajes resta $1{,}4$ en 2Wiki), por dos razones que la
sección~\ref{sec:palancas} desarrolla. La primera es de alcance: la
reponderación llega a una arista de las semillas; en la pregunta anterior
podría amplificar película$\to$Curtiz, pero la fila de Curtiz seguiría
siendo la estática y su biografía quedaría a una arista de la cadena,
sin refuerzo de pasaje, porque ningún triplete aprobado la respalda. La
segunda es de representación: opera sobre tripletes planos y entidades
fragmentadas, y cuando la mención del puente y su biografía son nodos
distintos sin arista fiable, ningún recorrido, por adaptado que esté,
puede cruzar. Nuestros resultados sugieren que la
palanca grande estaba en la representación: con el grafo reificado sobre
entidades canónicas, los mismos mecanismos de recorrido (que heredamos de
esa línea) rinden $90{,}6$ de cadena completa donde el recorrido solo
llegaba a $67{,}6$.

% -------------------------------------------------------------------------
\subsection{Dónde falla todavía}
\label{sec:falla-aun}

Quedan $94$ fallos de cadena completa a 5 ($9{,}4\%$ del banco de
pruebas): $38$ de comparación-puente, $37$ composicionales, $17$ de
inferencia y $2$ de comparación. De los $102$ pasajes oro perdidos, $88$
son el pasaje \emph{puente} ($64$ fuera del top-20, $18$ en puestos
6--10): el residuo sigue concentrado en el segundo salto.

\paragraph{Un fallo real, disecado.}
\guillemotleft{}\emph{Which film has the director who was born later, Playing It Wild or
I'll Be Going Now?}\guillemotright{}. El sistema devuelve tres de los cuatro oros en el
top-4, pero la biografía de Dino Risi (director de \emph{I'll Be Going
Now}) cae fuera del top-20. La disección muestra una cadena de causas
instructiva: el hecho puente \emph{existe} en el grafo, pero su frase está
verbalizada con el título italiano de la película
(\guillemotleft{}\emph{Tolgo il disturbo was directed by Dino Risi}\guillemotright{}, fiel al pasaje
original); el diccionario no puede fusionar \guillemotleft{}I'll Be Going Now\guillemotright{} con
\guillemotleft{}Tolgo il disturbo\guillemotright{} porque la fusión exige un término significativo
común y un alias \emph{translingüe} no comparte ninguno; en consecuencia la
frontera de la entidad enlazada no contiene el hecho de dirección, el
coseno tampoco lo acerca (la pregunta usa el título internacional), y el
filtro nunca lo ve. Como distractor, un hecho de otra película con
\guillemotleft{}Wild\guillemotright{} en el título coloca a Cornel Wilde en el puesto 2. La clase de
fallo ---alias que ninguna semejanza de superficie ni de contexto breve
puede unir--- es exactamente la que una capa de identidad \emph{externa}
(Wikidata aquí; SNOMED en el caso clínico) resolvería.

\paragraph{El cuello de botella medido: el juicio del filtro.}
El estudio de errores instrumentado sobre los fallos del sondeo (retenido)
atribuye la mayoría del residuo no al grafo sino al filtro: en $12$ de $15$
fallos de puente, el hecho correcto \emph{estaba} en la lista de candidatos
y el modelo pequeño no lo seleccionó. El caso trazado:
en \guillemotleft{}\emph{Which film has the director born earlier, Bílá spona or A Night
at Earl Carroll's?}\guillemotright{}, el filtro eligió como director al \emph{actor} Earl
Carroll y omitió \guillemotleft{}\emph{Bílá spona was directed by Kurt Neumann}\guillemotright{}
aunque figuraba en la lista; con una consigna de dos fases
(planificar la cadena y después seleccionar, tope de 8 hechos) el mismo
modelo seleccionó ambos puentes. El tope de 4 hechos es además
estructuralmente insuficiente en comparación-puente, que necesita al menos
4 (dos direcciones y dos fechas). Las líneas de mejora priorizadas
---filtro planifica-y-selecciona, extendido a la estrategia de comparación
(donde la compuerta blanda siembra homónimos, sección~\ref{sec:siembra}),
selección final de \emph{conjunto} sobre
el top-10 (30 oros del banco quedan en puestos 6--10), segundo salto
dirigido cuando el plan nombra un hueco, y enlazador por NER-LLM (la regla
de mayúsculas produjo el tramo roto \guillemotleft{}Magician (1958\guillemotright{} en un fallo)---
atacan exactamente este perfil; el Plan A, que implementa las tres
primeras, se describe y mide en la sección~\ref{sec:evoluciones}. Las
evoluciones de la memoria canónica se recogen, ordenadas, en la
sección~\ref{sec:memoria}, y todas juntas, con la adaptación clínica y la
generalización, en la sección~\ref{sec:futuro}.

% =========================================================================
\section{Evoluciones, adaptación al asistente GeSIDA y generalización}
\label{sec:futuro}

Los apartados anteriores dejan planteadas evoluciones dispersas por
componente (memoria canónica, E1--E9; enlazador, L1--L4; candidatos,
siembra, paseo). Esta sección las reúne en un solo lugar y las separa
con claridad de lo que ya está ejecutado y medido; después traduce los
mecanismos validados en 2Wiki a un plan de adaptación concreto para el
asistente clínico; y cierra con el protocolo de generalización a otros
conjuntos de datos, con las referencias publicadas contra las que
comparar. El criterio es el mismo en las tres partes: cada evolución
lleva la evidencia que la motiva, el efecto esperado, el coste y cómo se
mide, y nada se ajusta sobre el banco de pruebas.

\subsection{Evoluciones del sistema}
\label{sec:evoluciones}

\paragraph{Qué está ejecutado y qué no.}
El sistema descrito en las secciones~\ref{sec:metodos}--\ref{sec:resultados}
es la versión v3 ---memoria canónica, enlazador, frontera, filtro con
compuerta dura, pasaje propio---, medida en paridad estricta con \hippo{}
($96{,}0$ / $90{,}6$). Sobre ella se ha implementado y medido una
evolución más, el \emph{Plan A}, que ataca el cuello de botella
identificado en la sección~\ref{sec:falla-aun} ---el juicio del filtro---
sin tocar la representación. El resto de evoluciones de esta sección son
propuestas con evidencia, no implementadas.

\paragraph{Evolución ejecutada: el Plan A (analista, salto dirigido y conjunto).}
El Plan A extiende la v3 sin modificarla y sustituye la llamada del
filtro por tres mecanismos:
\begin{enumerate}[leftmargin=1.6em, itemsep=2pt]
\item \textbf{Analista de la pregunta} (una llamada, siempre). Recibe la
  pregunta y la lista de candidatos $\mathcal{K}(q)$ y devuelve, en JSON
  estricto: las \emph{menciones} de la pregunta con su forma canónica
  ---el reconocimiento de entidades que el enlazador por mayúsculas no
  hace (evolución L1)---; el \emph{plan} de la cadena como lista de
  huecos (\guillemotleft{}director de $A$ $\to$ $X$\guillemotright{}, \guillemotleft{}fecha de muerte de
  $X$\guillemotright{}), con la prohibición explícita de rellenar un hueco con
  conocimiento paramétrico (en la iteración previa el analista inventaba
  directores); la \emph{selección} de hechos que rellenan huecos, con
  tope dinámico igual al número de huecos (máximo $8$, frente al tope
  fijo de $4$ del filtro, insuficiente en cadenas de cuatro); y los
  \emph{huecos sin cubrir}, redactados como consulta de recuperación.
  Absorbe en una sola llamada el filtro selector y el enlazador.
\item \textbf{Salto dirigido} (una llamada, solo si quedan huecos). Cada
  hueco sin cubrir se incrusta y se buscan sus hechos más afines entre
  candidatos \emph{estructurales}: las aserciones del pasaje propio de
  las entidades enlazadas o descubiertas. Un mini-filtro los adjudica y
  se vuelve a sembrar. Es lo que rescata el puente verbalizado en otro
  idioma (\emph{Ansiktet} / \emph{The Magician}), inalcanzable para el
  coseno y para la frontera.
\item \textbf{Selección final de conjunto} (una llamada; se omite en
  comparación). Del top-12 del paseo, el modelo elige los $5$ pasajes que
  \emph{juntos} cubren el plan, conservando siempre el top-2 del paseo
  como guarda. Responde a los $30$ oros que la v3 dejaba en los puestos
  6--10.
\end{enumerate}
El coste es de $\sim 2{,}2$ llamadas por consulta; con
\texttt{deepseek-chat} (V3), $\sim 0{,}001$~USD por consulta y $\sim
1$~USD la medición completa. Todas las llamadas se guardan en caché
SQLite por (modelo, resumen del mensaje), de modo que las re-mediciones
son gratuitas.

La tabla~\ref{tab:plan-a} resume la evidencia. En el sondeo (retenido) la
ganancia es del \emph{diseño} y no del modelo: \texttt{gpt-4o-mini} pasa
de $92{,}0$ a $96{,}0$ de $\mathrm{FC@}5$, \texttt{gpt-4o} solo añade
$0{,}5$ más a quince veces el coste, y un modelo abierto de $20\,000$
millones de parámetros (\texttt{gpt-oss-20b}) empata de hecho con
DeepSeek; un modelo que no produce planes válidos (\texttt{qwen3.6-27b})
marca el suelo. Sobre el banco, con \texttt{deepseek-chat}, el Plan A
alcanza $\mathrm{R@}5$ $98{,}6$ y $\mathrm{FC@}5$ $97{,}1$;\footnote{La
medición del 25 de agosto dio $98{,}5$ / $96{,}9$; la re-ejecución
posterior con las mismas consignas dio $98{,}6$ / $97{,}1$. La
diferencia es de dos preguntas y refleja el no determinismo residual del
modelo.} pareado frente a \hippo{}, $\Delta\mathrm{FC@}5=+31{,}3$
$[+28{,}4,\,+34{,}4]$ (gana $322$, pierde $9$) y, frente a la v3,
$+6{,}5$ $[+4{,}7,\,+8{,}2]$ (gana $74$, pierde $9$), significativo en
todos los tipos salvo comparación ($+0{,}4$, ya saturada). Las
ablaciones aditivas en el sondeo atribuyen la ganancia: el analista solo
aporta $+0{,}5$ ---su valor es producir el plan que consumen los otros
dos---, el salto dirigido es el mayor contribuyente individual
($+3{,}0$; doble puente $86 \to 96$) y el conjunto final aporta orden
temprano ($+1{,}5$ de $\mathrm{FC@}5$ y $+6$ de $\mathrm{FC@}2$), con
ligera superaditividad ($5{,}0$ frente a $5{,}5$ observado).

El hallazgo conceptual es que el gradiente de dificultad de la
sección~\ref{sec:saltos} se \emph{invierte}: por estructura de salto, el
Plan A rinde $99{,}6$ (sin puente), $95{,}0$ (puente simple) y $99{,}1$
(doble puente). El segmento que era el peor de todos los sistemas
publicados ($0{,}9$ BM25, $12{,}3$ \hippo{}) es ahora el mejor, y queda
muy por encima del modelo multiplicativo ($p_B=0{,}975$ predice $90{,}3$
para cuatro oros): el plan asigna recursos \emph{por cadena} ---huecos
explícitos, tope dinámico, conjunto final---, de modo que la doble
cadena no comparte un presupuesto fijo sino que recibe el doble. Es la
refutación constructiva del colapso estructural de \hippo{}.

Dos reservas mantienen al Plan A fuera de las tablas principales de este
documento. Primera, paridad: \texttt{deepseek-chat} no es
\texttt{gpt-4o-mini}; la comparación en igualdad exacta de modelos es la
del sondeo ($96{,}0$ con \texttt{gpt-4o-mini}) y, en el banco, la v3.
Segunda, contaminación de diseño: las demostraciones del analista deben
re-derivarse del conjunto de calibración y medirse una sola vez
(\guillemotleft{}cifra de archivo\guillemotright{}, sección~\ref{sec:validez}) antes de
promoverlo a resultado principal.

\begin{table}[p]
\centering\footnotesize\setlength{\tabcolsep}{4pt}
\caption{Evidencia del Plan A. Arriba: matriz de modelos del analista en
el sondeo retenido ($200$ preguntas), con el coste aproximado por
consulta. Centro: ablaciones aditivas en el sondeo con
\texttt{deepseek-chat}. Abajo: banco de pruebas ($1\,000$), con la
comparación pareada (IC del $95\%$ por bootstrap; entre paréntesis,
preguntas que gana / pierde).}
\label{tab:plan-a}
\scriptsize\setlength{\tabcolsep}{3.5pt}
\begin{tabular}{lcccp{0.36\textwidth}}
\toprule
sondeo (200): modelo del analista & $\mathrm{R@}5$ & $\mathrm{FC@}5$ & USD/consulta & \\
\midrule
v3 (referencia; \texttt{gpt-4o-mini}) & $97{,}1$ & $92{,}0$ & $0{,}0004$ & \\
Plan A, \texttt{gpt-4o-mini} & $98{,}5$ & $96{,}0$ & $0{,}0008$ & misma paridad que la v3\\
Plan A, \texttt{gpt-4o} & $98{,}6$ & $96{,}5$ & $0{,}012$ & $+0{,}5$ a $15\times$ coste\\
Plan A, \texttt{gpt-oss-20b} (Groq) & $98{,}9$ & $97{,}0$ & $0{,}0005$ & alternativa abierta validada\\
Plan A, \texttt{deepseek-chat} & $\mathbf{99{,}0}$ & $\mathbf{97{,}5}$ & $0{,}0010$ & adoptado\\
Plan A, \texttt{qwen3.6-27b} & $85{,}6$ & $67{,}5$ & --- & planes vacíos: suelo\\
\midrule
ablaciones (sondeo, \texttt{deepseek-chat}) & $\mathrm{FC@}5$ & $\mathrm{FC@}2$ & doble puente & aporte\\
\midrule
v3 & $92{,}0$ & $46{,}0$ & $82{,}0$ & ---\\
$+$ analista & $92{,}5$ & $55{,}5$ & $86{,}0$ & $+0{,}5$ (habilitador)\\
$+$ analista $+$ conjunto & $94{,}0$ & $61{,}0$ & $86{,}0$ & $+1{,}5$\\
$+$ analista $+$ salto & $95{,}5$ & $55{,}0$ & $96{,}0$ & $+3{,}0$\\
Plan A completo & $\mathbf{97{,}5}$ & $\mathbf{61{,}5}$ & $\mathbf{96{,}0}$ & $+5{,}5$\\
\midrule
banco (1\,000) & $\mathrm{R@}2$ & $\mathrm{R@}5$ & $\mathrm{FC@}5$ & $\mathrm{FC@}5$ por tipo (comp./compos./infer./p.-comp.)\\
\midrule
\hippo{} & $70{,}7$ & $85{,}9$ & $65{,}8$ & $93{,}4$ / $77{,}2$ / $75{,}9$ / $12{,}3$\\
v3 (\texttt{gpt-4o-mini}) & $73{,}0$ & $96{,}0$ & $90{,}6$ & $99{,}2$ / $91{,}0$ / $84{,}3$ / $83{,}8$\\
Plan A (\texttt{deepseek-chat}) & $\mathbf{82{,}2}$ & $\mathbf{98{,}6}$ & $\mathbf{97{,}1}$ & $99{,}6$ / $96{,}1$ / $90{,}7$ / $99{,}1$\\
Plan A v4 (\texttt{gpt-4o-mini}; \S\ref{sec:v4}) & $80{,}2$ & $98{,}1$ & $96{,}0$ & $99{,}6$ / $95{,}9$ / $84{,}3$ / $97{,}9$\\
Plan A v4 (\texttt{deepseek-chat}; \S\ref{sec:v4}) & $\mathbf{84{,}2}$ & $\mathbf{98{,}9}$ & $\mathbf{97{,}7}$ & $100$ / $97{,}1$ / $92{,}6$ / $98{,}7$\\
\multicolumn{5}{p{0.93\textwidth}}{pareado Plan A $-$ \hippo{}: $\mathrm{FC@}5$ $+31{,}3$ $[+28{,}4,+34{,}4]$ (gana 322 / pierde 9); $\mathrm{R@}5$ $+12{,}7$; $\mathrm{R@}2$ $+11{,}5$; por tipo: p.-comp.\ $+86{,}8$, compos.\ $+18{,}9$, infer.\ $+14{,}8$, comp.\ $+6{,}1$ (todos significativos)}\\
\multicolumn{5}{p{0.93\textwidth}}{pareado Plan A $-$ v3: $\mathrm{FC@}5$ $+6{,}5$ $[+4{,}7,+8{,}2]$ (gana 74 / pierde 9); $\mathrm{R@}2$ $+9{,}2$; por tipo: p.-comp.\ $+15{,}3$, compos.\ $+5{,}1$, infer.\ $+6{,}5$ $[0{,}0,+13{,}0]$, comp.\ $+0{,}4$ (n.s.)}\\
\bottomrule
\end{tabular}
\end{table}

\paragraph{Evoluciones propuestas.}
La tabla~\ref{tab:evoluciones} reúne las evoluciones no implementadas,
ordenadas por componente, con la evidencia que las motiva, el efecto
esperado y cómo se mide, el coste y el estado. Tres criterios las
gobiernan. Primero, \emph{generalidad}: se descartan los parches
específicos del banco (reintentar el alias anteponiendo el artículo,
reglas léxicas por tipo) en favor de mecanismos que no dependen del
idioma ni de la forma de los títulos. Segundo, \emph{medida antes que
intuición}: donde una alternativa se pudo medir sin coste (enlazador,
frontera, presa densa) la cifra figura; donde no, se indica qué medición
la decidiría. Tercero, \emph{protocolo}: toda evolución se valida en el
sondeo y se mide una sola vez en el banco, con comparación pareada.

\begin{table}[p]
\centering\scriptsize\setlength{\tabcolsep}{3pt}
\caption{Evoluciones del sistema por componente. E$n$ y L$n$ remiten a
las listas de las secciones~\ref{sec:memoria} y~\ref{sec:linker}. Coste en
USD aproximados; \guillemotleft{}0\guillemotright{} significa sin llamadas nuevas (caché).}
\label{tab:evoluciones}
\begin{tabular}{p{0.09\textwidth}p{0.27\textwidth}p{0.19\textwidth}p{0.19\textwidth}p{0.06\textwidth}p{0.10\textwidth}}
\toprule
componente & evolución & evidencia & efecto esperado y medida & coste & estado\\
\midrule
extracción & consigna v2: sujeto $=$ forma del título $+$ alias; roles $n$-arios completos; descripciones de ficha redactadas por el modelo & el canal contexto mide pasaje y no entidad; $1\,342$ títulos sintéticos & menos uniones título$\to$sujeto necesarias; auditar con QIDs & $\sim 2{,}5$ & propuesta\\
\addlinespace
memoria canónica & E1 auditoría con QIDs; E2 procedencia y estado por alias; E3 decisiones registradas y regeneración determinista; E4 re-adjudicación de las $1\,549$ uniones de título; E5 candidatos por umbral; E6 canal contexto por rol; E7 colisiones; E8 identidad externa; E9 bandeja unificada. Nuevo: la ficha-título hereda las fechas del sujeto, y un nombre base idéntico con desambiguador entre paréntesis va al juez & sobre-fusión medida: parientes, secuelas, homónimos (\texttt{can:christopher newton}) & precisión y exhaustividad de fusión con oro externo; re-medición pareada & $1$--$2$ & E1, E4 ejecutadas; E2, E7 parciales (\S\ref{sec:v4})\\
\addlinespace
enlazador & L1 reconocimiento de entidades (absorbido por el analista); L2 escalón 2 de nombre a nombre: Dice $\ge 0{,}80$ $\vee$ coseno de nombre $\ge 0{,}85$, ficha como desempate; L3 BM25 solo como candidatos; L4 colisiones & escalón 2 inactivo ($3$ de $275$); la unión medida da $179$ correctos / $8$ erróneos & $+176$ enlaces correctos netos; $\mathrm{FC@}5$ pareado & $0$ & L1, L2 ejecutadas; L4 parcial; L3 propuesta (\S\ref{sec:v4})\\
\addlinespace
candidatos $\mathcal{K}(q)$ & frontera completa para entidades con $\le 30$ hechos; candidatos estructurales del pasaje propio (como en el salto dirigido) & $9\%$ de los hechos oro de primer salto quedan fuera de $\mathcal{K}(q)$ & cobertura de $\mathcal{K}(q)$ sin llamadas; después $\mathrm{FC@}5$ & $0$ & propuesta\\
\addlinespace
filtro & analista con plan y tope dinámico (ejecutado); extenderlo a la estrategia de comparación & la compuerta blanda siembra homónimos (David Newton) & orden temprano en comparación ($\mathrm{FC@}2$) & 1 llamada en el $24\%$ & ejecutada (\S\ref{sec:v4})\\
\addlinespace
siembra & (i) ablación de la presa densa y de $w_{\text{pas}}$; (ii) normalización por grado del pasaje, $\mathbf{p}^{*}(c)/d(c)^{\beta}$; (iii) agregación conjuntiva de semillas en los tipos con dos entidades nombradas & la presa densa es el $94$--$97\%$ de la masa; sesgo a pasajes con muchos hechos & $\beta\in\{0,\,0{,}25,\,0{,}5,\,1\}$ y conjunción condicionada por tipo, todo con caché & $0$ & propuesta\\
\addlinespace
paseo & ablación de la modulación ($\mu\equiv 1$) sobre el banco; perfil de pesos por longitud solo en los grafos con ciclos (clínico) & en la traza la modulación no cambia el orden; su contribución no está aislada & cuantificar $\mu$; en el grafo por capas la profundidad es $2$ & $0$ & propuesta\\
\addlinespace
coste (Plan B) & comparación sin LLM ($24\%$ de consultas a coste cero); compuerta de confianza (filtro solo si las entidades enlazadas no tienen su pasaje en el top-5); caché de prefijo de consigna; modelo más barato con regla dura $\mathrm{FC@}5\ge 92$ en sondeo & comparación $99{,}2$ en v3 sin depender del juez & $\sim 0{,}5$ llamadas por consulta a igual calidad & $0$ & propuesta\\
\addlinespace
evaluación & QA extremo a extremo (EM/F1 y JSR con lector compartido; referencias: JSR $53{,}0$ \hippo{}, $55{,}0$ CatRAG con Llama-70B); cifra de archivo con demostraciones re-derivadas de calibración; MuSiQue y HotpotQA (\S\ref{sec:generalizacion}) & columna F1 pendiente en la tabla~\ref{tab:sota} & comparabilidad completa con los artículos & $1$ / $1$ / $20$ & propuesta\\
\bottomrule
\end{tabular}
\end{table}

\paragraph{Evoluciones ejecutadas el 30 de agosto: curación auditada del
diccionario, enlazador de nombre a nombre y Plan A generalizado.}
\label{sec:v4}
Antes de indexar otro conjunto de datos se ejecutaron, con el protocolo de
siempre (validación en el sondeo, una sola medición en el banco), las
evoluciones que la tabla~\ref{tab:evoluciones} señalaba como prerrequisito de
la generalización: las que corrigen el índice ---porque un índice nuevo
heredaría sus defectos--- y las que retiran del flujo de consulta las reglas
específicas de 2Wiki. Todos los artefactos nuevos llevan el sufijo
\texttt{\_v4}; la versión v3 descrita en las secciones anteriores queda
intacta.

\emph{Auditoría con oro externo (E1).} La tabla~\ref{tab:auditoria-qid} mide
el diccionario contra los identificadores Wikidata que el propio conjunto trae
en cada pregunta. La unidad es el par (forma, QID) de las entidades de la
cadena ---sujetos de evidencia y títulos oro---; los objetos-literal
(gentilicios, lugares) quedan fuera porque su QID en el oro es inconsistente
(\guillemotleft{}American\guillemotright{} aparece bajo cuatro QIDs), y una misma forma bajo dos
QIDs (\guillemotleft{}Michael Curtiz\guillemotright{} bajo Q1559143 y Q51491) se excluye como ruido
del oro. La atribución usa la \emph{procedencia} de cada unión, que la v4
registra (E2) y que para la v3 se reprodujo en seco con la lógica original
($26\,932$ de $26\,934$ grupos idénticos). El resultado confirma el
diagnóstico de la sección~\ref{sec:diccionario}: en la v3, el origen que con
más frecuencia tiende el puente entre dos QIDs distintos es la unión
título$\to$sujeto ejecutada cuando el título ya era mención propia ($40$ de
$74$ entradas sobre-fundidas).

\begin{table}[t]
\centering\footnotesize\setlength{\tabcolsep}{4pt}
\caption{Auditoría del diccionario con los QIDs del oro (banco; $2\,487$ pares
forma--QID de entidades de la cadena, $2\,293$ resueltos a una entrada).
Sobre-fusión: entradas con dos o más QIDs de formas distintas; infra-fusión:
QIDs repartidos en dos o más entradas. \guillemotleft{}Origen puente\guillemotright{}: origen de la
unión menos frecuente en el camino que conecta las fichas de los dos QIDs.}
\label{tab:auditoria-qid}
\begin{tabular}{lrrrl}
\toprule
diccionario & entradas & sobre-fundidas & infra & origen puente de la sobre-fusión\\
\midrule
v3 & $26\,960$ & $74$ de $2\,211$ tocadas ($3{,}3\%$) & $6$ & título$\to$sujeto no sintética $40$; juez $19$; casi idéntico $18$; idéntico $16$\\
v4 & $27\,966$ & $\mathbf{32}$ de $2\,249$ & $6$ & juez $9$; título$\to$sujeto juzgada $8$; idéntico $8$; casi idéntico $5$\\
\bottomrule
\end{tabular}
\end{table}

\emph{Curación del diccionario (v4).} Cuatro cambios en la consolidación,
todos como decisiones registradas y no como parches: (i) la unión
título$\to$sujeto solo es automática cuando el título es mención
\emph{sintética} y el sujeto candidato es único; si el título ya es mención
propia ($1\,549$ casos) o hay empate ($202$), el par va al juez, que rechazó
$893$ de los $1\,549$ y aceptó $192$ de los $202$; (ii) la ficha-título sin
fechas hereda las del sujeto único del artículo ($911$ fichas), lo que
desactiva la fusión automática de homónimos como \texttt{can:christopher
newton}; (iii) un nombre base idéntico ---o casi idéntico--- con
desambiguador entre paréntesis distinto va al juez en lugar de fundirse
(\guillemotleft{}The Girl of the Golden West\guillemotright{} de 1915, 1922, 1923, 1938 y 1942 eran
una sola entrada); (iv) cada entrada guarda el recuento de uniones por
origen. Coste: $2\,325$ adjudicaciones nuevas, unos céntimos. La sobre-fusión
baja de $74$ a $32$ entradas; el residuo son parientes con nombre casi
igual que el propio juez acepta (\guillemotleft{}Edsel Ford\guillemotright{}/\guillemotleft{}Henry Ford\guillemotright{},
\guillemotleft{}Otto\guillemotright{}/\guillemotleft{}Karl von Habsburg\guillemotright{}), materia de la señalización
estructural E4.

\emph{Enlazador de nombre a nombre (L2, con L4 parcial).} El escalón 2 pasa a
comparar el tramo con los \emph{alias}: manda el coeficiente de Dice de
trigramas ($\ge 0{,}80$), que mide truncamientos y penaliza la palabra que
falta (\guillemotleft{}bad education movie\guillemotright{} no es \guillemotleft{}Bad Education\guillemotright{}); el coseno de
nombre ($\ge 0{,}85$) solo admite candidatos cuando el léxico no admite a
nadie; entre entradas a menos de $0{,}02$ de la mejor decide la ficha con
descripción. Dos correcciones acompañan: el escalón exacto prueba primero la
forma con el desambiguador conservado (\guillemotleft{}X (1922 film)\guillemotright{} ya no colisiona
con \guillemotleft{}X (1923 film)\guillemotright{}) y las colisiones de alias se resuelven por ficha y
no por orden de carga; y el posesivo se retira solo al final del tramo
(\guillemotleft{}Lothair II's mother\guillemotright{}), porque dentro del tramo es parte del nombre
(\guillemotleft{}God's Gift to Women\guillemotright{}). Sobre $18$ títulos del sondeo truncados como
fallan en el banco, L2 enlaza $16$ frente a $4$ del escalón vectorial de la
v3.

\emph{Plan A generalizado.} Cuatro cambios retiran del flujo lo específico de
2Wiki: la selección final de conjunto se ejecuta en todos los tipos (el
enrutador léxico queda solo como etiqueta del desglose); el analista devuelve,
por hueco sin cubrir, la entidad \emph{ancla} a la que se refiere, y los
candidatos estructurales del salto dirigido salen del pasaje propio de esa
ancla (antes, de una subcadena del nombre en el texto del hueco); las formas
canónicas del analista se enlazan con L2; y el salto dirigido admite varias
rondas: si tras un salto quedan huecos con comodín $[X]$, el analista
re-planifica con los hechos descubiertos a la vista y se salta de nuevo (a lo
sumo tres rondas; en cadenas de dos saltos degenera en una). Las
demostraciones del analista se re-derivaron del conjunto de calibración
(\guillemotleft{}Convoy Buddies / Tip Toes\guillemotright{}, \guillemotleft{}Zoolander 2\guillemotright{}), lo que retira la
primera de las dos reservas de la tabla~\ref{tab:plan-a}.

\begin{table}[t]
\centering\footnotesize\setlength{\tabcolsep}{4pt}
\caption{Escalera del paquete v4 en el sondeo retenido ($200$ preguntas,
analista \texttt{gpt-4o-mini} = paridad). \guillemotleft{}dicc v3\guillemotright{}: mismo plan y enlazador
con el diccionario de la v3. Pareado: diferencia de $\mathrm{FC@}5$ frente al
Plan A del 25 de agosto con el mismo modelo (IC del $95\%$; gana / pierde).}
\label{tab:v4-sondeo}
\begin{tabular}{lcccl}
\toprule
configuración & $\mathrm{R@}5$ & $\mathrm{FC@}5$ & $\mathrm{FC@}10$ & pareado\\
\midrule
Plan A, 25 de agosto & $98{,}5$ & $96{,}0$ & $96{,}5$ & ---\\
v4a: diccionario curado (1.ª) $+$ L2 $+$ plan generalizado & $98{,}6$ & $96{,}5$ & $98{,}0$ & $+0{,}5$ $[-2{,}5,+3{,}5]$ (5/4)\\
\quad v4a con el enlazador de la v3 & $98{,}5$ & $96{,}5$ & $97{,}5$ & L2 neutro en 2Wiki\\
\quad plan generalizado $+$ L2 con dicc v3 & $99{,}2$ & $98{,}0$ & $98{,}5$ & $+2{,}0$ n.s.\\
\textbf{v4} (desambiguador en \guillemotleft{}casi idéntico\guillemotright{}; exacto con paréntesis) & $\mathbf{98{,}8}$ & $\mathbf{97{,}0}$ & $98{,}0$ & $+1{,}0$ $[-1{,}5,+4{,}0]$ (5/3)\\
\bottomrule
\end{tabular}
\end{table}

\emph{Resultado en el sondeo.} La tabla~\ref{tab:v4-sondeo} muestra que
generalizar no cuesta nada en 2Wiki: la v4 queda en paridad estadística con
el Plan A, ligeramente por encima. La ablación con el diccionario de la v3
($98{,}0$) obliga a una lectura honesta: el diccionario curado \emph{pierde}
dos preguntas que la v3 acertaba por fusiones erróneas afortunadas. En
\guillemotleft{}\emph{Where was the director of film Evil Eyes born?}\guillemotright{} el oro es el
pasaje \guillemotleft{}Mark Atkins (footballer)\guillemotright{}: el conjunto de datos enlazó a un
futbolista como director, y la v3 lo fundía con el director homónimo; en
\guillemotleft{}\emph{Who is the husband of Lady Gouyi?}\guillemotright{} la v3 fundía al emperador Wu (el
marido) con su hijo el emperador Zhao (el oro). La tercera pérdida de la
primera v4 ---las cinco películas homónimas fundidas--- era un fallo real y
está corregida. Se adopta la v4: es el principio correcto, es la que
necesitan los conjuntos sin identificadores externos, y las dos preguntas
que cede se documentan como ruido del oro.

\emph{Generalización del banco de pruebas.} El harness admite ya varios
conjuntos con un esquema común (\texttt{\_id}, tipo, pasajes oro) y un
prefijo por conjunto en splits y artefactos; MuSiQue se normaliza desde sus
\texttt{paragraphs} y el número de saltos de su descomposición. Para HotpotQA
se generó el sondeo de $200$ ($100$ de puente, $100$ de comparación, nivel
\emph{hard}, desde el conjunto de entrenamiento; mini-corpus de $1\,990$
pasajes), con BM25 como suelo: banco $\mathrm{R@}5$ $72{,}8$ / $\mathrm{FC@}5$
$49{,}5$ (puente $45{,}4$, comparación $67{,}2$); sondeo $77{,}2$ /
$58{,}5$. Sobre ese sondeo, el Plan A v4 con la configuración congelada y
\texttt{gpt-4o-mini} ---cero ajuste al conjunto--- obtiene $\mathrm{R@}5$
$98{,}8$ / $\mathrm{FC@}5$ $97{,}5$ (puente $95{,}0$, comparación $100$), lo
que sostiene H1 en el sondeo; el banco, con corpus cinco veces mayor, es la
prueba que cuenta. Las ablaciones en ese mismo sondeo reparten la
contribución como en 2Wiki: el salto dirigido aporta $+2{,}0$ de
$\mathrm{FC@}5$ $[+0{,}5,+4{,}0]$ (puente $91{,}0 \to 95{,}0$) y la selección
de conjunto $+4{,}0$ de $\mathrm{R@}2$ $[+1{,}5,+6{,}8]$ ---orden temprano--- y
$+2$ puntos de $\mathrm{FC@}5$ en comparación, lo que justifica ejecutarla en
todos los tipos. 

\emph{Cifra de archivo en paridad estricta.} La medición única del banco de
2Wiki con la configuración v4 congelada y \texttt{gpt-4o-mini} ---el mismo
modelo que la v3, que \hippo{} reproducido y que CatRAG publicado--- da
$\mathrm{R@}2$ $80{,}2$, $\mathrm{R@}5$ $98{,}1$, $\mathrm{FC@}5$ $96{,}0$ y
$\mathrm{FC@}10$ $96{,}7$ (por tipo, $\mathrm{FC@}5$: comparación $99{,}6$,
composicional $95{,}9$, inferencia $84{,}3$, comparación-puente $97{,}9$; fila
añadida a la tabla~\ref{tab:plan-a}). Pareada, supera a la v3 en $+5{,}4$ de
$\mathrm{FC@}5$ $[+3{,}7,+7{,}1]$ (gana $67$, pierde $13$; comparación-puente
$+14{,}0$, composicional $+4{,}8$, ambas significativas) y a \hippo{} en
$+30{,}2$ $[+27{,}2,+33{,}2]$ (gana $317$, pierde $15$); frente a CatRAG
publicado con los mismos modelos ($87{,}0$ / $67{,}6$), $+11{,}1$ de
$\mathrm{R@}5$ y $+28{,}4$ de cadena completa. Esta fila retira la reserva de
paridad de la tabla~\ref{tab:plan-a}: la ganancia del diseño se sostiene con
el modelo de la comparación. Frente al Plan A con \texttt{deepseek-chat} del
25 de agosto queda $-1{,}1$ $[-2{,}2,-0{,}1]$, casi todo en inferencia
($-6{,}5$, n.s.): el modelo adoptado sigue un punto por encima. Su medición
con la v4 congelada, hecha la misma noche, da $\mathrm{R@}2$ $84{,}2$,
$\mathrm{R@}5$ $98{,}9$, $\mathrm{FC@}5$ $97{,}7$ y $\mathrm{FC@}10$
$97{,}8$ (comparación $100$, composicional $97{,}1$, inferencia $92{,}6$,
comparación-puente $98{,}7$; por estructura $100$ / $96{,}2$ / $98{,}7$),
la mejor cifra del proyecto: $+1{,}7$ $[+0{,}8,+2{,}7]$ sobre la v4 con
\texttt{gpt-4o-mini}, $+0{,}6$ $[-0{,}1,+1{,}3]$ (n.s.) sobre el Plan A del
25 de agosto con el mismo modelo, $+7{,}1$ sobre la v3 y $+31{,}9$
$[+29{,}0,+34{,}9]$ sobre \hippo{} (gana $326$, pierde $7$).

\emph{HotpotQA: la primera generalización.} Con el índice construido con el
mismo procedimiento (OpenIE de $9\,811$ pasajes, $87\,629$ aserciones;
diccionario de $48\,856$ entradas; grafo de $146\,274$ nodos y $311\,979$
aristas) y la configuración v4 congelada ---ni un dial, ni una regla, ni una
demostración cambiados---, la medición única del banco de HotpotQA con
\texttt{gpt-4o-mini} da $\mathrm{R@}2$ $77{,}1$, $\mathrm{R@}5$ $94{,}7$,
$\mathrm{FC@}5$ $90{,}2$ y $\mathrm{FC@}10$ $93{,}6$ (puente $93{,}6$ /
$88{,}2$; comparación $99{,}5$ / $98{,}9$). Frente a las referencias
publicadas con los mismos modelos (tabla~\ref{tab:otros-datasets}, fila
añadida), supera a CatRAG en $+5{,}2$ de $\mathrm{R@}5$ y $+9{,}8$ de cadena
completa, y a \hippo{} en $+7{,}6$ / $+14{,}7$; pareada frente a BM25 sobre el
mismo corpus, $+40{,}7$ de $\mathrm{FC@}5$ $[+37{,}1,+44{,}1]$ (gana $432$,
pierde $25$). H1 queda confirmada en el banco, no solo en el sondeo, y con un
margen mayor que el previsto: la hipótesis pedía superar el $80{,}4$ de CatRAG
con la v3, y la v4 lo supera en diez puntos. La distancia con el sondeo
($97{,}5 \to 90{,}2$) es la del corpus cinco veces mayor: de los $98$
fallos, $96$ son de puente y en $90$ falta un solo pasaje oro, el puente,
que queda en los puestos 6--10 en $37$ casos, en 11--20 en $23$ y fuera
del top-20 en $46$; las $35$ preguntas sin mayúsculas interiores fallan
algo más ($6$, un $17\%$), señal de que el reconocimiento de entidades del
analista debería preceder también a la frontera (L1 completo). La reproducción de
\hippo{} sobre el corpus de HotpotQA para la comparación pareada queda
pendiente.

\subsection{Adaptación al asistente GeSIDA}
\label{sec:adaptacion-gesida}

La tabla~\ref{tab:gesida-wiki} fijó qué es patrón y qué es dominio. Lo
que sigue es el plan de porte de los mecanismos validados en 2Wiki al
asistente clínico, en el orden en que conviene ejecutarlo, con el efecto
esperado y la forma de medirlo. El objetivo cuantitativo es el residuo
duro medido en producción (sección~\ref{sec:gesida}): la precisión
temprana ---todas@5 y sección@5 en los niveles 3--4--- sin perder la
cobertura profunda que el modo grafo ya tiene ($84{,}9\%$ en nivel 4 a
$k=20$). Tres restricciones lo enmarcan. Los datos clínicos no salen a
proveedores sin garantías contractuales: queda excluido
\texttt{deepseek-chat}, el modelo del Plan A en 2Wiki, y se usan los
proveedores cubiertos por las garantías documentadas o un modelo local
(la matriz del sondeo muestra que un modelo abierto de $20\,000$ millones
de parámetros basta). La polaridad deóntica y la vigencia siguen siendo
guardarraíles deterministas por encima de cualquier recuperador. Y
ninguna decisión se ajusta sobre el banco completo de $505$ preguntas.

\paragraph{Paso 0: que la medición mida lo desplegado.}
La auditoría del repositorio dejó defectos que invalidan cualquier
comparación previa: el reordenador neuronal no está en la imagen de
despliegue (degradación silenciosa al orden de fusión); la evaluación de
respuestas mide un recuperador denso directo, no el flujo de producción;
la lente numérica comprueba el respaldo por subcadena (\guillemotleft{}50\guillemotright{} casa con
\guillemotleft{}150\guillemotright{}); y la corrección anti-concentrador solo cuenta las aristas de
mención, de modo que las entidades que solo aparecen en recomendaciones
reciben semilla máxima. Se corrigen primero y se fija la línea de base
con el flujo real sobre la muestra estratificada de $100$ y las $73$
preguntas de nivel 4.

\paragraph{Paso 1: filtro selector con compuerta dura.}
Sustituir la compuerta blanda anti-concentrador ---la que en 2Wiki
sembró homónimos--- por la memoria de reconocimiento con compuerta dura:
un selector con demostraciones clínicas ---conservar los saltos
intermedios (fármaco $\to$ clase $\to$ recomendación $\to$ población) y
todos los candidatos de un papel--- sobre $\mathcal{K}(q)$ $=$ Top-40
$\cup$ frontera de las entidades ancladas. Efecto esperado: precisión
temprana, porque solo siembran las entidades de hechos juzgados útiles.
Coste: una llamada por consulta. Se valida en la muestra de $100$ y en
las $73$ de nivel 4 antes de tocar el flujo.

\paragraph{Paso 2: la unidad propia clínica.}
El pasaje propio no tiene equivalente directo: un fármaco no tiene
\guillemotleft{}biografía\guillemotright{}. Su análogo son los nodos \emph{Recomendación} y los
fragmentos de tabla posológica de la entidad: cada entidad sembrada
deposita masa directa en sus recomendaciones vigentes (por la arista
\textsc{mencionada\_en\_rec}, ya existente) y, para las clases, en las de
sus miembros por IS-A. Se mide como en 2Wiki: qué fracción de las
secciones oro entra por siembra directa y cuál por el paseo.

\paragraph{Paso 3: memoria canónica como respaldo del anclaje terminológico.}
El $70\%$ de las entidades sin SCTID sigue resolviéndose con aristas de
sinonimia a umbral, el esquema cuyos fallos motivaron el diccionario. Se
construye la memoria canónica sobre el corpus clínico ---fichas por
mención con las frases de sus aserciones, dos canales, restricciones
duras propias del dominio (dosis o vía distintas no son la misma
entidad), juez estricto y cierre transitivo---, con la estrella por SCTID
como fusión \emph{verificada} que la precede; y el enlazador de consulta
se unifica en la cascada exacto $\to$ trigramas Dice $\to$ vector de
nombre $\to$ juez (evolución L2), que el enlazador terminológico ya
implementa en parte. Prerrequisitos que la auditoría dejó medidos: los
alias de la extracción llegan vacíos y la normalización de mayúsculas
rompe los acrónimos (\guillemotleft{}Vih\guillemotright{}, \guillemotleft{}Dtg\guillemotright{}); y los umbrales de coseno se
recalibran en español con una muestra etiquetada, porque los de 2Wiki
valen para inglés y \texttt{text-embedding-3-small}.

\paragraph{Paso 4: analista, salto dirigido y conjunto con diversidad.}
El Plan A se porta con la semántica del dominio: los huecos del plan son
clínicos (fármaco, clase, población, polaridad, vigencia); los
candidatos estructurales del salto dirigido son las aserciones de la
sección y de las recomendaciones de la entidad; y la selección final de
conjunto incorpora la diversidad por guía que hoy hace el reordenador
(primero el mejor fragmento de cada guía implicada), de modo que un solo
componente cubra ambas funciones. Los huecos sin cubrir son además una
señal natural para la abstención y para la recuperación dirigida del
verificador.

\paragraph{Paso 5: lo que ya está y se conserva.}
La modulación global de transiciones (origen del mecanismo), las
semillas conceptuales IS-A y los sesgos por intención se mantienen; se
añade la caché SQLite de llamadas adoptada en 2Wiki, que hace
reproducibles las mediciones y estabiliza la varianza del verificador.
Del análisis del código de \hippo{} se hereda un aviso para la fase
incremental: su reindexado duplica las aristas de sinonimia; la vía es la
consolidación nocturna por lotes, no el incremental ingenuo.

\paragraph{Protocolo y expectativa.}
Cada paso se mide de forma pareada sobre la muestra de $100$ y las $73$
de nivel 4 (cobertura de guías y de secciones por nivel, alguna@1, MRR de
guía), con intervalos bootstrap como en 2Wiki; el banco completo de
$505$ se mide una sola vez al final. La expectativa, por analogía con la
escalera de 2Wiki, es que el filtro con compuerta dura y la memoria
canónica muevan la precisión temprana ---el análogo del salto de la v2 a
la v3--- y que el Plan A mueva las preguntas multi-guía de nivel 4 ---el
análogo del doble puente---. La validación clínica con patrón oro sigue
siendo el paso previo a cualquier uso asistencial.

\subsection{Generalización y comparación con otros conjuntos de datos}
\label{sec:generalizacion}

\paragraph{Conjuntos y referencias.}
Los subconjuntos de reproducción publicados con \hippo{} para MuSiQue
($1\,000$ preguntas, $11\,656$ pasajes; saltos etiquetados: $518$ de dos,
$316$ de tres y $166$ de cuatro) y HotpotQA ($1\,000$ preguntas, $9\,811$
pasajes; $811$ de puente y $189$ de comparación) están vendorizados junto
al de 2Wiki y son exactamente los que usa CatRAG, de modo que sus números
publicados con \texttt{gpt-4o-mini} + \texttt{text-embedding-3-small}
son la referencia directa (tabla~\ref{tab:otros-datasets}). Las
diferencias con 2Wiki que importan al método son tres. (i)~Profundidad:
MuSiQue exige cadenas de tres y cuatro pasajes con un puente por
eslabón, no dos cadenas de dos como comparación-puente, y ofrece la
descomposición oro de cada pregunta (\texttt{question\_decomposition},
con \guillemotleft{}\#1\guillemotright{} como comodín del puente), que permite medir por primera
vez la \emph{calidad del plan} del analista y no solo la recuperación.
(ii)~Corpus: HotpotQA y 2Wiki tienen un pasaje por título; MuSiQue tiene
$647$ títulos repetidos ($11\,656$ pasajes, $9\,838$ títulos), de modo que
el pasaje propio deja de ser único y debe definirse como el
\emph{conjunto} de pasajes con el título de la entidad (masa repartida o
completa: decisión a fijar en el sondeo). (iii)~Tipos: no existe la
taxonomía de cuatro tipos; el enrutador léxico de 2Wiki no aplica.

\begin{table}[t]
\centering\footnotesize\setlength{\tabcolsep}{4pt}
\caption{Referencias publicadas por~\cite{catrag} en MuSiQue y HotpotQA
con \texttt{gpt-4o-mini} + \texttt{text-embedding-3-small} (los mismos
modelos de este trabajo) sobre los mismos subconjuntos de $1\,000$
preguntas; FCR $=\mathrm{FC@}5$; F1 con lector Llama-3.3-70B. Última
fila: \hippo{} con NV-Embed-v2 (7B), de~\cite{hipporag2}, como techo de
hardware.}
\label{tab:otros-datasets}
\begin{tabular}{lcccccc}
\toprule
 & \multicolumn{3}{c}{MuSiQue} & \multicolumn{3}{c}{HotpotQA}\\
\cmidrule(lr){2-4}\cmidrule(lr){5-7}
sistema & $\mathrm{R@}5$ & FCR & F1 & $\mathrm{R@}5$ & FCR & F1\\
\midrule
BM25 & $31{,}6$ & $6{,}4$ & $22{,}9$ & $64{,}8$ & $38{,}3$ & $54{,}1$\\
Contriever & $46{,}6$ & $11{,}3$ & $31{,}3$ & $75{,}3$ & $54{,}1$ & $62{,}3$\\
GTR & $49{,}1$ & $15{,}0$ & $34{,}6$ & $73{,}9$ & $53{,}0$ & $62{,}8$\\
denso \texttt{te3-small} & $55{,}4$ & $21{,}1$ & $36{,}1$ & $81{,}3$ & $64{,}9$ & $64{,}6$\\
RAPTOR & $53{,}3$ & $19{,}6$ & $36{,}0$ & $79{,}5$ & $61{,}4$ & $64{,}4$\\
\hippo{} & $61{,}4$ & $30{,}5$ & $43{,}2$ & $87{,}1$ & $75{,}5$ & $69{,}4$\\
CatRAG & $64{,}9$ & $34{,}6$ & $45{,}0$ & $89{,}5$ & $80{,}4$ & $71{,}4$\\
\textbf{Plan A v4 (este trabajo; \S\ref{sec:v4})} & --- & --- & --- & $\mathbf{94{,}7}$ & $\mathbf{90{,}2}$ & ---\\
\midrule
\hippo{} con NV-Embed-v2 (7B)~\cite{hipporag2} & $74{,}7$ & --- & --- & $96{,}3$ & --- & ---\\
\bottomrule
\end{tabular}
\end{table}

\paragraph{Qué debe generalizar y cómo.}
Cuatro piezas de la versión evaluada son específicas de 2Wiki y se
sustituyen por sus equivalentes generales antes de medir, con la
configuración \emph{congelada} para los tres conjuntos (cero ajuste por
conjunto de datos).
\begin{enumerate}[leftmargin=1.6em, itemsep=2pt]
\item \textbf{Enrutador.} Una única estrategia para todo: analista con
  plan (Plan A) o, en la versión de paridad, filtro selector. La variante
  de descomposición, que en 2Wiki solo servía a comparación, queda como
  ablación; en HotpotQA la comparación se detecta con la única marca
  general (\emph{are both / same}) o no se detecta, y la diferencia se
  mide.
\item \textbf{Enlazador.} Reconocimiento de entidades por el analista
  (L1) y escalón 2 de nombre a nombre (L2) en lugar del recorte por
  mayúsculas, que en MuSiQue ---preguntas largas, anidadas, con
  menciones descriptivas (\guillemotleft{}the person who \dots\guillemotright{})--- fallaría más
  que en 2Wiki.
\item \textbf{Memoria canónica.} El patrón (ficha, dos canales,
  restricciones, juez, cierre) es general; la heurística
  título$\to$sujeto y el pasaje propio se conservan porque los tres
  corpus son enciclopédicos, con el pasaje propio definido como conjunto
  en MuSiQue. La auditoría con identificadores externos (E1) no es
  posible en MuSiQue ni en HotpotQA, que no traen QIDs: la precisión de
  fusión se estima con una muestra etiquetada por corpus.
\item \textbf{Profundidad.} El grafo por capas no encadena saltos en el
  paseo (sección~\ref{sec:ppr}): cada puente lo descubre el filtro y lo
  siembra el pasaje propio. Para cadenas de tres y cuatro, el salto
  dirigido pasa de una ronda a \emph{una por hueco pendiente} (a lo sumo
  saltos$-1$ rondas) y el tope dinámico del analista crece con el número
  de huecos. Es la única modificación del método motivada por un
  conjunto nuevo, y se decide en el sondeo.
\end{enumerate}

\paragraph{Protocolo.}
Idéntico al de la sección~\ref{sec:protocolo}: para cada conjunto, un
sondeo de $200$ preguntas extraído del \emph{split} de entrenamiento
etiquetado (disjunto de las $1\,000$) con su mini-corpus, para las
decisiones anteriores; diales pre-registrados (los de 2Wiki); una única
medición sobre las $1\,000$; \hippo{} reproducido con su código sobre los
mismos corpus (dos llamadas por pasaje) para la comparación pareada; y
las cifras de CatRAG como referencia externa. Métricas: $\mathrm{R@}5$ y
$\mathrm{FC@}5$ ---y $\mathrm{FC@}10$, porque en MuSiQue una cadena de
cuatro deja una sola posición de holgura a $k=5$, como
comparación-puente---, desglosadas por número de saltos, y EM/F1/JSR
cuando se añada el lector. Coste estimado con \texttt{gpt-4o-mini} en
indexación ($21\,467$ pasajes, una llamada cada uno) y consulta
($2\,000$ preguntas, $\sim 2$ llamadas): $15$--$20$~USD, la mitad de ellos
la reproducción de \hippo{}.

\paragraph{Hipótesis falsables.}
\begin{enumerate}[leftmargin=1.6em, itemsep=1pt, label=H\arabic*.]
\item En HotpotQA (dos saltos, $81\%$ de puente) la v3 supera la cadena
  completa publicada de CatRAG ($80{,}4$) con los mismos modelos, porque
  el segmento equivale al composicional de 2Wiki ($91{,}0$ frente a
  $77{,}2$ de \hippo{}).
\item En MuSiQue la ventaja se concentra en tres y cuatro saltos y
  depende del salto dirigido: sin él, la cadena completa cae con la
  profundidad según el modelo multiplicativo $p^{h}$; con rondas por
  hueco se mantiene por encima de esa predicción, como ocurrió en el
  doble puente de 2Wiki.
\item La calidad del plan del analista, medida contra la descomposición
  oro de MuSiQue (huecos correctos y en orden), predice la cadena
  completa mejor que la afinidad máxima de la pregunta: el error residual
  está en el plan, no en el grafo.
\end{enumerate}
El caso clínico (sección~\ref{sec:adaptacion-gesida}) es la
generalización en la otra dirección ---corpus no enciclopédico, otro
idioma, terminología externa--- y comparte con esta el mismo protocolo.

% =========================================================================
\section{Consideraciones de validez y limitaciones}
\label{sec:validez}

\begin{enumerate}[leftmargin=1.6em, itemsep=2pt]
\item \textbf{Contaminación de diseño.} Tres preguntas del banco de pruebas
  se inspeccionaron durante el desarrollo (trazas de la iteración~2), y las
  demostraciones del filtro se redactaron a partir de sus \emph{modos} de
  fallo (patrones generales, sin respuestas); el análisis de errores que
  motivó el diccionario se hizo sobre agregados del banco a nivel de
  mecanismo. Todas las mejoras se validaron primero en el sondeo
  (retenido), donde reproducen el mismo efecto
  ($\mathrm{FC@}5$ $64{,}0 \to 83{,}5 \to 92{,}0$), lo que descarta la
  memorización como explicación. Para la versión de archivo, las
  demostraciones deben re-derivarse del conjunto de calibración y las tres
  preguntas inspeccionadas, excluirse o señalarse.
\item \textbf{Cuello de botella actual.} El estudio de errores de la
  versión final (sobre fallos del sondeo, retenido) atribuye la mayoría de
  los fallos restantes al juicio del filtro: en $12$ de $15$ casos el hecho
  puente \emph{estaba} en la lista de candidatos y el modelo pequeño no lo
  seleccionó (con el tope de $4$ hechos, insuficiente en cadenas de
  comparación-puente que requieren al menos $4$). Las heurísticas léxicas
  (enrutador, tramos capitalizados del enlazador) no son causa medible de
  fallo en este conjunto, pero sí un límite de generalización a otros
  idiomas y estilos de pregunta. El respaldo vectorial del enlazador está
  además inactivo en la práctica (acepta $3$ de $275$ tramos): en esos
  casos el sistema evaluado se sostiene sobre el filtro y la presa densa,
  y su corrección (sección~\ref{sec:linker}, evoluciones L1--L2) queda
  pendiente.
\item \textbf{Alcance.} Los resultados corresponden a un único conjunto de
  datos (2WikiMultihopQA) y a la métrica de recuperación; la extensión a
  MuSiQue y HotpotQA y la evaluación extremo a extremo de respuesta
  (EM/F1 con lector compartido) están planificadas con el mismo protocolo
  (sección~\ref{sec:generalizacion}).
\item \textbf{Memoria canónica de entidades.} La heurística título$\to$sujeto
  del diccionario está sobreaplicada ($1\,549$ uniones fuera de su caso
  previsto, con parientes y secuelas fundidos; sección~\ref{sec:diccionario})
  y la precisión de las fusiones no está medida; la auditoría con los
  identificadores Wikidata del propio conjunto de datos y la corrección de la
  heurística, con re-medición, quedan pendientes para la versión de archivo.
\item \textbf{Reproducibilidad.} Todos los artefactos intermedios
  (extracción, diccionario, grafos, incrustaciones, resultados por
  pregunta) y las cachés de las llamadas al modelo se conservan; cada
  medición es re-ejecutable de forma determinista y sin coste de API.
\end{enumerate}

% =========================================================================
\begin{thebibliography}{99}\small

\bibitem{hipporag} B.~Jiménez Gutiérrez, Y.~Shu, Y.~Gu, M.~Yasunaga,
Y.~Su. HippoRAG: Neurobiologically inspired long-term memory for large
language models. \emph{NeurIPS}, 2024. arXiv:2405.14831.

\bibitem{hipporag2} B.~Jiménez Gutiérrez, Y.~Shu, W.~Qi, S.~Zhou, Y.~Su.
From RAG to memory: Non-parametric continual learning for large language
models. 2025. arXiv:2502.14802.

\bibitem{catrag} K.\,H.~Lau, F.~Zhang, B.~Ruan, Y.~Zhou, Q.~Guo,
R.~Zhang, X.~Zhou. Breaking the static graph: Context-aware traversal for
graph-based RAG. \emph{Findings of ACL}, págs.~5849--5863, 2026.

\bibitem{proprag} J.~Wang, J.~Han. PropRAG: Guiding retrieval with beam
search over proposition paths. 2025.

\bibitem{raptor} P.~Sarthi, S.~Abdullah, A.~Tuli, S.~Khanna, A.~Goldie,
C.\,D.~Manning. RAPTOR: Recursive abstractive processing for
tree-organized retrieval. \emph{ICLR}, 2024.

\bibitem{graphrag} D.~Edge, H.~Trinh, N.~Cheng, J.~Bradley, A.~Chao,
A.~Mody, S.~Truitt, J.~Larson. From local to global: A Graph RAG approach
to query-focused summarization. 2024. arXiv:2404.16130.

\bibitem{lightrag} Z.~Guo, L.~Xia, Y.~Yu, T.~Ao, C.~Huang. LightRAG:
Simple and fast retrieval-augmented generation. 2024. arXiv:2410.05779.

\bibitem{hypergraphrag} H.~Luo et~al. HyperGraphRAG: Retrieval-augmented
generation via hypergraph-structured knowledge representation. 2025.
arXiv:2503.21322.

\bibitem{hoprag} H.~Liu, Z.~Wang, X.~Chen, Z.~Li, F.~Xiong, Q.~Yu,
W.~Zhang. HopRAG: Multi-hop reasoning for logic-aware retrieval-augmented
generation. 2025. arXiv:2502.12442.

\bibitem{contriever} G.~Izacard, M.~Caron, L.~Hosseini, S.~Riedel,
P.~Bojanowski, A.~Joulin, E.~Grave. Unsupervised dense information
retrieval with contrastive learning. \emph{TMLR}, 2022.

\bibitem{gtr} J.~Ni, C.~Qu, J.~Lu, Z.~Dai, G.~Hernández Ábrego, J.~Ma,
V.~Zhao, Y.~Luan, K.~Hall, M.-W.~Chang, Y.~Yang. Large dual encoders are
generalizable retrievers. \emph{EMNLP}, 2022.

\bibitem{nvembed} C.~Lee, R.~Roy, M.~Xu, J.~Raiman, M.~Shoeybi,
B.~Catanzaro, W.~Ping. NV-Embed: Improved techniques for training LLMs as
generalist embedding models. \emph{ICLR}, 2025.

\bibitem{pprsurvey} M.~Yang, H.~Wang, Z.~Wei, S.~Wang, J.-R.~Wen.
Efficient algorithms for personalized PageRank computation: A survey.
\emph{IEEE TKDE}, 2024. arXiv:2403.05198.

\bibitem{zep} P.~Rasmussen, P.~Paliychuk, T.~Beauvais, J.~Ryan,
D.~Chalef. Zep: A temporal knowledge graph architecture for agent memory.
2025. arXiv:2501.13956.

\bibitem{magma} D.~Jiang, Y.~Li, G.~Li, B.~Li. MAGMA: A multi-graph based
agentic memory architecture for AI agents. 2026. arXiv:2601.03236.

\bibitem{agentmemory} C.~Yang, C.~Zhou, Y.~Xiao et~al. Graph-based agent
memory: Taxonomy, techniques, and applications. 2026. arXiv:2602.05665.

\bibitem{wiki2} X.~Ho, A.-K.~Duong Nguyen, S.~Sugawara, A.~Aizawa.
Constructing a multi-hop QA dataset for comprehensive evaluation of
reasoning steps. \emph{COLING}, 2020.

\bibitem{umls} O.~Bodenreider. The Unified Medical Language System (UMLS):
integrating biomedical terminology. \emph{Nucleic Acids Research},
32(suppl.~1):D267--D270, 2004.

\end{thebibliography}

\end{document}
